\documentclass[reqno]{amsart}
% Self-contained preamble, derived on 2026-11-25 from common.tex (the author's shared macro file), keeping only what this paper uses.
\usepackage{amssymb}
\usepackage{enumerate}
\usepackage{hyperref}

\DeclareMathOperator{\vol}{vol}
\newcommand{\qr}[2]{\left( \frac{#1}{#2} \right)}

\theoremstyle{plain}
\newtheorem{theorem}{Theorem}
\newtheorem{corollary}[theorem]{Corollary}
\newtheorem{proposition}[theorem]{Proposition}
\theoremstyle{definition}
\newtheorem{definition}[theorem]{Definition}
\newtheorem{example}[theorem]{Example}
\newtheorem{remark}[theorem]{Remark}
\newtheoremstyle{itplain}{6pt}{5pt\topsep}{\itshape}{}{\itshape}{.}{5pt plus 1pt minus 1pt}{}
\theoremstyle{itplain}
\newtheorem{lemma}[theorem]{Lemma}
\setlength{\emergencystretch}{3em}

% Links to the Lean formalization in lean/ (Appendix~\ref{sec:lean}).
% \lean{Decl} follows the syntax of leanblueprint.  In the PDF it links to the
% source line of Decl at commit \leancommit; in the HTML version it links to the
% row of Decl in Appendix C, which carries the same link.  The rows
% \leanentry{Decl}{file}{line} are kept current by scripts/lean-links.py.
\newcommand{\leanrepo}{https://github.com/ultronozm/metaplectic-summation}
\newcommand{\leancommit}{3b301a97286276f336773d5d16b00d96d020b185}
\newcommand{\leanurl}[2]{\leanrepo/blob/\leancommit/lean/#1\#L#2}
\newcommand{\leanmark}{\textsuperscript{\textsf{Lean}}}
\newcommand{\lean}[1]{\hyperlink{lean:#1}{\leanmark}}
\newcommand{\leanentry}[3]{\hypertarget{lean:#1}{\href{\leanurl{#2}{#3}}{\texttt{#1}}}}
\makeatletter
\newcommand{\leandecl}[3]{\global\@namedef{leanloc@#1}{{#2}{#3}}}
\AtBeginDocument{%
  \newcommand{\leanlink}[2]{\href{\leanurl{#1}{#2}}{\leanmark}}%
  \renewcommand{\lean}[1]{%
    \@ifundefined{leanloc@#1}{\leanmark\PackageWarning{lean}{Unknown declaration #1}}%
      {\expandafter\expandafter\expandafter\leanlink\csname leanloc@#1\endcsname}}%
  \renewcommand{\leanentry}[3]{%
    \immediate\write\@auxout{\string\leandecl{#1}{#2}{#3}}%
    \href{\leanurl{#2}{#3}}{\nolinkurl{#1}}}}
\makeatother

\numberwithin{equation}{section}
\numberwithin{theorem}{section}

\begin{document}

\author{Paul D. Nelson}

% \address{Aarhus University, Denmark}

\address{Aarhus University, Ny Munkegade 118,
  DK-8000 Aarhus, Denmark}
\email{paul.nelson@math.au.dk}

\title{A Schwartz space for the $n$-fold cover of $\mathrm{SL}_2$}

% \date{\today.  Preliminary draft; comments welcome.}

\subjclass[2020]{Primary 11F70; Secondary 11F27, 11L05, 22E50}
\keywords{Metaplectic groups, Poisson summation, Schwartz spaces, Eisenstein series, Gauss sums, Voronoi summation}

\thanks{The author is supported by a research grant (VIL54509) from VILLUM FONDEN}

\begin{abstract}
  Poisson summation on the plane gives the functional equation for Eisenstein series on $\mathrm{SL}_2$.  We develop an analogue of Poisson summation for Kubota's $n$-fold cover of $\mathrm{SL}_2$, in the spirit of Braverman--Kazhdan.  Locally, we define a Schwartz space consisting of genuine functions on Deligne's $\mu_n$-torsor over the punctured plane.  It comes with a Fourier transform $\mathcal{F}$ defined like the symplectic Fourier transform, but with the symplectic form replaced by its $n$th roots.  At almost all places, we define a basic vector fixed by $\mathcal{F}$.  Globally, we prove a Poisson summation formula for the resulting adelic Schwartz space.  Taking a Fourier coefficient gives a general Voronoi-type summation formula for Gauss sums of order $n$, with an explicit integral transform relating the weights.  For cubic Gauss sums, the formula takes sextic residue twists to quadratic twists, as in recent work on the quasi-Riemann hypothesis.
\end{abstract}

\maketitle

\setcounter{tocdepth}{1} \tableofcontents

\section{Introduction}\label{sec:introduction}

\subsection{Poisson summation on the plane}\label{sec:poisson-summation-plane}
Let $\Phi$ be a Schwartz function on $\mathbb{R}^2$, and let
\begin{equation*}
  \mathcal{F} \Phi(x', y') := \int_{\mathbb{R}^2} \Phi(x, y) e(x y' - y x') \, d x \, d y,
  \qquad
  e(z) := e^{2 \pi i z}
\end{equation*}
denote its symplectic Fourier transform.  The operator $\mathcal{F}$ commutes with the right translations $\Phi \mapsto \Phi(\cdot \, g)$ by $g \in \mathrm{SL}_2(\mathbb{R})$.  The Poisson summation formula reads
\begin{equation}\label{eq:intro-poisson-Z}
  \sum_{v \in \mathbb{Z}^2} \Phi(v) = \sum_{v \in \mathbb{Z}^2} \mathcal{F} \Phi(v).
\end{equation}
Applying the formula to $\Phi(\cdot \, g)$ gives an identity of functions on \(\mathrm{SL}_2(\mathbb Z)\backslash\mathrm{SL}_2(\mathbb R)\).

Every nonzero vector in $\mathbb{Z}^2$ is a positive integer multiple of a unique primitive vector.  The primitive vectors are the bottom rows of the matrices in $\Gamma := \mathrm{SL}_2(\mathbb{Z})$, so, writing $\Gamma_\infty$ for the stabilizer of $(0, 1)$ in $\Gamma$, we have
\begin{equation*}
  E_\Phi(g)
  := \sum_{v \in \mathbb{Z}^2 - \{0\}} \Phi(v g)
  = \sum_{\gamma \in \Gamma_\infty \backslash \Gamma} \,
  \sum_{t \geq 1} \Phi \bigl( t (0, 1) \gamma g \bigr).
\end{equation*}
Thus $E_\Phi$ is a pseudo-Eisenstein series on $\Gamma \backslash \mathrm{SL}_2(\mathbb{R})$.  Poisson summation gives its functional equation
\begin{equation}\label{eq:intro-pseudo-eisenstein}
  E_\Phi(g) + \Phi(0) = E_{\mathcal{F} \Phi}(g) + \mathcal{F} \Phi(0).
\end{equation}
Replacing $\Phi$ by its dilates $v \mapsto \Phi(r v)$ and taking the Mellin transform in $r$ gives the meromorphic continuation and functional equation of the completed Eisenstein series for $\Gamma$ (see~\cite[\S\S3--5]{MR0210827}); the terms $\Phi(0)$ and $\mathcal{F} \Phi(0)$ account for its poles.  Conversely, given the analytic properties of the Eisenstein series, which may be established independently \cite{Se56}, Mellin inversion recovers the Poisson formula, with the terms $\Phi(0)$ and $\mathcal{F} \Phi(0)$ arising as residues.

More generally, let $F$ be a global field with adele ring $\mathbb{A}$, let $\psi$ be a nontrivial unitary character of $\mathbb{A} / F$, and let $\mathcal{F}$ denote the symplectic Fourier transform on $\mathbb{A}^2$ defined using $\psi$.  For a Schwartz--Bruhat function $\Phi = \otimes_{\mathfrak{p}} \Phi_{\mathfrak{p}}$ on $\mathbb{A}^2$, we have
\begin{equation}\label{eq:intro-poisson-A}
  \sum_{v \in F^2} \Phi(v) = \sum_{v \in F^2} \mathcal{F} \Phi(v).
\end{equation}
Let $S$ be a nonempty finite set of places of $F$ containing the archimedean places, let $F_S := \prod_{\mathfrak{p} \in S} F_{\mathfrak{p}}$, and let $\mathcal{O}_S$ denote the ring of $S$-integers.  If $\psi_{\mathfrak{p}}$ is unramified and $\Phi_{\mathfrak{p}}$ is the characteristic function of $\mathfrak{o}_{\mathfrak{p}}^2$ for each $\mathfrak{p} \notin S$, then $\mathcal{F} \Phi_{\mathfrak{p}} = \Phi_{\mathfrak{p}}$ for such $\mathfrak{p}$, and \eqref{eq:intro-poisson-A} is the Poisson summation formula for the lattice $\mathcal{O}_S^2$ in $F_S^2$.  As before, applying Poisson summation to \(\Phi(\cdot\,g)\), now with \(g\in\mathrm{SL}_2(\mathbb A)\), gives the functional equation of a pseudo-Eisenstein series on \(\mathrm{SL}_2(F)\backslash\mathrm{SL}_2(\mathbb A)\) (see~\cite[(7.8)--(7.10)]{MR1694894}, \cite[\S1.1]{MR871663}).

This approach to principal Eisenstein series on split reductive groups $G$ goes back to Gelfand, Graev and Piatetskii-Shapiro \cite[\S\S3.4.10--3.4.13]{MR3468638}.  Braverman and Kazhdan \cite{MR1694894} formulate it in terms of a Schwartz space on the quotient \(U\backslash G\), where $U$ is a maximal unipotent subgroup of $G$.  For \(G=\mathrm{SL}_2\), this quotient is the punctured plane.  Recent work in this direction includes constructions of basic functions and Schwartz spaces on reductive monoids \cite{BouthierNgoSakellaridis2016, HsuYao2026}, constructions of Fourier transforms on spherical varieties \cite{MR3969881, 2021arXiv210310261G, GurevichKazhdan2025}, and proofs of Poisson summation formulas related to Siegel and Schubert Eisenstein series \cite{GetzLiu2021, 2021arXiv2107.01874}.

\subsection{The metaplectic case}\label{sec:metaplectic-case}
Let $n \geq 1$, and suppose that $F$ contains a primitive $n$th root of unity.  We write $\mu_n$ for the group of $n$th roots of unity in $F$.  What is the analogue of Poisson summation on the plane for Kubota's Eisenstein series \cite{MR0255490} on the $n$-fold metaplectic cover $\widetilde{\mathrm{SL}}_2$ of $\mathrm{SL}_2$?

At each completion of $F$ (indeed, over any local field containing $\mu_n$), Deligne \cite[\S\S1, 3.7]{SB_1978-1979__21__244_0} constructs a $\mu_n$-torsor over the punctured plane.  The metaplectic group acts on the torsor, lifting the usual action of $\mathrm{SL}_2$ on the plane.  We define a Schwartz space consisting of genuine functions on this torsor, equipped with a Fourier transform, and prove a Poisson summation formula for the resulting adelic Schwartz space that reduces to \eqref{eq:intro-poisson-A} when $n = 1$.

The local ingredients, defined precisely in \S\S\ref{sec:hilbert-symbol}--\ref{sec:germs}, are as follows:
\begin{itemize}
\item Fixing an embedding $\mu_n \hookrightarrow \mathbb{C}^\times$, a complex-valued function on the torsor is called \emph{genuine} if it is $\mu_n$-equivariant.  In coordinates, a genuine function is a pair $(\Phi_{\mathfrak{p}, X}, \Phi_{\mathfrak{p}, Y})$ of smooth functions, defined away from the lines $x = 0$ and $y = 0$, respectively, and related by
  \begin{equation*}
    \Phi_{\mathfrak{p}, X}(x, y) = \{x, y\}_{\mathfrak{p}, n} \Phi_{\mathfrak{p}, Y}(x, y),
  \end{equation*}
  where $\{x, y\}_{\mathfrak{p}, n}$ denotes the Hilbert symbol with values in $\mu_n$.
\item Deligne \cite[\S6.1]{SB_1978-1979__21__244_0} realizes the principal series representations as spaces of homogeneous genuine functions on the torsor.  His formula for the standard intertwining operator \cite[\S6.3]{SB_1978-1979__21__244_0} involves an integral over the line $x y ' - y x ' = 1$, using his trivialization of the torsor along that line \cite[\S1.13]{SB_1978-1979__21__244_0}.  To define $\mathcal{F}$, we integrate over the line
  \begin{equation*}
    x y' - y x' = \tau^n
  \end{equation*}
  and take the Fourier transform in $\tau$ at $1$.
\item Write \(\mathcal D_{\mathfrak p}\) for the space of genuine test functions at \(\mathfrak p\) (compact support in $F_{\mathfrak{p}}^2 - \{0\}$ if $\mathfrak{p}$ is finite, rapid decay at the origin and at infinity if $\mathfrak{p}$ is archimedean).  Following \cite[\S3.1]{MR1694894}, we define the Schwartz space
  \begin{equation*}
    \mathcal S_{\mathfrak p} := \mathcal{D}_{\mathfrak{p}} + \mathcal{F} \mathcal{D}_{\mathfrak{p}}.
  \end{equation*}
  It is stable under $\widetilde{\mathrm{SL}}_2$ and $\mathcal{F}$ (Theorem~\ref{theorem:schwartz-space}).
\item At a finite place $\mathfrak{p} \nmid n$, the characteristic function of $\mathfrak{o}_{\mathfrak{p}}^2$ is replaced by a \emph{basic vector} $\Phi_{\mathfrak{p}}^0$.  It is supported on the vectors $\varpi^k (x_0, y_0)$ with $(x_0, y_0) \in \mathfrak{o}_{\mathfrak{p}}^2$ primitive and $k \in n \mathbb{Z}_{\geq 0}$, and its $X$-component there is $q^{(n-1) k / n}$ or $q^{(n-1) k / n} \{x_0, y_0\}_{\mathfrak{p}, n}$, according as $x_0$ is a unit or not.  It belongs to $\mathcal{S}_{\mathfrak{p}}$ and is fixed by $\mathcal{F}$ when $\psi_{\mathfrak{p}}$ is unramified (Proposition~\ref{proposition:basic-vector-fixed}).
\end{itemize}
For $n = 1$, these are the objects considered in \S\ref{sec:poisson-summation-plane}.

Let $\Phi = \otimes_{\mathfrak{p}} \Phi_{\mathfrak{p}}$ be a pure tensor in the restricted tensor product $\mathcal{S}_{\mathbb{A}} := \bigotimes'_{\mathfrak{p}} \mathcal{S}_{\mathfrak{p}}$ with respect to the basic vectors $\Phi_{\mathfrak{p}}^0$.  By Hilbert's reciprocity law, the products of the local values of the $\Phi_{\mathfrak{p}}$ define values $\Phi(v)$ at the points $v \in F^2 - \{0\}$.

For the sake of illustration, suppose that $S$ contains the places dividing $n$, that $\Phi_{\mathfrak{p}} = \Phi_{\mathfrak{p}}^0$ and $\psi_{\mathfrak{p}}$ is unramified for each $\mathfrak{p} \notin S$, and that $\mathcal{O}_S$ is a principal ideal domain.  Evaluating the basic vectors outside \(S\), we obtain (see Example~\ref{example:S-arithmetic})
\begin{equation}\label{eq:intro-theta-S}
  \sum_{v \in F^2 - \{0\}} \Phi(v)
  =
  \sum_{\tau} \lvert \tau \rvert_S^{n - 1}
  \sum_{(c, d)} \qr{d}{c}_n \Phi_S(\tau^n c, \tau^n d),
\end{equation}
where $\tau$ runs over representatives for the nonzero elements of $\mathcal{O}_S$ modulo units, $(c, d)$ runs over the pairs of coprime elements of $\mathcal{O}_S$, $\lvert \tau \rvert_S := \prod_{\mathfrak{p} \in S} \lvert \tau \rvert_{\mathfrak{p}}$, and $\qr{d}{c}_n$ denotes the $n$th power residue symbol (Definition~\ref{definition:residue-symbol}).  Here $\Phi_S$ denotes the product of the $X$-components of the $\Phi_{\mathfrak{p}}$ with $\mathfrak{p} \in S$; when $c = 0$, it denotes the product of their $Y$-components, and the residue symbol is omitted.

For \(n\geq 2\), an element \(\Phi_{\mathfrak p}\in\mathcal S_{\mathfrak p}\) has no natural value at the origin.  The role of this value is played by the \emph{leading term} \(b\Phi_{\mathfrak p}\), the genuine function on the punctured plane given by \(u\mapsto\lim_{\mu\to0}|\mu|_{\mathfrak p}^{n-1}\Phi_{\mathfrak p}(\mu^n u)\).  It vanishes if $\Phi_{\mathfrak{p}} \in \mathcal{D}_{\mathfrak{p}}$.  Our main result (Theorem \ref{theorem:poisson}) asserts that
\begin{equation}\label{eq:intro-poisson-boundary}
  \sum_{v \in F^2 - \{0\}} \Phi(v) + B(b \Phi)
  =
  \sum_{v \in F^2 - \{0\}} \mathcal{F} \Phi(v) + B(b \mathcal{F} \Phi),
\end{equation}
where $\mathcal{F} \Phi := \otimes_{\mathfrak{p}} \mathcal{F} \Phi_{\mathfrak{p}}$ and $b \Phi := \otimes_{\mathfrak{p}} b \Phi_{\mathfrak{p}}$.  Here $B$ is a linear functional on the restricted tensor product of the spaces of local leading terms, and the identity \eqref{eq:intro-poisson-boundary} determines $B$ uniquely.  For $n = 1$, we have $B(b \Phi) = \prod_{\mathfrak{p}} \Phi_{\mathfrak{p}}(0)$, and \eqref{eq:intro-poisson-boundary} specializes to \eqref{eq:intro-poisson-A}.  For $n \geq 2$, we do not determine $B$ explicitly.

Both boundary terms vanish if some factor of $\Phi$ and some factor of $\mathcal{F} \Phi$ have leading term zero.  This holds if some factor $\Phi_{\mathfrak{p}}$ lies in $\mathcal{D}_{\mathfrak{p}}$ and satisfies a vanishing condition (Definition \ref{definition:genuine-vanishing}), in which case (Corollary \ref{corollary:poisson-vanishing})
\begin{equation}\label{eq:intro-poisson-metaplectic}
  \sum_{v \in F^2 - \{0\}} \Phi(v) = \sum_{v \in F^2 - \{0\}} \mathcal{F} \Phi(v).
\end{equation}

Kubota's Eisenstein series are the sums, over the coprime pairs $(c, d)$ modulo $n$th powers of units, of $\qr{d}{c}_n$ times a function of $(c, d)$ that is homogeneous with respect to the dilations by $n$th powers.  The case $n = 1$ suggests that applying \eqref{eq:intro-poisson-boundary} to dilates of $\Phi$ and taking a Mellin transform should recover Kubota's functional equations, with the boundary terms corresponding to the residues of the Eisenstein series.  We do not pursue this Mellin interpretation here.

\subsection{Sums of Gauss sums}\label{sec:intro-gauss-sums}
We take the Fourier coefficient of the sum $\sum_v \Phi(v)$ in \eqref{eq:intro-poisson-boundary} with respect to the translations $(x, y) \mapsto (x, y + u x)$, integrating against $\psi(-u)$ over $u \in \mathbb{A} / F$.  The terms with $x = 0$ drop out, and those with $x = c \in F^\times$ unfold to the product over $\mathfrak{p}$ of the local integrals
\begin{equation*}
  W_{\mathfrak{p}}(c) := \int_{F_{\mathfrak{p}}} \Phi_{\mathfrak{p}, X}(c, y) \psi_{\mathfrak{p}}(-y/c) \, d y.
\end{equation*}
Every $W_{\mathfrak{p}} \in C_c^\infty(F_{\mathfrak{p}}^\times)$ arises in this way from some $\Phi_{\mathfrak{p}} \in \mathcal{D}_{\mathfrak{p}}$ (Lemma \ref{lemma:jacquet-surjective}).  Replacing $\Phi_{\mathfrak{p}}$ by $\mathcal{F} \Phi_{\mathfrak{p}}$ in the integral gives a function of $c$ depending only on $W_{\mathfrak{p}}$, namely
\begin{equation*}
  W_{\mathfrak{p}}^\vee(c) :=
  \lvert n \rvert_{\mathfrak{p}}^{-1/2}
  \int_{F_{\mathfrak{p}}} \psi_{\mathfrak{p}}(t)
  \left( \int_{F_{\mathfrak{p}}^\times}
    W_{\mathfrak{p}}(x) \{x,-c\}_{\mathfrak{p},n}
    \psi_{\mathfrak{p}}\!\left(-\frac{t^n}{cx}\right) \, d^\times x \right) \, d t
\end{equation*}
(Proposition \ref{proposition:transform-formula}).  Here $d t$ is self-dual for $\psi_{\mathfrak{p}}$ and $d^\times x = d x / \lvert x \rvert_{\mathfrak{p}}$.

Take such a test function at each place in $S$.  Suppose for simplicity that $S$ contains the places dividing $n$ and that $\psi_{\mathfrak{p}}$ is unramified outside $S$.  At each $\mathfrak{p} \notin S$, take $\Phi_{\mathfrak{p}}$ to be the basic vector.  Then $W_{\mathfrak{p}}$ is a generating function for Gauss sums of order $n$ (Lemma \ref{lemma:jacquet-basic}), and $W_{\mathfrak{p}}^\vee = W_{\mathfrak{p}}$ since $\mathcal{F} \Phi_{\mathfrak{p}} = \Phi_{\mathfrak{p}}$.  Write $W := \prod_{\mathfrak{p}} W_{\mathfrak{p}}$ and $W^\vee := \prod_{\mathfrak{p}} W_{\mathfrak{p}}^\vee$.  Theorem \ref{theorem:main} asserts that
\begin{equation}\label{eq:intro-formula}
  \sum_{c \in F^\times} W(c) = \sum_{c \in F^\times} W^\vee(c) + M(W),
\end{equation}
where the main term $M(W)$ is a linear combination of the products over $\mathfrak{p} \in S$ of the integrals $\int_{\beta_{\mathfrak{p}}} W_{\mathfrak{p}}(x) \lvert x \rvert_{\mathfrak{p}}^{1/n} \, d^\times x$, with $\beta_{\mathfrak{p}}$ running over the cosets of $F_{\mathfrak{p}}^{\times n}$ in $F_{\mathfrak{p}}^\times$.  We do not determine its coefficients.  If all these integrals vanish at one place, then $M(W) = 0$, and \eqref{eq:intro-formula} is the Fourier coefficient of \eqref{eq:intro-poisson-metaplectic}.

In \S\S\ref{sec:local-computations}--\ref{sec:quadratic-dual}, we write \eqref{eq:intro-formula} classically.  Suppose as in \eqref{eq:intro-theta-S} that $\mathcal{O}_S$ is a principal ideal domain.  Then the product of the factors $W_{\mathfrak{p}}(c)$ with $\mathfrak{p} \notin S$ vanishes unless $c \in \mathcal{O}_S$, in which case it is (Lemma \ref{lemma:product-of-basic-functions})
\begin{equation*}
  \tilde{g}(c) := \frac{1}{\lvert c \rvert_S} \sum_{\tau^n \mid c} \lvert \tau \rvert_S^{n-1} \, g(c / \tau^n),
  \qquad
  g(c) := \sum_{d \in (\mathcal{O}_S / c)^\times} \qr{d}{c}_n \psi_S(d / c).
\end{equation*}
Here $\psi_S := \prod_{\mathfrak{p} \in S} \psi_{\mathfrak{p}}$, and $\tau$ runs over the nonzero elements of $\mathcal{O}_S$, modulo units, whose $n$th power divides $c$; the sum over $\tau$ corresponds to that in \eqref{eq:intro-theta-S}.  Thus \eqref{eq:intro-formula} reads (Example \ref{example:simplest})
\begin{equation*}
  \sum_{c \in \mathcal{O}_S - \{0\}} W_S(c) \, \tilde{g}(c)
  =
  \sum_{c \in \mathcal{O}_S - \{0\}} W_S^\vee(c) \, \tilde{g}(c)
  + M(W),
\end{equation*}
where $W_S := \prod_{\mathfrak{p} \in S} W_{\mathfrak{p}}$ and $W_S^\vee := \prod_{\mathfrak{p} \in S} W_{\mathfrak{p}}^\vee$.  More generally, Theorem \ref{theorem:gauss-sums} treats the Gauss sums with $\psi_S(d / c)$ replaced by $\psi_S(b d / c)$, and twists of the sums over $c$ by characters at finitely many places outside $S$.  For odd $n$ and suitable $h \in \mathcal{O}_S$, Corollary \ref{corollary:quadratic-dual} shows that the sums twisted by $\qr{c}{h}_2 \qr{c}{h}_n^2$ have dual sums twisted by the quadratic character $\qr{c}{h}_2$, together with an explicit Hilbert symbol factor at the places in $S$.  For $n = 3$, the twist on the original sum is the sextic residue character.

Our interest in this example arose from the appearance of the same sextic-to-quadatic twist in recent work on the quasi-Riemann hypothesis \cite{OpenAI2026QuasiRHOctober5}.  There, with $\chi_{\mathfrak{p}}$ denoting the sextic residue character at each prime $\mathfrak{p}$ dividing $h$, the quadratic twist appears as the product $\chi_{\mathfrak{p}}^{-1} \chi_{\mathfrak{p}}^{-2} = \chi_{\mathfrak{p}}^3$ of the Fourier and theta factors \cite[(2.14)]{OpenAI2026QuasiRHOctober5}.  The summation formulas needed there are derived from the automorphy of the cubic theta function, using Patterson's formula for its Fourier coefficients \cite{Patterson1977CubicTheta} and a variant of the theta transformation identities established by Dunn and Radziwi\l\l{} \cite[\S5]{MR4816436}.

\subsection{Relation to the literature}\label{sec:literature}
Following Kubota, one obtains functional equations for Dirichlet series of Gauss sums by taking a Fourier coefficient of the functional equation for a metaplectic Eisenstein series.  See, for example, Diaconu \cite[(2.48)]{Diaconu_2004} and Brubaker and Bump \cite[Theorem 1]{MR2270571}.  Here we take the Fourier coefficient of the Poisson formula \eqref{eq:intro-poisson-boundary} to obtain the summation formula \eqref{eq:intro-formula} without passing through Dirichlet series.

In the quartic case, David, Dunn, Hamieh and Lin \cite{2023arXiv2306.11875} use Diaconu's functional equation to derive a Voronoi formula over \(\mathbb Q(i)\).  Ji \cite{2026arXiv2611.16665} applies the same functional equation to the quasi-Riemann hypothesis.  For cubic Gauss sums, the analytic properties of these Dirichlet series underlie the work of Heath-Brown and Patterson \cite{HeathBrownPatterson1979} and of Heath-Brown \cite{MR1815372} on their distribution at prime arguments.  For cubic Gauss sums over $\mathbb{Q}(\mu_3)$ twisted by a cubic character, Dunn and Radziwi\l\l{} \cite[\S5]{MR4816436} establish a Voronoi summation formula, generalizing a functional equation of Patterson \cite[Theorem~6.1]{Patterson1977CubicTheta}; Yoshimoto \cite{Yoshimoto1987} treats twists by non-cubic characters.  Both kinds of twist are covered by our results: a cubic twist can be absorbed into the frequency $b$, as in Remark \ref{remark:frequency-twist}, while Theorem \ref{theorem:gauss-sums} treats the twists by characters $\chi_{\mathfrak{p}}$ for which $\chi_{\mathfrak{p}}^3$ is ramified.

The local theory developed here rests on Deligne's account \cite{SB_1978-1979__21__244_0} of Kubota's cover, in which the principal series are realized on the torsor and the standard intertwining operator is an integral over lines.  That operator has been studied from the point of view of representation theory by several authors, over $p$-adic fields and (in some cases) for covers of general split reductive groups: see \cite{McNamara2012PrincipalSeries} for its action on spherical vectors (compare \S\ref{sec:intertwining}), \cite{Szpruch2019, GaoShahidiSzpruch2023} for its action on Whittaker functionals (compare \S\ref{sec:line-transform}), and \cite{MR3544420, Szpruch2019} for the determination of Plancherel measures (compare \S\ref{sec:inversion}).  On homogeneous genuine functions, our Fourier transform is a scalar multiple of the standard intertwining operator.

Kubota \cite{Kubota+1974+180+189} proves a summation formula for the sums $\sum_m c_m \Phi(m^{1/n})$, where the $c_m$ are the Fourier coefficients of his theta function on the $n$-fold cover of $\mathrm{SL}_2$, a residue of his Eisenstein series.  His formula is equivalent to the transformation law of that theta function under the Weyl element.  For $n = 3$ and $F = \mathbb{Q}(\mu_3)$, Patterson \cite{Patterson1977CubicTheta} expresses the $c_m$ in terms of cubic Gauss sums, so one expects Kubota's formula to be a specialization of \eqref{eq:intro-formula} in that case.  We do not verify this, but we check that Kubota's transform agrees with ours at the complex place after an explicit change of variables (Remark \ref{remark:kubota-transform}).  For general $n$, Kazhdan and Patterson \cite[pp.~36--37]{PMIHES_1984__59__35_0} relate Gauss sums of order $n$ to a theta function on the $n$-fold cover of $\mathrm{GL}_{n-1}$.

Sawin and Whitehead \cite[\S\S1.3, 1.5]{2025arXiv2507.08662} ask for functional equations for Dirichlet series over $\mathbb{F}_q[T]$ whose coefficients are Dirichlet characters multiplied by Gauss sums, as an input to their study of multiple Dirichlet series.  Theorem~\ref{theorem:gauss-sums} gives a summation formula for such coefficients (Remark~\ref{remark:sawin-whitehead}).  Taking a Mellin transform should give the desired functional equations.

\subsection{Idea of the proof}\label{sec:idea-of-the-proof}
As in the linear case, one would expect to recover \eqref{eq:intro-poisson-boundary} by Mellin inversion from the analytic properties of Kubota's Eisenstein series, established using Selberg's theory.  We record a direct proof in Appendix \ref{sec:proof-poisson}.  For $g \in \widetilde{\mathrm{SL}}_2(\mathbb{A})$, the difference
\begin{equation*}
   \sum_{v \in F^2 - \{0\}} (g \Phi)(v)
   -
   \sum_{v \in F^2 - \{0\}} (g \mathcal{F} \Phi)(v)
\end{equation*}
is an automorphic function of $g$.  Ordinary Poisson summation on $\mathbb{A}$ and on $\mathbb{A}^\times / \mathbb{A}^{\times n}$ gives its constant term explicitly in terms of $b \Phi$ and $b \mathcal{F} \Phi$ (Proposition \ref{proposition:constant-terms}).  If each of
$b \Phi$ and $b \mathcal{F} \Phi$ has a zero factor, then the constant term vanishes and the difference decays rapidly in the cusp; unfolding then shows that the difference is orthogonal to itself, hence vanishes.  For general $\Phi$, dilation at one finite place scales the two leading terms differently, allowing us to separate their contributions to the difference.

\subsection{Acknowledgments}\label{sec:acknowledgments}

We used the coding agents Codex and Claude extensively throughout this work.

% We are grateful to Valentin Blomer, Alex Dunn, Alex De Faveri, and Maksym Radziwill for helpful discussions.  


\section{The Hilbert symbol}\label{sec:hilbert-symbol}
Throughout \S\S\ref{sec:hilbert-symbol}--\ref{sec:germs}, let $F$ be a local field containing $\mu_n$; thus $F$ is non-archimedean, or $F = \mathbb{C}$, or $F = \mathbb{R}$ and $n \leq 2$.  Let $\psi$ be a nontrivial unitary character of $F$.  Let $d x$ denote the $\psi$-self-dual Haar measure on $F$, and $d^\times x := \frac{d x}{\lvert x \rvert}$ the corresponding multiplicative Haar measure.  We write $\lvert \cdot \rvert$ for the normalized absolute value on $F$, which for $F = \mathbb{C}$ is the square of the usual one.  When $F$ is non-archimedean, we write $\mathfrak{o}$ for its ring of integers, $\varpi$ for a uniformizer, $v$ for the valuation, and $q$ for the residue field cardinality.

We denote by $\{x, y\}_n \in \mu_n$ the Hilbert symbol of $x, y \in F^\times$ \cite[Ch.~V, \S3]{MR1697859}.  We regard it as complex-valued via a fixed embedding $\mu_n \hookrightarrow \mathbb{C}^\times$.  (The normalizations of the Hilbert symbol in the literature differ by inversion, which may be absorbed into the choice of embedding.)  We use the following standard properties.
\begin{enumerate}[(a)]
\item\label{item:hilbert-bimultiplicative} The symbol is bimultiplicative and continuous, and $\{x, y\}_n \{y, x\}_n = 1$.
\item\label{item:hilbert-powers} We have $\{x, y\}_n = 1$ for all $x \in F^\times$ if and only if $y \in F^{\times n}$.
\item\label{item:hilbert-units} If $F$ is non-archimedean of residue characteristic not dividing $n$, then $\{x, y\}_n = 1$ for $x, y \in \mathfrak{o}^\times$.
\item\label{item:hilbert-steinberg} We have $\{x, y\}_n = 1$ when $x + y = 1$ (the \emph{Steinberg relation}), and $\{x, -x\}_n = 1$ for all $x$.
\item\label{item:hilbert-open} The subgroup $F^{\times n}$ of $F^\times$ is open.
\item\label{item:hilbert-tame} If $F$ is non-archimedean and of residue characteristic not dividing $n$, then $1 + \varpi \mathfrak{o} \subset F^{\times n}$, and $u \mapsto \{\varpi, u\}_n$ is a character of $\mathfrak{o}^\times$ of order $n$.
\end{enumerate}
By properties \eqref{item:hilbert-bimultiplicative} and \eqref{item:hilbert-powers}, the symbol $\{x, y\}_n$ is trivial if either argument lies in $F^{\times n}$.

For an integrable function $g$ on $F$, we write $\hat{g}(y) := \int_F g(x) \psi(x y) \, d x$.  We equip the open subgroup $F^{\times n}$ of $F^\times$ with the restriction of the Haar measure $d^\times x$.  The map $\mu \mapsto \mu^n$ from $F^\times$ onto $F^{\times n}$ is $n$-to-one and multiplies $d^\times \mu$ by $\lvert n \rvert$, so for $\phi \in L^1(F^{\times n})$,
\begin{equation}\label{eq:haar-on-powers}
  \int_{F^{\times n}} \phi(a) \, d^\times a = \frac{\lvert n \rvert}{n} \int_{F^\times} \phi(\mu^n) \, d^\times \mu.
\end{equation}

\section{Genuine functions}\label{sec:genuine}
We denote by
\begin{equation*}
  U := F^2 - \{0\}
\end{equation*}
the punctured plane, with coordinates $(x, y)$, and by $U_X$ and $U_Y$ the open subsets on which $x \neq 0$ and $y \neq 0$, respectively.  For $v = (x, y)$, write $\lVert v \rVert := \max(\lvert x \rvert, \lvert y \rvert)$.

Deligne's $\mu_n$-torsor over $U$ is obtained by gluing the trivial torsors on $U_X$ and $U_Y$ via the Hilbert symbol \cite[\S\S1.12, 3.7]{SB_1978-1979__21__244_0}.  We describe sections of the associated line bundle, or equivalently functions on the torsor equivariant for $\mu_n \hookrightarrow \mathbb{C}^\times$, in these two charts.

\begin{definition}\lean{MetaplecticPoisson.Glues}
  By a \emph{genuine function} on $U$, we mean a pair $\Phi = (\Phi_X, \Phi_Y)$ of smooth functions $\Phi_X : U_X \rightarrow \mathbb{C}$ and $\Phi_Y : U_Y \rightarrow \mathbb{C}$ such that
  \begin{equation}\label{eq:transition-relation}
    \Phi_X(x, y) = \{x, y\}_n \Phi_Y(x, y) \quad \text{ on } U_X \cap U_Y.
  \end{equation}
  We denote by $\mathcal{D}$ the following space of genuine functions $\Phi$.\lean{MetaplecticPoisson.Datum.testSet}
  \begin{itemize}
  \item If $F$ is non-archimedean, we require that $\Phi_X$ and $\Phi_Y$ vanish outside some compact subset of $U$.
  \item If $F$ is archimedean, we require that $\Phi_X$ and $\Phi_Y$, together with their partial derivatives of all orders with respect to the real coordinates, decay rapidly at the origin and at infinity: each is $O(\min(\lVert v \rVert^{N}, \lVert v \rVert^{-N}))$ for each $N \geq 0$.
  \end{itemize}
\end{definition}
For $a \in F^{\times n}$, the dilate $\Phi(a \, \cdot) := (\Phi_X(a x, a y), \Phi_Y(a x, a y))$ of a genuine function is genuine, by property \eqref{item:hilbert-powers} of the Hilbert symbol.

For a genuine function $\Phi$, the functions $\lvert \Phi_X \rvert$ and $\lvert \Phi_Y \rvert$ agree on $U_X \cap U_Y$; we denote by $\lvert \Phi \rvert$ the resulting function on $U$.

\section{The Fourier transform and the Schwartz space}\label{sec:transform}
\subsection{Integrals over lines}\label{sec:lines}
For $n = 1$, a genuine function is a function $\Phi$ on $U$, and one may compute its symplectic Fourier transform at $(x', y') \neq(0, 0)$ by first integrating $\Phi$ over the line $x y' - y x' = \tau$, then taking the Fourier transform in the variable $\tau$ at $1$.  For homogeneous $\Phi$, integrating over $x y ' - y x ' = 1$ gives the standard intertwining operator  \cite[(1.4) and Example~1.7]{MR1988971}.

Deligne \cite[\S\S1.13--1.14, 6.3]{SB_1978-1979__21__244_0} uses the trivialization of his torsor along $a x + b y = 1$ to realize the standard intertwining operator for the metaplectic principal series as a Radon transform.  Dilation by $\tau^n$ transports this trivialization to $a x + b y = \tau^n$ for $\tau \neq 0$.  Writing $(a, b) = (y', -x')$ in his formulas leads to the following definition:
\begin{definition}\label{definition:lines}\lean{MetaplecticPoisson.radon}
  Let $\Phi$ be a genuine function.  For $\tau \in F$ and $v' = (x', y') \in U$, we define $R_X \Phi(v', \tau)$ and $R_Y \Phi(v', \tau)$ by the following integrals whenever they converge absolutely:
  \begin{align*}
    R_X \Phi(v', \tau)
    &:=
      \int_F \Phi_X \left( x, \frac{x y' - \tau^n}{x'} \right) \{x, -x'\}_n
      \,
      \frac{ d x}{\lvert x' \rvert}
    && \text{if } x' \neq 0, \\
    R_Y \Phi(v', \tau)
    &:=
      \int_F \Phi_Y \left( \frac{x' y + \tau^n}{y'}, y \right) \{y, y'\}_n \{-1, y'\}_n \, 
      \frac{d y}{\lvert y' \rvert}
    && \text{if } y' \neq 0.
  \end{align*}
\end{definition}

The two formulas use the measures $d x / \lvert x' \rvert$ and $d y / \lvert y' \rvert$ on the line $x y' - y x' = \tau^n$, each obtained by disintegrating $d x \, d y$ with respect to $x y' - y x'$.  Write $(a, b) := (y', -x')$, so that the line is defined by $a x + b y = \tau^n$.  The above integrands are then
\begin{equation*}
  I_X(x) := \Phi_X \left( x, y \right) \{x, b\}_n,
  \quad I_Y(y) := \Phi_Y(x, y) \{y, a\}_n \{- 1, a\}_n.
\end{equation*}
When $a b x y \tau \neq 0$, the Steinberg relation applied to $a x / \tau^n$ and $b y / \tau^n$ gives $\{a x, b y\}_n = 1$.  Using also that $\{b, a\}_n = \{- 1, a\}_n \{x', y'\}_n$, we obtain
\begin{equation}\label{eq:symbol-on-lines}
  I_X(x) = \{x', y'\}_n I_Y(y).
\end{equation}
Near $x = 0$ with $y \neq 0$, the quantity $b y = \tau^n - a x$ lies in the open subgroup $F^{\times n}$, so $I_X(x) = \{x, b y\}_n \Phi_Y(x, y) = \Phi_Y(x, y)$ for $x \neq 0$.  This extends $I_X$ smoothly across $x = 0$, and likewise $I_Y$ extends smoothly across $y = 0$ away from the origin.  The Hilbert symbol factors in Definition~\ref{definition:lines} may be motivated by the continuity of the integrands, the transition relation \eqref{eq:transition-relation} (see Lemma~\ref{lemma:lines}\eqref{item:lines-compatibility} below), and the special case
\begin{equation*}
  R_Y \Phi((0, 1), \tau) = \int_{F} \Phi_Y(\tau^n, y) \, d y.
\end{equation*}

\begin{lemma}\label{lemma:lines}
  Let $\Phi$ be a genuine function satisfying
  \begin{equation}\label{eq:growth}
    \lvert \Phi(v) \rvert \ll_N \min \left( \lVert v \rVert^{-1 + 1/n}, \lVert v \rVert^{-N} \right)
    \qquad (N \geq 1).
  \end{equation}
  \begin{enumerate}[(i)]
  \item\label{item:lines-convergence} The integrals defining $R_X \Phi(v', \tau)$ and $R_Y \Phi(v', \tau)$ converge absolutely and define continuous functions on $U_X \times F$ (resp.\ $U_Y \times F$).  They decay rapidly in $\tau$, uniformly for $v'$ in a compact subset of $U_X$ (resp.\ $U_Y$).
  \item\label{item:lines-compatibility} If $x' y' \neq 0$, then $R_X \Phi(v', \tau) = \{x', y'\}_n R_Y \Phi(v', \tau)$.
  \item\label{item:lines-smooth}\lean{MetaplecticPoisson.Datum.radon_smooth_real}\lean{MetaplecticPoisson.Datum.radon_smooth_complex} If $\Phi \in \mathcal{D}$, then $R_X \Phi$ is smooth on $U_X \times F$, and $R_Y \Phi$ is smooth on $U_Y \times F$.  For $v'$ in a compact subset of $U_X$ (resp.\ $U_Y$), the functions $R_X \Phi(v', \cdot)$ (resp.\ $R_Y \Phi(v', \cdot)$) form a bounded subset of the space of Schwartz--Bruhat functions on $F$.  If $F$ is archimedean, the same holds for all derivatives of $R_X \Phi$ (resp.\ $R_Y \Phi$).
  \end{enumerate}
\end{lemma}
\begin{proof}
  \eqref{item:lines-convergence} We treat $R_X$.  Let $K$ be a compact subset of $U_X$, and let $v'$ range over $K$.  For $v = (x, y)$ on the line $x y' - y x' = \tau^n$, we have
  \begin{equation}\label{eq:norm-on-lines}
    \lVert v \rVert \asymp_K \max(\lvert x \rvert, \lvert \tau \rvert^n).
  \end{equation}
  Indeed, $x'$ and $y'$ are bounded on $K$, and $x'$ is bounded away from zero, so the line equation bounds $\lvert \tau \rvert^n$ in terms of $\lVert v \rVert$, and $\lvert y \rvert$ in terms of $\max(\lvert x \rvert, \lvert \tau \rvert^n)$.  For $\lVert v \rVert \leq 1$, both $x$ and $\tau$ are bounded and the integrand is $O(\lvert x \rvert^{-1+1/n})$.  For $\lVert v \rVert > 1$, the integrand is
  \begin{equation*}
    O((1+\lvert x \rvert)^{-2}(1+\lvert \tau \rvert)^{-N}) \qquad(N \geq 1)
  \end{equation*}
  by \eqref{eq:growth}.  Integration gives the bound, and dominated convergence gives continuity since the integrand is continuous in $(v',\tau)$ for $x \neq 0$.

  \eqref{item:lines-compatibility} For $\tau \neq 0$, we integrate \eqref{eq:symbol-on-lines}, using that $d x / \lvert x' \rvert = d y / \lvert y' \rvert$ on the line.  The case $\tau = 0$ follows by continuity.

  \eqref{item:lines-smooth} We treat $R_X$.  Let $\Phi \in \mathcal{D}$.  We first check that the integrand defining $R_X \Phi(v', \tau)$ extends to a smooth function of $(x, v', \tau) \in F \times U_X \times F$.  For $x \neq 0$, it is smooth, also when $\tau = 0$, because $\Phi_X$ is smooth on $U_X$ and $\{x, -x'\}_n$ is locally constant.  At a point with $x = 0$ and $\tau \neq 0$, the identity $I_X(x) = \Phi_Y(x, y)$ holds for all nearby $(x, v', \tau)$ with $x \neq 0$, as observed after Definition \ref{definition:lines}.  This gives the smooth extension there.  At the points with $x = \tau = 0$, the line passes through the origin, and we use that $\Phi \in \mathcal{D}$.  Let $v'$ range over $K$, as above.  If $F$ is non-archimedean, then $\Phi$ vanishes near the origin, so by \eqref{eq:norm-on-lines}, the integrand vanishes for $(x, \tau)$ in a neighborhood of $(0, 0)$ that does not depend upon $v'$.  If $F$ is archimedean, then by the chain rule, each derivative of the integrand with respect to $(x, v', \tau)$ is a combination of derivatives of $\Phi_X$ at $v$, with coefficients that are polynomial in $(x, \tau)$ and smooth in $v'$.  Set $m := \max(\lvert x \rvert, \lvert \tau \rvert^n)$.  By the definition of $\mathcal{D}$ and \eqref{eq:norm-on-lines}, each such derivative is $O(\min(m^N, m^{-N}))$ for each $N \geq 0$, uniformly in $v'$.  The integrand thus extends by zero to a smooth function near the points with $x = \tau = 0$.

  The integrand is in fact Schwartz--Bruhat in $(x, \tau)$, uniformly in $v'$.  If $F$ is non-archimedean, then since $\Phi$ has compact support, \eqref{eq:norm-on-lines} shows that the integrand vanishes for $(x, \tau)$ outside a fixed compact set.  Being locally constant in $(x, v', \tau)$, it is then, by compactness, invariant under a fixed open subgroup of translations of $(x, \tau)$, and locally constant in $v'$ uniformly in $(x, \tau)$.  If $F$ is archimedean, the bound just given shows that the integrand and its derivatives decay rapidly in $(x, \tau)$.  Integration in $x$ carries Schwartz--Bruhat functions on $F^2$ continuously to Schwartz--Bruhat functions on $F$, giving \eqref{item:lines-smooth}.
\end{proof}

\subsection{Normalization of the intertwining operator}\label{sec:intertwining}
To motivate our definitions of both the Fourier transform and the basic vector, we record how the standard intertwining operator acts on normalized spherical vectors in the principal series.

Let $\xi$ be a unitary character of $F^\times$, trivial on $\mu_n$, and let $0 < \Re(s) < 1$.  Suppose that a genuine function $f$ satisfies
\begin{equation}\label{eq:principal-series-homogeneity}
  f(\mu^n v) = \xi(\mu)^{-1} \lvert \mu \rvert^{-n-s} f(v)
  \qquad (\mu \in F^\times,\ v \in U).
\end{equation}
Write $R$ for $R_X$ or $R_Y$.  For $\tau \neq 0$, the line integrals $R f(v', \tau)$ converge absolutely.  Set
\begin{equation*}
  A f := R f(\cdot, 1).
\end{equation*}
Thus $A$ is the standard intertwining operator in Deligne's model \cite[\S6.3]{SB_1978-1979__21__244_0}.  The substitution $x \mapsto \tau^n x$ gives
\begin{equation}\label{eq:radon-homogeneity}
  R f(v', \tau) = \xi(\tau)^{-1} \lvert \tau \rvert^{-s} A f(v') \qquad (\tau \neq 0).
\end{equation}

Suppose now that $F$ is non-archimedean, of residue characteristic not dividing $n$, and that $\psi$ is unramified.  Let $\Phi^1$ denote the genuine function supported on the primitive vectors of $\mathfrak{o}^2$, with $\Phi^1_X(x, y) = 1$ when $x$ is a unit and $\Phi^1_Y(x, y) = 1$ when $y$ is a unit; the two conditions are compatible by property \eqref{item:hilbert-units} of the Hilbert symbol.  For a genuine function $\Phi$ invariant under dilation by $n$th powers of units, we define its Mellin coefficient at $s \in \mathbb{C}$ by
\begin{equation*}
  \Phi[s](v) := \sum_{k \in \mathbb{Z}} q^{k (n + s)} \, \Phi(\varpi^{-n k} v)
\end{equation*}
whenever the sum converges.  The spherical section $e_s := \Phi^1[s]$ satisfies \eqref{eq:principal-series-homogeneity} with $\xi = 1$ and is normalized by $(e_s)_Y(0, 1) = 1$.  Write $L(s) := (1 - q^{-s})^{-1}$.

\begin{lemma}\label{lemma:spherical-axis}
  For $\Re(s) > 0$, we have $(A e_s)_Y(0, 1) = L(s) / L(1 + s)$.
\end{lemma}
\begin{proof}
  At $(0, 1)$ the value of $A e_s$ is the integral of $(e_s)_X(1, y)$ over $y \in F$.  The contribution from $\mathfrak{o}$ is $1$.  The $y$ of valuation $- n k$, with $k \geq 1$, form a set of measure $q^{n k} (1 - q^{-1})$, on which $(e_s)_X(1, y) = q^{- k (n + s)}$; the remaining $y$ contribute zero.  The total is
  \begin{equation*}
    1 + (1 - q^{-1}) \sum_{k \geq 1} q^{- k s} = \frac{1 - q^{-1-s}}{1 - q^{-s}}. \qedhere
  \end{equation*}
\end{proof}
This gives the value at $(0, 1)$ of the identity
\begin{equation}\label{eq:spherical}
  A e_s = \frac{L(s)}{L(1 + s)} \, e_{-s},
\end{equation}
which we verify on all of $U$ in Lemma \ref{lemma:spherical}.  Compare with the formula of McNamara \cite[Theorem~6]{McNamara2012PrincipalSeries}, which concerns a renormalized intertwining operator, under the assumption $2 n \mid q - 1$.

For $n = 1$, the characteristic function of $\mathfrak{o}^2$ is the sum of the functions $\Phi^1(\varpi^{-k} \, \cdot)$ over $k \geq 0$ (cf.\ the proof of \cite[Theorem~3.13(3)]{MR1694894}), and its Mellin coefficient at $s$ is $L(1 + s) e_s$.  For general $n$, the Mellin coefficient of the dilate $\Phi^1(\varpi^{-n k} \, \cdot)$ is $q^{-k (n + s)} e_s$.  To obtain the same factor $L(1 + s)$, the term with index $k$ should contribute $q^{-k (1 + s)} e_s$.  We therefore set
\begin{equation*}
  \Phi^0 := \sum_{k \geq 0} q^{(n-1) k} \, \Phi^1(\varpi^{- n k} \, \cdot),
\end{equation*}
so that $\Phi^0[s] = L(1 + s) e_s$.  We call $\Phi^0$ the \emph{basic vector}, and study it in \S\ref{sec:basic-vector}.

As in the linear case, we seek a Fourier transform that reverses dilations and fixes $\Phi^0$.  Such a transform should carry $\Phi^0[s]$ to $\Phi^0[-s] = L(1 - s) e_{-s}$.  By \eqref{eq:spherical}, $A$ carries $\Phi^0[s]$ to $L(s) e_{-s}$.  A one-variable Fourier transform in $\tau$ supplies the missing factor $L(1 - s) / L(s)$.  Indeed, for $f$ satisfying \eqref{eq:principal-series-homogeneity}, we have by \eqref{eq:radon-homogeneity} and Tate's local functional equation \cite{MR0217026}
\begin{equation*}
  \int_F \psi(\tau) R f(v', \tau) \, d \tau
  = A f(v') \int_F \psi(\tau) \xi(\tau)^{-1} \lvert \tau \rvert^{-s} \, d \tau
  = \gamma(s, \xi, \psi) A f(v'),
\end{equation*}
where the Fourier integrals are understood distributionally at frequency $1$.  Since $\gamma(s, 1, \psi) = L(1 - s) / L(s)$ for unramified $\psi$, this gives
\begin{equation}\label{eq:cuue7ixq1x}
  \int_F \psi(\tau) R(\Phi^0[s])(v', \tau) \, d \tau
  = \frac{L(1 - s)}{L(s)} \, L(s) e_{-s}(v')
  = \Phi^0[-s](v').
\end{equation}

\subsection{The Fourier transform}\label{sec:fourier-transform}

We now apply the integral on the left-hand side of \eqref{eq:cuue7ixq1x} to define the Fourier transform of any genuine function $\Phi$ satisfying \eqref{eq:growth}.
\begin{definition}\label{definition:fourier-transform}\lean{MetaplecticPoisson.transform}
  We define functions $(\mathcal{F} \Phi)_X$ on $U_X$ and $(\mathcal{F} \Phi)_Y$ on $U_Y$ by
  \begin{equation}\label{eq:transform-via-lines}
    \begin{aligned}
      (\mathcal{F} \Phi)_X(v') &:= \lvert n \rvert^{-1/2} \int_F \psi(\tau) R_X \Phi(v', \tau) \, d \tau,
      \\
      (\mathcal{F} \Phi)_Y(v') &:= \lvert n \rvert^{-1/2} \int_F \psi(\tau) R_Y \Phi(v', \tau) \, d \tau.
    \end{aligned}
  \end{equation}
\end{definition}
The integrals converge absolutely by Lemma \ref{lemma:lines}.  For $n = 1$, we recover the symplectic Fourier transform.  For $f$ satisfying \eqref{eq:principal-series-homogeneity} with $0 < \Re(s) < 1$, the distributional interpretation gives $\mathcal{F} f = \lvert n \rvert^{-1/2} \gamma(s, \xi, \psi) A f$.

\begin{proposition}\label{proposition:fourier-transform-genuine}
  The functions $(\mathcal{F} \Phi)_X$ and $(\mathcal{F} \Phi)_Y$ are continuous, and $(\mathcal{F} \Phi)_X = \{x', y'\}_n (\mathcal{F} \Phi)_Y$ on $U_X \cap U_Y$.  If $\Phi \in \mathcal{D}$, then $(\mathcal{F} \Phi)_X$ and $(\mathcal{F} \Phi)_Y$ are smooth.\lean{MetaplecticPoisson.Datum.arch_transform_smooth}
\end{proposition}
\begin{proof}
  This follows from Lemma \ref{lemma:lines} and the definition \eqref{eq:transform-via-lines}, by dominated convergence.
\end{proof}
We write $\mathcal{F} \Phi := ((\mathcal{F} \Phi)_X, (\mathcal{F} \Phi)_Y)$.  It is smooth, and thus a genuine function, if $\Phi \in \mathcal{D}$.  The factor $\lvert n \rvert^{-1/2}$, which equals one at places of residue characteristic prime to $n$, makes $\mathcal{F}$ unitary (Theorem \ref{theorem:schwartz-space}).

We record the integral kernel of $\mathcal{F}$.
\begin{lemma}\label{lemma:root-kernel}
  Write $[v, w] := x y' - y x'$ for $v = (x, y)$ and $w = (x', y')$.  For $v, w \in U_X$ with $[v, w] \neq 0$, set
  \begin{equation*}
    K(w, v) := \{x, -x'\}_n \, \lvert n \rvert^{-3/2} \sum_{\tau^n = [v, w]} \psi(\tau) \, \lvert \tau \rvert^{1-n},
  \end{equation*}
  where $\tau$ runs over the $n$th roots of $[v, w]$ in $F$.  Then $K(w, v) \ll \lvert [v, w] \rvert^{-1 + 1/n}$, and for each genuine function $\Phi$ satisfying \eqref{eq:growth} and each $w \in U_X$,
  \begin{equation}\label{eq:root-kernel}
    (\mathcal{F} \Phi)_X(w) = \int_{U_X} \Phi_X(v) K(w, v) \, d v,
  \end{equation}
  where the integral converges absolutely.
\end{lemma}
\begin{proof}
  The bound holds because $\lvert \tau \rvert^n = \lvert [v, w] \rvert$.  In the definition of $(\mathcal{F} \Phi)_X(w)$, the double integral over $x$ and $\tau$ converges absolutely, by the estimates in the proof of Lemma \ref{lemma:lines}\eqref{item:lines-convergence}.  Fix $x \neq 0$, and change variables from $\tau$ to $y = (x y' - \tau^n) / x'$.  Then each $y$ with $[v, w] \in F^{\times n}$ arises from $n$ values of $\tau$, and $d \tau \, d x / \lvert x' \rvert = \lvert n \tau^{n-1} \rvert^{-1} \, d x \, d y$.  This gives \eqref{eq:root-kernel}, since the locus $[v, w] = 0$ has measure zero.
\end{proof}

\begin{remark}\label{remark:half-density}
  The weight $\lvert n \rvert^{-3/2} \lvert \tau \rvert^{1-n}$ in $K$ is the square root of the absolute value of a Hessian determinant.  Indeed, write $u := [v, w]$, and let $\tau = f(u)$ be a local branch of the $n$th root.  Since $u f'(u) = f(u) / n$, we have
  \begin{equation*}
    \det \left( \partial_v \partial_w f(u) \right)
    = f'(u) \bigl( f'(u) + u f''(u) \bigr)
    = \frac{f'(u)^2}{n},
  \end{equation*}
  and $f'(u) = (n \tau^{n-1})^{-1}$.  Thus
  \begin{equation*}
    K(w, v) = \{x, -x'\}_n \sum_{\tau^n = [v, w]} \psi(\tau) \, \bigl\lvert \det ( \partial_v \partial_w \tau ) \bigr\rvert^{1/2}.
  \end{equation*}
  For $n = 1$, the determinant is $1$, and the kernel is that of the symplectic Fourier transform.  For general $n$, the factor $\lvert n \rvert^{-1/2}$ in Definition \ref{definition:fourier-transform} is precisely what makes the weight of each root equal to the square root of its Hessian determinant.

  This is the weight attached to a phase function in the theory of Fourier integral operators.  For a function $\phi$ of $(v, w)$, the measure $\lvert \det(\partial_v \partial_w \phi) \rvert \, d v \, d w$ is unchanged by changes of coordinates in $v$ and in $w$.  The kernel $\psi(\phi) \, \lvert \det(\partial_v \partial_w \phi) \rvert^{1/2}$ thus defines an operator on half-densities that depends only upon $\phi$.

  To see how this weight enters Fourier inversion, fix $v$ and compare the phases at $v$ and $v + z$ on the same branch.  The linear term in $\tau(v + z, w) - \tau(v, w)$ is $\eta \cdot z$, where $\eta := \partial_v \tau$.  In the composition $\mathcal{F}^* \mathcal{F}$, the product of the two weights at $z = 0$ is $\lvert \det(\partial_w \eta) \rvert$.  Thus the change of variables $w \mapsto \eta$ gives $d \eta$ as the measure for the linearized phase $\eta \cdot z$, as in ordinary Fourier inversion.

  To obtain the whole integral over $\eta$, we need all the roots.  Indeed, $n \tau^{n-1} \eta = (y', -x')$, whence $\tau = n \, \eta \cdot v$.  Each $\eta$ with $\eta \cdot v \neq 0$ therefore determines exactly one pair $(w, \tau)$, and the excluded line has measure zero.  This accounts for the normalization at the level of the linearized phase.  Theorem \ref{theorem:schwartz-space} gives the exact unitarity statement.
\end{remark}

We record how $\mathcal{F}$ interacts with the dilations by $F^{\times n}$, and with a change of the character $\psi$; understanding the latter will allow us to reduce some computations to the case of a convenient $\psi$.

\begin{lemma}\label{lemma:transform-dilation}
  \begin{enumerate}[(i)]
  \item\label{item:transform-dilation} For $a \in F^{\times n}$, we have $\mathcal{F}(\Phi(a \, \cdot)) = \lvert a \rvert^{-2} (\mathcal{F} \Phi)(a^{-1} \, \cdot)$.
  \item\label{item:transform-character} For $b \in F^\times$, let $\mathcal{F}_{\psi_b}$ denote the transform defined using the character $\psi_b := \psi(b \, \cdot)$ and the $\psi_b$-self-dual measures.  Then $\mathcal{F}_{\psi_b} \Phi = \lvert b \rvert^n (\mathcal{F} \Phi)(b^n \, \cdot)$.
  \end{enumerate}
\end{lemma}
\begin{proof}
  We verify both identities for the $X$-components; the $Y$-components are treated in the same way.

  \eqref{item:transform-dilation} The substitution $x \mapsto x / a$ in Definition \ref{definition:lines} gives $R_X (\Phi(a \, \cdot))(v', \tau) = \lvert a \rvert^{-2} R_X \Phi(a^{-1} v', \tau)$; the symbol $\{x, -x'\}_n$ is unchanged, since $a \in F^{\times n}$.  We conclude by integrating against $\psi(\tau)$.

  \eqref{item:transform-character} The $\psi_b$-self-dual measure is $\lvert b \rvert^{1/2} \, d x$.  It occurs once in the line integral and once in the integral over $\tau$, so $(\mathcal{F}_{\psi_b} \Phi)_X(v') = \lvert n \rvert^{-1/2} \lvert b \rvert \int_F \psi(b \tau) R_X \Phi(v', \tau) \, d \tau$.  The substitution $\tau \mapsto \tau / b$ then gives $\lvert b \rvert^n (\mathcal{F} \Phi)_X(b^n v')$, because $R_X \Phi(v', \tau / b) = \lvert b \rvert^n R_X \Phi(b^n v', \tau)$.
\end{proof}

We next observe that along each orbit of dilations by $F^{\times n}$, the function $\mathcal{F} \Phi$ is given by an ordinary Fourier transform in an $n$th root variable.
\begin{lemma}\label{lemma:transform-fourier}
  Let $v' \in U_X$ and $\mu \in F^\times$.  Then
  \begin{equation*}
    (\mathcal{F} \Phi)_X(\mu^n v') = \lvert n \rvert^{-1/2} \lvert \mu \rvert^{1 - n} \int_F R_X \Phi(v', \tau) \psi(\mu \tau) \, d \tau.
  \end{equation*}
  The same holds with $Y$ in place of $X$.
\end{lemma}
\begin{proof}
  We have $R_X \Phi(\mu^n v', \mu \tau) = \lvert \mu \rvert^{-n} R_X \Phi(v', \tau)$.  We substitute $\tau \mapsto \mu \tau$ in \eqref{eq:transform-via-lines}.
\end{proof}

\begin{corollary}\label{corollary:transform-decay}
  Let $\Phi \in \mathcal{D}$.  Then $\mathcal{F} \Phi$ satisfies \eqref{eq:growth}.  If $F$ is non-archimedean, then $\mathcal{F} \Phi$ vanishes outside a compact subset of $F^2$.
\end{corollary}
\begin{proof}
  Choose compact subsets $K_X \subset U_X$ and $K_Y \subset U_Y$ whose union meets every orbit of $F^{\times n}$ in $U$.  Writing $v = \mu^n v_0$ with $v_0 \in K_X \cup K_Y$, we have $\lVert v \rVert \asymp \lvert \mu \rvert^n$.  By Lemma \ref{lemma:transform-fourier}, it suffices to bound the Fourier transform of $R_X \Phi(v_0, \cdot)$ by $O((1 + \lvert \mu \rvert)^{-N})$ uniformly for $v_0 \in K_X$, and likewise in the $Y$-chart.  In the non-archimedean case we also need these Fourier transforms to vanish outside a fixed compact set.  Both assertions follow from the boundedness in the Schwartz--Bruhat space given by Lemma \ref{lemma:lines}.
\end{proof}

\subsection{The Schwartz space}\label{sec:schwartz-space}
The space $\mathcal{D}$ is not preserved by $\mathcal{F}$.  Following \cite[\S3.1]{MR1694894}, we enlarge it by the transforms of its elements.
\begin{definition}\label{definition:schwartz-space}\lean{MetaplecticPoisson.Datum.schwartz}
  We define the \emph{Schwartz space} by
  \begin{equation*}
    \mathcal{S} := \mathcal{D} + \mathcal{F} \mathcal{D}.
  \end{equation*}
\end{definition}
By Proposition \ref{proposition:fourier-transform-genuine} and Corollary \ref{corollary:transform-decay}, the elements of $\mathcal{S}$ are genuine functions satisfying \eqref{eq:growth}.\lean{MetaplecticPoisson.Datum.schwartz_smooth_real}\lean{MetaplecticPoisson.Datum.schwartz_smooth_complex}  The integral defining $\mathcal{F}$ thus applies to every element of $\mathcal{S}$.

\begin{lemma}\label{lemma:schwartz-dilation}
  The space $\mathcal{S}$ is preserved by the dilations by $F^{\times n}$, and does not depend upon the choice of $\psi$.
\end{lemma}
\begin{proof}
  The dilations by $F^{\times n}$ preserve $\mathcal{D}$, hence also $\mathcal{F} \mathcal{D}$, by Lemma \ref{lemma:transform-dilation}.  Each nontrivial unitary character of $F$ is of the form $\psi_b$ with $b \in F^\times$, and by the same lemma, the transform defined using $\psi_b$ maps $\mathcal{D}$ onto the dilate of $\mathcal{F} \mathcal{D}$ by $b^n$.
\end{proof}

\section{Symmetries}\label{sec:symmetries}
\subsection{The metaplectic action}\label{sec:metaplectic}
Deligne \cite[\S\S1.7, 1.9, 1.11--1.12, 3.7]{SB_1978-1979__21__244_0} lifts the action of $\mathrm{SL}_2(F)$ on the punctured plane to an action of Kubota's metaplectic group \cite{MR0255490} on the torsor.  The lift of each unipotent matrix preserves one of the two trivializations \cite[(1.14.3)]{SB_1978-1979__21__244_0}.  We let matrices act on row vectors from the right.  In terms of the components of a genuine function $\Phi$, the lifts $n(u)$ and $n^-(u)$ of the upper and lower unipotent matrices with entry $u \in F$ are then given by
\begin{equation}\label{eq:unipotent-charts}
  (n(u) \Phi)_X(x, y) = \Phi_X(x, y + u x),
  \qquad
  (n^-(u) \Phi)_Y(x, y) = \Phi_Y(x + u y, y).
\end{equation}
A genuine function is determined by either of its components, by continuity.  The transition relation \eqref{eq:transition-relation} gives the other component of $n(u) \Phi$: for $x \neq 0$ and $y (y + u x) \neq 0$,
\begin{equation}\label{eq:unipotent-other-chart}
  (n(u) \Phi)_Y(x, y) = \{x, 1 + u x / y\}_n \Phi_Y(x, y + u x).
\end{equation}
Since $F^{\times n}$ is open, the symbol is trivial for $(u, x, y)$ near any given $(u_0, 0, y_0)$ with $y_0 \neq 0$, so that $(n(u) \Phi)_Y(x, y) = \Phi_Y(x, y + u x)$ there.  The same applies to $n^-(u)$ near the line $y = 0$.  In particular,
\begin{equation}\label{eq:unipotent-axis}
  (n(u) \Phi)_Y(0, y) = \Phi_Y(0, y),
  \qquad
  (n^-(u) \Phi)_X(x, 0) = \Phi_X(x, 0).
\end{equation}
In these terms, Deligne's construction \cite[\S\S1.11--1.12, 3.7]{SB_1978-1979__21__244_0} takes the following form:
\begin{proposition}\label{proposition:kubota}\lean{MetaplecticPoisson.Datum.kubota_projection}
  Let $\widetilde{\mathrm{SL}}_2(F)$ denote the group of automorphisms of the space of genuine functions generated by the scalars in $\mu_n$ and by the operators $n(u)$ and $n^-(u)$ with $u \in F$.  There is a surjective homomorphism $\pi : \widetilde{\mathrm{SL}}_2(F) \rightarrow \mathrm{SL}_2(F)$ with
  \begin{equation*}
    \pi(n(u)) = \begin{pmatrix} 1 & u \\ 0 & 1 \end{pmatrix},
    \qquad
    \pi(n^-(u)) = \begin{pmatrix} 1 & 0 \\ u & 1 \end{pmatrix},
  \end{equation*}
  whose kernel is $\mu_n$.
\end{proposition}

The unipotent lifts commute with the line integrals, hence with $\mathcal{F}$, in the following sense:
\begin{lemma}\label{lemma:unipotent-properties}
  Let $\Phi$ be a genuine function and $u \in F$.  Then
  \begin{equation}\label{eq:unipotent-lines}
    \begin{aligned}
      R_X (n(u) \Phi)((x', y'), \tau) &= R_X \Phi((x', y' + u x'), \tau), \\
      R_Y (n^-(u) \Phi)((x', y'), \tau) &= R_Y \Phi((x' + u y', y'), \tau),
    \end{aligned}
  \end{equation}
  whenever the integrals on either side converge absolutely.  If $\Phi$ satisfies \eqref{eq:growth}, then $(\mathcal{F} n(u) \Phi)_X(x', y') = (\mathcal{F} \Phi)_X(x', y' + u x')$ and $(\mathcal{F} n^-(u) \Phi)_Y(x', y') = (\mathcal{F} \Phi)_Y(x' + u y', y')$.
\end{lemma}
\begin{proof}
  The shear $(x, y) \mapsto (x, y + u x)$, applied to both $v$ and $v'$, preserves the pairing $x y' - y x'$.  It fixes $x$ and $x'$, hence the measure $d x / \lvert x' \rvert$ and the factor $\{x, -x'\}_n$ in Definition \ref{definition:lines}.  The integrands on the two sides of the first identity thus coincide.  The second identity follows in the same way from the shear $(x, y) \mapsto (x + u y, y)$.  The shears preserve the bound \eqref{eq:growth}, so integrating against $\psi(\tau)$ gives the identities for $\mathcal{F}$.
\end{proof}

Let $\Phi \in \mathcal{D}$.  Then $n(u) \Phi \in \mathcal{D}$, and $\mathcal{F} \Phi$ is a genuine function.  By Lemma \ref{lemma:unipotent-properties}, the genuine functions $\mathcal{F} n(u) \Phi$ and $n(u) \mathcal{F} \Phi$ have the same $X$-component, so they are equal.  Likewise, $\mathcal{F} n^-(u) \Phi = n^-(u) \mathcal{F} \Phi$.  In particular, $n(u)$ and $n^-(u)$ preserve $\mathcal{F} \mathcal{D}$, hence $\mathcal{S} = \mathcal{D} + \mathcal{F} \mathcal{D}$.  They commute with the dilations by $F^{\times n}$, and $\lvert n(u) \Phi \rvert(v) = \lvert \Phi \rvert(v \, \pi(n(u)))$ for each genuine function $\Phi$, and likewise for $n^-(u)$.  Passing from the generators to the group, we obtain the following:
\begin{corollary}\label{corollary:metaplectic}
  Each $g \in \widetilde{\mathrm{SL}}_2(F)$ preserves $\mathcal{D}$ and $\mathcal{S}$, commutes with the dilations by $F^{\times n}$, and satisfies $\mathcal{F} g \Phi = g \mathcal{F} \Phi$ for $\Phi \in \mathcal{D}$.  For each genuine function $\Phi$, we have $\lvert g \Phi \rvert(v) = \lvert \Phi \rvert(v \, \pi(g))$.
\end{corollary}

\subsection{Torus elements}\label{sec:torus}
Let $a \in F^\times$.  The diagonal matrix with entries $a, a^{-1}$ sends $(x, y)$ to $(a x, y/a)$.  This substitution does not preserve the symbol $\{x, y\}_n$, so it does not by itself act on genuine functions.  Following Deligne \cite[\S1.9]{SB_1978-1979__21__244_0}, we lift the matrix by writing it as a product of unipotent matrices.  Set
\begin{equation*}
  w(a) := n(a) n^-(-1/a) n(a),
  \qquad
  h(a) := w(a) w(1)^{-1}.
\end{equation*}
Then $h(a)$ projects to the diagonal matrix with entries $a, a^{-1}$.

We record the action of $h(a)$ in the two charts, and how it commutes with the unipotent lifts.

\begin{lemma}\label{lemma:torus}
  Let $a \in F^\times$.
  \begin{enumerate}[(i)]
  \item\label{item:torus-formula} For each genuine function $\Psi$, we have
    \begin{align*}
      (h(a) \Psi)_X(x, y) &= \{x, a\}_n \{-1, a\}_n \Psi_X(a x, y / a), \\
      (h(a) \Psi)_Y(x, y) &= \{a, y\}_n \Psi_Y(a x, y / a).
    \end{align*}
  \item\label{item:torus-unipotent} For $u \in F$, we have $n(u) h(a) = h(a) n(u / a^2)$ and $n^-(u) h(a) = h(a) n^-(a^2 u)$.
  \end{enumerate}
\end{lemma}
\begin{proof}
  \eqref{item:torus-formula} Let $d_1(a)$ and $d_2(a)$ denote Deligne's lifts of the diagonal matrices with entries $a, 1$ and $1, a$.  They preserve the trivializations over $U_X$ and $U_Y$, respectively \cite[\S1.12, (1.14.3)]{SB_1978-1979__21__244_0}:
  \begin{equation*}
    (d_1(a) \Psi)_X(x, y) = \Psi_X(a x, y),
    \qquad
    (d_2(a) \Psi)_Y(x, y) = \Psi_Y(x, a y).
  \end{equation*}
  Moreover, $h(a) = d_1(a) d_2(a^{-1})$ \cite[(1.9.3)]{SB_1978-1979__21__244_0}.  By the transition relation \eqref{eq:transition-relation}, the multiplier of $d_1(a)$ in the other chart is
  \begin{equation*}
    \frac{\{a x, y\}_n}{\{x, y\}_n} = \{a, y\}_n,
  \end{equation*}
  that is, $(d_1(a) \Psi)_Y(x, y) = \{a, y\}_n \Psi_Y(a x, y)$.  Applying this to $d_2(a^{-1}) \Psi$ gives the formula for the $Y$-component of $h(a) \Psi$.  The formula for the $X$-component follows, first on $U_X \cap U_Y$ and then on $U_X$ by continuity, from the transition relation and the identity
  \begin{equation*}
    \{a x, y / a\}_n = \{x, y\}_n \{a, y\}_n \{x, a\}_n^{-1} \{-1, a\}_n,
  \end{equation*}
  together with $\{-1, a\}_n^2 = 1$.

  \eqref{item:torus-unipotent} By \eqref{item:torus-formula}, the multiplier of $h(a)$ in the $X$-chart depends only upon $x$, which the shear $(x, y) \mapsto (x, y + u x)$ fixes.  Both sides of the first identity thus send $\Psi$ to the genuine function with $X$-component $\{x, a\}_n \{-1, a\}_n \Psi_X(a x, (y + u x) / a)$.  The second identity follows in the same way in the $Y$-chart, where the multiplier depends only upon $y$.
\end{proof}

\section{The basic vector}\label{sec:basic-vector}
In this section, we assume that $F$ is non-archimedean, of residue characteristic not dividing $n$.

For $n = 1$ and $\psi$ unramified, the characteristic function of $\mathfrak{o}^2$ is fixed by $\mathrm{SL}_2(\mathfrak{o})$ and by the Fourier transform.  For general $n$, its role is played by the basic vector
\begin{equation*}
  \Phi^0 = \sum_{k \geq 0} q^{(n-1) k} \, \Phi^1(\varpi^{- n k} \, \cdot)
\end{equation*}
of \S\ref{sec:intertwining}.  Explicitly, write $(x, y) = \varpi^k (x_0, y_0)$ with $(x_0, y_0) \in \mathfrak{o}^2$ primitive.  If $k$ is a nonnegative multiple of $n$, then
\begin{equation*}
  \begin{aligned}
    \Phi^0_X(x, y) &= q^{(n-1) k / n} \cdot
    \begin{cases}
      1 & \text{ if } x_0 \in \mathfrak{o}^\times, \\
      \{x, y\}_n & \text{ otherwise,}
    \end{cases}
    \\
    \Phi^0_Y(x, y) &= q^{(n-1) k / n} \cdot
    \begin{cases}
      1 & \text{ if } y_0 \in \mathfrak{o}^\times, \\
      \{y, x\}_n & \text{ otherwise.}
    \end{cases}
  \end{aligned}
\end{equation*}
Otherwise, $\Phi^0_X(x, y) = 0$ and $\Phi^0_Y(x, y) = 0$.  The function $\Phi^0$ does not depend upon the choice of $\varpi$.  The values at the primitive vectors are chosen so that $\Phi^0$ is invariant under the unipotent translations by elements of $\mathfrak{o}$ (Lemma \ref{lemma:basic-vector-invariant}).  For $n = 1$, we obtain the characteristic function of $\mathfrak{o}^2$, restricted to $U$.

The bound $\lvert \Phi^0(v) \rvert \leq \lVert v \rVert^{-1 + 1/n} 1_{\lVert v \rVert \leq 1}$ ensures that $\mathcal{F} \Phi^0$ is defined.  For every $n$, the support of $\Phi^0$ accumulates at the origin, so $\Phi^0 \notin \mathcal{D}$.  The function $\Phi^1$ of \S\ref{sec:intertwining} agrees with $\Phi^0$ at the primitive vectors of $\mathfrak{o}^2$ and vanishes elsewhere.  It lies in $\mathcal{D}$, and
\begin{equation}\label{eq:primitive-layer}
  \Phi^1 = \Phi^0 - q^{n-1} \Phi^0(\varpi^{-n} \, \cdot).
\end{equation}

\begin{lemma}\label{lemma:basic-vector-invariant}
  For $u \in \mathfrak{o}$, we have $n(u) \Phi^0 = \Phi^0$ and $n^-(u) \Phi^0 = \Phi^0$.  Consequently, $h(a) \Phi^0 = \Phi^0$ for $a \in \mathfrak{o}^\times$.  The same holds for $\Phi^1$.
\end{lemma}
\begin{proof}
  We treat $n(u)$, starting with $\Phi^1$.  The shear $(x, y) \mapsto (x, y + u x)$ preserves the set of primitive vectors of $\mathfrak{o}^2$.  By \eqref{eq:unipotent-charts}, it thus suffices to check that $\Phi^1_X(x, y + u x) = \Phi^1_X(x, y)$ when $(x, y)$ is primitive and $x \neq 0$.  If $x$ is a unit, then both sides are $1$.  Otherwise, $y$ is a unit, and the two sides differ by the factor $\{x, 1 + u x / y\}_n$, which is trivial by property \eqref{item:hilbert-tame} of the Hilbert symbol, since $1 + u x / y \in 1 + \varpi \mathfrak{o}$.  Thus $n(u) \Phi^1 = \Phi^1$.  Since $n(u)$ commutes with the dilations by $F^{\times n}$, it fixes each term of the sum defining $\Phi^0$, hence $\Phi^0$.  The assertions concerning $h(a)$ follow because, for $a \in \mathfrak{o}^\times$, the operator $h(a) = w(a) w(1)^{-1}$ is a product of operators $n(u)$ and $n^-(u)$ with $u \in \mathfrak{o}$.
\end{proof}

We return to the spherical section $e_s$ of \S\ref{sec:intertwining}.

\begin{lemma}\label{lemma:spherical}
  Suppose that $\psi$ is unramified, and let $\Re(s) > 0$.  Then $A e_s = \frac{L(s)}{L(1 + s)} \, e_{-s}$.
\end{lemma}
\begin{proof}
  By Lemma \ref{lemma:basic-vector-invariant}, the shears $n(u)$ and $n^-(u)$ for $u \in \mathfrak{o}$ fix $e_s$ and $e_{-s}$.  They commute with $A$ by \eqref{eq:unipotent-lines} with $\tau = 1$.  The shears carry every vector to the line $x = 0$, on which the elements $h(a)$ for $a \in \mathfrak{o}^\times$ act by unit dilations with the factor in Lemma \ref{lemma:torus}.  Under $v \mapsto \mu^n v$, both sides scale by $\lvert \mu \rvert^{-n+s}$, as one sees for $A e_s$ by scaling the integration variable by $\mu^{-n}$.  It therefore suffices to check the identity at $(0, \varpi^j)$ for $0 \leq j < n$.

  The case $j = 0$ is Lemma \ref{lemma:spherical-axis}.  For $0 < j < n$, the function $e_{-s}$ vanishes at $(0, \varpi^j)$.  By Definition \ref{definition:lines}, the value of $(A e_s)_Y$ there is proportional to the integral of $(e_s)_X(\varpi^{-j}, y)$ over $y$.  Only the shells $y = \varpi^{-nk} u$, with $nk > j$ and $u \in \mathfrak{o}^\times$, meet the support of this integrand.  On each such shell its dependence on $u$ is through $\{\varpi, u\}_n^{-j}$, a nontrivial character by property \eqref{item:hilbert-tame} of the Hilbert symbol.  Its integral is therefore zero.
\end{proof}

\begin{proposition}\label{proposition:basic-vector-fixed}
  \begin{enumerate}[(i)]
  \item\label{item:basic-fixed}\lean{MetaplecticPoisson.Datum.basicCharts_fourier_fixed} If $\psi$ is unramified, then $\mathcal{F} \Phi^0 = \Phi^0$.
  \item\label{item:basic-schwartz}\lean{MetaplecticPoisson.Datum.basic_mem} The basic vector $\Phi^0$ lies in $\mathcal{S}$.
  \end{enumerate}
\end{proposition}
\begin{proof}
  \eqref{item:basic-fixed} One can deduce this identity by Mellin inversion from \eqref{eq:cuue7ixq1x}, which gives $\mathcal{F}(\Phi^0[s]) = \Phi^0[-s]$.  We give a direct proof in \S\ref{sec:basic-vector-transform} by computing $\mathcal{F} \Phi^0$ on the line $x = 0$ and using the invariance of Lemma \ref{lemma:basic-vector-invariant}.

  \eqref{item:basic-schwartz} By Lemma \ref{lemma:schwartz-dilation}, we may suppose that $\psi$ is unramified.  By \eqref{eq:primitive-layer}, Lemma \ref{lemma:transform-dilation} and \eqref{item:basic-fixed},
  \begin{equation*}
    \mathcal{F} \Phi^1
    = \mathcal{F} \Phi^0 - q^{n-1} \mathcal{F} \bigl( \Phi^0(\varpi^{-n} \, \cdot) \bigr)
    = \Phi^0 - q^{-n-1} \, \Phi^0(\varpi^n \, \cdot).
  \end{equation*}
  Dilating \eqref{eq:primitive-layer}, we have $\Phi^0(\varpi^n \, \cdot) = q^{n-1} \Phi^0 + \Phi^1(\varpi^n \, \cdot)$.  Hence
  \begin{equation*}
    (1 - q^{-2}) \, \Phi^0 = \mathcal{F} \Phi^1 + q^{-n-1} \, \Phi^1(\varpi^n \, \cdot),
  \end{equation*}
  which lies in $\mathcal{S}$.
\end{proof}

In the global setting, the following situation occurs at almost all places:
\begin{definition}\label{definition:good}
  We say that the triple $(F, \psi, n)$ is \emph{good} if $F$ is non-archimedean, of residue characteristic not dividing $n$, and $\psi$ is trivial on $\mathfrak{o}$.
\end{definition}

\begin{corollary}\label{corollary:basic-vector-general}
  Suppose that $(F, \psi, n)$ is good.  Let $b \in \mathfrak{o} - \{0\}$ be such that $b^{-1} \mathfrak{o}$ is the largest fractional ideal on which $\psi$ is trivial.  Then $\mathcal{F} \Phi^0 = \lvert b \rvert^n \Phi^0(b^n \, \cdot)$.
\end{corollary}
\begin{proof}
  The character $\psi(b^{-1} \, \cdot)$ is unramified.  We apply Proposition \ref{proposition:basic-vector-fixed} to it, and then Lemma \ref{lemma:transform-dilation}.
\end{proof}

\section{Fourier inversion}\label{sec:inversion}
For $n = 1$, the symplectic Fourier transform is an involution commuting with $\mathrm{SL}_2(F)$.  The analogous properties of $\mathcal{F}$ are as follows:

\begin{theorem}\label{theorem:schwartz-space}\lean{MetaplecticPoisson.Datum.schwartz_inversion}
  Set $\varepsilon := (-1)^{n+1}$.
  \begin{enumerate}[(i)]
  \item\label{item:schwartz-inversion} The transform $\mathcal{F}$ maps $\mathcal{S}$ bijectively onto itself, and $\mathcal{F}^2 = h(\varepsilon)$ on $\mathcal{S}$.
  \item\label{item:schwartz-stable} Each $g \in \widetilde{\mathrm{SL}}_2(F)$ preserves $\mathcal{S}$, and $\mathcal{F} g \Phi = g \mathcal{F} \Phi$ for $\Phi \in \mathcal{S}$.
  \item\label{item:schwartz-plancherel} For $\Phi \in \mathcal{S}$, we have $\lVert \mathcal{F} \Phi \rVert_2 = \lVert \Phi \rVert_2$, where $\lVert \Phi \rVert_2^2 := \int_U \lvert \Phi \rvert(v)^2 \, d v$.
  \end{enumerate}
\end{theorem}
\begin{proof}
  We prove in Appendix \ref{sec:marginals} that $\mathcal{F}^2 \Phi = h(\varepsilon) \Phi$ for $\Phi \in \mathcal{D}$ (Proposition \ref{proposition:inversion}).  The operator $h(\varepsilon)$ preserves $\mathcal{D}$ and commutes with $\mathcal{F}$ on $\mathcal{D}$ (Corollary \ref{corollary:metaplectic}).

  \eqref{item:schwartz-inversion} Since $\mathcal{F}^2 \mathcal{D} = \mathcal{D}$, we have $\mathcal{F}\mathcal{S} = \mathcal{S}$.  For $\Phi \in \mathcal{D}$,
  \begin{equation*}
    \mathcal{F}^2 (\mathcal{F} \Phi) = \mathcal{F} (h(\varepsilon) \Phi) = h(\varepsilon) \mathcal{F} \Phi.
  \end{equation*}
  The inversion formula therefore holds on $\mathcal{S}$ and implies injectivity there.

  \eqref{item:schwartz-stable} Each $g$ preserves $\mathcal{S}$ (Corollary \ref{corollary:metaplectic}).  Let $\Phi \in \mathcal{S}$ and $u \in F$.  By \eqref{item:schwartz-inversion}, both $\mathcal{F} n(u) \Phi$ and $n(u) \mathcal{F} \Phi$ are genuine functions.  They have the same $X$-component, by Lemma \ref{lemma:unipotent-properties}, so they are equal.  Likewise, $\mathcal{F} n^-(u) \Phi = n^-(u) \mathcal{F} \Phi$.  Thus $\mathcal{F}$ commutes with the generators of $\widetilde{\mathrm{SL}}_2(F)$.

  \eqref{item:schwartz-plancherel} The elements of $\mathcal{S}$ are square-integrable, by \eqref{eq:growth}.  For $\Phi, \Psi \in \mathcal{S}$, write $\langle \Phi, \Psi \rangle := \int_U \Phi \overline{\Psi}$, the integrand being $\Phi_X \overline{\Psi_X}$ on $U_X$.  Let $K$ be as in Lemma \ref{lemma:root-kernel}.  We claim that, for $v, w \in U_X$ with $[v, w] \neq 0$,
  \begin{equation}\label{eq:kernel-adjoint}
    \overline{K(v, w)} = \{\varepsilon, x'\}_n \, K(\varepsilon w, v).
  \end{equation}
  Indeed, exchanging $v$ and $w$ replaces $[v, w]$ by its negative, and conjugating $\psi(\tau)$ replaces $\tau$ by $- \tau$, so the sum over the roots of $[w, v]$ becomes the sum over the roots of $(-1)^{n+1} [v, w] = [v, \varepsilon w]$.  The symbols on the two sides differ by $\{-x, x'\}_n \{x, - \varepsilon x'\}_n^{-1} = \{-1, x'\}_n \{- \varepsilon, x\}_n$, which equals $\{\varepsilon, x'\}_n$ both for odd and for even $n$.  By Lemma \ref{lemma:torus}, we have $(h(\varepsilon)^{-1} \Psi)_X(v) = \{\varepsilon, x\}_n \Psi_X(\varepsilon v)$.  By \eqref{eq:root-kernel} and \eqref{eq:kernel-adjoint}, it follows that
  \begin{equation*}
    \langle \mathcal{F} \Phi, \Psi \rangle = \langle \Phi, h(\varepsilon)^{-1} \mathcal{F} \Psi \rangle.
  \end{equation*}
  The interchange of integrals is justified because $K(w, v) \ll \lvert [v, w] \rvert^{-1 + 1/n}$, which is locally integrable transverse to the locus $[v, w] = 0$, while the bound \eqref{eq:growth} for $\Phi$ and $\Psi$ gives the convergence at the origin and at infinity.  Taking $\Psi = \mathcal{F} \Phi$ and applying \eqref{item:schwartz-inversion}, we obtain $\lVert \mathcal{F} \Phi \rVert_2^2 = \langle \Phi, h(\varepsilon)^{-1} \mathcal{F}^2 \Phi \rangle = \lVert \Phi \rVert_2^2$.
\end{proof}
Since $\mathcal{D}$ is dense in the space of square-integrable genuine functions and $\mathcal{F} \mathcal{S} = \mathcal{S}$, the transform $\mathcal{F}$ extends to a unitary operator on that space.

For odd $n$, the operator $h(\varepsilon) = h(1)$ is the identity; in particular, for $n = 1$, we recover the involutivity of the symplectic Fourier transform.  For even $n$, Lemma \ref{lemma:torus} gives $(h(-1) \Phi)_Y(x, y) = \{-1, y\}_n \Phi_Y(-x, -y)$.  The element $\varepsilon = - (-1)^n$ reflects two signs: Fourier inversion on $F$, applied in the variable $\tau$, replaces $\tau^n$ by $(-1)^n \tau^n$, and the antisymmetry of $x y' - y x'$ contributes a further sign.

\begin{remark}
  On functions satisfying \eqref{eq:principal-series-homogeneity}, $\mathcal{F}$ is a multiple of the standard intertwining operator (\S\ref{sec:intertwining}), and the inversion formula corresponds to the fact that the composition of two such operators in opposite directions is a multiple of $h(\varepsilon)$.  The condition~\eqref{eq:principal-series-homogeneity} prescribes the action of the dilations by $F^{\times n}$ only, and $h(\varepsilon)$ need not act on such functions by a scalar: for $F = \mathbb{R}$ and $n = 2$, the operator $h(-1)$ exchanges opposite rays.  The element $h(\varepsilon)$ is central in $\widetilde{\mathrm{SL}}_2(F)$ (by Lemma \ref{lemma:torus}) and acts on a principal series representation by the central character of the inducing representation.  For $p$-adic $F$, Goldberg and Szpruch \cite{MR3544420} (resp.\ Szpruch~\cite{Szpruch2019}) compute the corresponding Plancherel measures when $p \nmid n$ (resp.\ when $4 \nmid n$).  One would expect to recover the inversion formula from such results by Mellin inversion with respect to the dilations by $F^{\times n}$.  The proof recorded in Appendix \ref{sec:marginals} is direct.
\end{remark}

\section{Leading terms at the origin}\label{sec:germs}
For $n = 1$, the Poisson summation formula involves the values of $\Phi$ and $\mathcal{F} \Phi$ at the origin.  For $n \geq 2$, a general element $\Psi$ of $\mathcal{S}$ has no natural value at the origin.  We consider instead, for each $v \in U$, the limiting value at the origin of the function
\begin{equation}\label{eq:normalized-ray}
  \mu \mapsto \lvert \mu \rvert^{n-1} \Psi(\mu^n v)
  \qquad (\mu \in F^\times),
\end{equation}
formed in either chart.

\begin{proposition}\label{proposition:schwartz-origin}\lean{MetaplecticPoisson.Datum.exists_leading}
  Let $\Psi \in \mathcal{S}$.  In either chart, the function \eqref{eq:normalized-ray} extends to a Schwartz--Bruhat function on $F$.  As $v$ varies over a compact subset of the chart, these extensions form a bounded subset of the space of Schwartz--Bruhat functions.  At an archimedean place, the same holds for their derivatives with respect to $v$.\lean{MetaplecticPoisson.Datum.schwartz_rayBound_real}\lean{MetaplecticPoisson.Datum.schwartz_rayBound_complex}
\end{proposition}
\begin{proof}
  For $\Psi \in \mathcal{D}$, the extension by zero has the required properties by the definition of $\mathcal{D}$.  For $\Psi = \mathcal{F} \Phi$ with $\Phi \in \mathcal{D}$, Lemma \ref{lemma:transform-fourier} identifies \eqref{eq:normalized-ray} in the $X$-chart with $\lvert n \rvert^{-1/2}$ times the Fourier transform of $R_X \Phi(v, \cdot)$, and likewise in the $Y$-chart.  The assertions follow from Lemma \ref{lemma:lines}\eqref{item:lines-smooth} and $\mathcal{S} = \mathcal{D} + \mathcal{F} \mathcal{D}$.
\end{proof}

Evaluating the extension at $\mu = 0$ in each chart gives a pair of functions, which we call the \emph{leading term} $b \Psi$ of $\Psi$ at the origin.  For $\Phi \in \mathcal{D}$, we have $b \Phi = 0$ by the defining conditions on $\mathcal{D}$, while letting \(\mu\to0\) in Lemma~\ref{lemma:transform-fourier} gives
\begin{equation}\label{eq:leading-term-transform}
  (b \mathcal{F} \Phi)_X(v) = \lvert n \rvert^{-1/2} \int_F R_X \Phi(v, \tau) \, d \tau,
  \qquad
  (b \mathcal{F} \Phi)_Y(v) = \lvert n \rvert^{-1/2} \int_F R_Y \Phi(v, \tau) \, d \tau.
\end{equation}
For $n = 1$, this is the constant $\int_{F^2} \Phi$, which is the value of $\mathcal{F} \Phi$ at the origin.  The leading term has the following properties:
\begin{itemize}
\item The pair $b \Psi$ is a genuine function, and $(b \Psi)(\mu^n \, \cdot) = \lvert \mu \rvert^{1-n} \, b \Psi$ for $\mu \in F^\times$.  The relation between the two components and the homogeneity follow from the definition, and the smoothness from \eqref{eq:leading-term-transform} and Lemma \ref{lemma:lines}\eqref{item:lines-smooth}.\lean{MetaplecticPoisson.Datum.leading_smooth_real}\lean{MetaplecticPoisson.Datum.leading_smooth_complex}
\item For $a \in F^{\times n}$, we have $b(\Psi(a \, \cdot)) = \lvert a \rvert^{-1 + 1/n} \, b \Psi$.
\item For $g \in \widetilde{\mathrm{SL}}_2(F)$, we have $b (g \Psi) = g (b \Psi)$.  It suffices to verify this for $g = n(u)$ and $g = n^-(u)$, where it follows from the definitions.
\item If the residue characteristic of $F$ does not divide $n$, then $b \Phi^0$ agrees with $\Phi^0$ at the vectors $v$ with $\lVert v \rVert \leq 1$.  In particular, $(b \Phi^0)_Y(0, y)$ equals $\lvert y \rvert^{-1 + 1/n}$ if $n$ divides $v(y)$, and vanishes otherwise.
\end{itemize}

\begin{proposition}\label{proposition:finite-germs}
  Suppose that $F$ is non-archimedean.  Each $\Psi \in \mathcal{S}$ vanishes outside a compact subset of $F^2$ and agrees with $b \Psi$ near the origin.  The kernel of $b$ on $\mathcal{S}$ is $\mathcal{D}$.
\end{proposition}
\begin{proof}
  Choose compact subsets of $U_X$ and $U_Y$ whose union meets every $F^{\times n}$-orbit in $U$.  As $v$ ranges over either compact set, the extensions in Proposition \ref{proposition:schwartz-origin} form a bounded family of Schwartz--Bruhat functions on $F$.  They thus have common compact support and are constant on a common neighborhood of zero.  The support assertion follows, and for $\lvert \mu \rvert$ sufficiently small we obtain
  \begin{equation*}
    \Psi(\mu^n v) = \lvert \mu \rvert^{1-n} b \Psi(v) = b \Psi(\mu^n v)
  \end{equation*}
  in either chart, by homogeneity.  If $b \Psi = 0$, then $\Psi$ has compact support in $U$, so $\Psi \in \mathcal{D}$.
\end{proof}

For $n = 1$, every constant occurs as $b \mathcal{F} \Phi = \int_{F^2} \Phi$ for some $\Phi \in \mathcal{D}$.  At a finite place, it follows from Proposition \ref{proposition:finite-germs} that $\mathcal{S}$ consists of the restrictions to $U$ of the elements of $C_c^\infty(F^2)$, with $b \Psi$ giving the value at the origin.  The same holds for $F$ archimedean, with the Schwartz space of $F^2$ in place of $C_c^\infty(F^2)$.  Indeed, by a theorem of Estrada \cite[Theorem~4.1]{Estrada1998VectorMoment}, each Schwartz function on $F^2$ has the same moments as some $\Phi \in \mathcal{D}$.  The Fourier transform of their difference therefore vanishes at the origin together with all its derivatives, so the difference belongs to $\mathcal{F} \mathcal{D}$.

For general $n$, the leading term of a transform is the image of a homogeneous function under the intertwining operator $A$ of \S\ref{sec:intertwining}.  For $\Phi \in \mathcal{D}$, set
\begin{equation*}
  P \Phi := \int_{F^{\times n}} \Phi(a \, \cdot) \lvert a \rvert^{1 + 1/n} \, d^\times a.
\end{equation*}
It is a genuine function, and $(P \Phi)(a \, \cdot) = \lvert a \rvert^{-1 - 1/n} P \Phi$ for $a \in F^{\times n}$.

\begin{lemma}\label{lemma:germ-via-average}
  For each $\Phi \in \mathcal{D}$, we have
  \begin{equation}\label{eq:germ-via-average}
    b \mathcal{F} \Phi = \frac{n}{\lvert n \rvert^{3/2}} A(P \Phi).
  \end{equation}
\end{lemma}
\begin{proof}
  The average $P(\lvert \Phi \rvert)$ is bounded on a compact set meeting every $F^{\times n}$-orbit, hence $\ll \lVert v \rVert^{-1 - 1/n}$ by homogeneity.  It is thus integrable on the lines defining $A$, which avoid the origin, so we may interchange the line integral with the integral over $a$.  Let $a = \mu^n$.  The substitution $x \mapsto x / a$ in Definition \ref{definition:lines} gives $R_X (\Phi(a \, \cdot))(v, 1) = \lvert a \rvert^{-1} R_X \Phi(v, \mu)$.  We multiply by $\lvert a \rvert^{1 + 1/n}$, integrate over $a \in F^{\times n}$, and apply \eqref{eq:haar-on-powers} and \eqref{eq:leading-term-transform}.
\end{proof}

\begin{corollary}\label{corollary:finite-schwartz}
  Suppose that $F$ is non-archimedean.  The map $b$ identifies $\mathcal{S}/\mathcal{D}$ with the image of $A$ on the space of genuine functions $f$ satisfying
  \begin{equation*}
    f(a v) = \lvert a \rvert^{-1 - 1/n} f(v)
    \qquad (a \in F^{\times n}).
  \end{equation*}
  A genuine function belongs to $\mathcal{S}$ if and only if it vanishes outside a compact subset of $F^2$ and agrees near the origin with $A f$ for some such $f$.
\end{corollary}
\begin{proof}
  Every such $f$ occurs as $P \Phi$.  Indeed, take $\Phi$ equal to $\vol((\mathfrak{o}^\times)^n)^{-1} f$ on $\lvert \varpi \rvert^n < \lVert v \rVert \leq 1$ and zero elsewhere.  For each $v$, the $a \in F^{\times n}$ for which $a v$ lies in this annulus form one coset of $(\mathfrak{o}^\times)^n$, so $P \Phi = f$.

  The assertion about $\mathcal{S}/\mathcal{D}$ follows from Lemma \ref{lemma:germ-via-average} and Proposition \ref{proposition:finite-germs}.  By that proposition, each element of $\mathcal{S}$ has the stated support and germ.  Conversely, given a function with those properties, subtracting an element of $\mathcal{S}$ with the same leading term leaves an element of $\mathcal{D}$.
\end{proof}

\begin{definition}\label{definition:genuine-vanishing}
  We say that $\Phi \in \mathcal{D}$ satisfies the \emph{vanishing condition} if $P \Phi = 0$.
\end{definition}
If $\Phi$ satisfies the vanishing condition, then $b \mathcal{F} \Phi = 0$, by \eqref{eq:germ-via-average}.

\section{The Poisson formula}\label{sec:poisson}
We now work globally.

\subsection{Setting}\label{sec:main-formula}
Let $F$ be a global field containing a primitive $n$th root of unity; thus $F$ is a number field, or a function field of characteristic not dividing $n$.  We denote by $\mathbb{A}$ the adele ring of $F$.  Let $\psi : \mathbb{A} / F \rightarrow \mathbb{C}^{(1)}$ be a nontrivial unitary character.  Let $S$ be a nonempty finite set of places, large enough that for each $\mathfrak{p} \notin S$, the triple $(F_{\mathfrak{p}}, \psi_{\mathfrak{p}}, n)$ is good (Definition~\ref{definition:good}).  We write $F_S := \prod_{\mathfrak{p} \in S} F_{\mathfrak{p}}$, and $\mathcal{O}_S \subset F$ for the ring of $S$-integers.  We apply the local definitions at each place $\mathfrak{p}$, using the character $\psi_{\mathfrak{p}}$, and indicate the place by a subscript when necessary.

We recall Hilbert's reciprocity law: for $x, y \in F^\times$, we have $\{x, y\}_{\mathfrak{p}, n} = 1$ for almost all $\mathfrak{p}$, and $\prod_{\mathfrak{p}} \{x, y\}_{\mathfrak{p}, n} = 1$.

\subsection{The global Schwartz space}\label{sec:rational-points}
For each place $\mathfrak{p}$, we write $\mathcal{S}_{\mathfrak{p}}$ for the Schwartz space over $F_{\mathfrak{p}}$ (Definition \ref{definition:schwartz-space}).  For each finite place $\mathfrak{p}$ not dividing $n$, we write $\Phi_{\mathfrak{p}}^0 \in \mathcal{S}_{\mathfrak{p}}$ for the basic vector (Proposition \ref{proposition:basic-vector-fixed}).  We define
\begin{equation*}
  \mathcal{S}_{\mathbb{A}} := \sideset{}{'}\bigotimes_{\mathfrak{p}} \mathcal{S}_{\mathfrak{p}},
\end{equation*}
the algebraic restricted tensor product with respect to the basic vectors $\Phi_{\mathfrak{p}}^0$.  For a pure tensor $\Phi = \otimes_{\mathfrak{p}} \Phi_{\mathfrak{p}}$ in $\mathcal{S}_{\mathbb{A}}$, we set
\begin{align*}
  \Phi_X(x, y) &:= \prod_{\mathfrak{p}} \Phi_{\mathfrak{p}, X}(x_{\mathfrak{p}}, y_{\mathfrak{p}})
  && \text{for } (x, y) \in \mathbb{A}^\times \times \mathbb{A},
  \\
  \Phi_Y(x, y) &:= \prod_{\mathfrak{p}} \Phi_{\mathfrak{p}, Y}(x_{\mathfrak{p}}, y_{\mathfrak{p}})
  && \text{for } (x, y) \in \mathbb{A} \times \mathbb{A}^\times.
\end{align*}
In the first product, for almost every $\mathfrak{p}$, the factor of $\Phi$ is the basic vector, $x_{\mathfrak{p}}$ is a unit and $y_{\mathfrak{p}}$ is integral, so the factor of the product is one; likewise for the second.  At a rational point $(x, y)$ with $x y \neq 0$, the two values agree by the reciprocity law, so together they define the value $\Phi(v)$ at each point $v \in F^2 - \{0\}$.  These values are linear in each factor of $\Phi$, and we extend them linearly to all $\Phi \in \mathcal{S}_{\mathbb{A}}$.

Since $\psi_{\mathfrak{p}}$ is unramified for almost every $\mathfrak{p}$, we have $\mathcal{F} \Phi_{\mathfrak{p}}^0 = \Phi_{\mathfrak{p}}^0$ for almost every $\mathfrak{p}$ (Proposition \ref{proposition:basic-vector-fixed}).  By Theorem \ref{theorem:schwartz-space}, there is thus a linear map $\mathcal{F} : \mathcal{S}_{\mathbb{A}} \rightarrow \mathcal{S}_{\mathbb{A}}$ with $\mathcal{F} \Phi = \otimes_{\mathfrak{p}} \mathcal{F} \Phi_{\mathfrak{p}}$ for each pure tensor $\Phi$.  The local normalizing factors $\lvert n \rvert_{\mathfrak{p}}^{-1/2}$ equal one for $\mathfrak{p} \notin S$, and their product over all places is one, by the product formula.  The values of $\mathcal{F} \Phi$ at the points of $F^2 - \{0\}$ are thus the same as for the transform defined without these factors.

\subsection{Statement}\label{sec:poisson-statement}
For each place $\mathfrak{p}$, let $b(\mathcal{S}_{\mathfrak{p}})$ denote the space of local leading terms defined in \S\ref{sec:germs}.  The leading term $b \Phi_{\mathfrak{p}}^0$ of the basic vector is nonzero (\S\ref{sec:germs}).  We write
\begin{equation*}
  b(\mathcal{S}_{\mathbb{A}}) := \sideset{}{'}\bigotimes_{\mathfrak{p}} b(\mathcal{S}_{\mathfrak{p}})
\end{equation*}
for the restricted tensor product with respect to the vectors $b \Phi_{\mathfrak{p}}^0$, and
\begin{equation*}
  b : \mathcal{S}_{\mathbb{A}} \rightarrow b(\mathcal{S}_{\mathbb{A}})
\end{equation*}
for the surjective linear map given on pure tensors by $b \Phi = \otimes_{\mathfrak{p}} b \Phi_{\mathfrak{p}}$.

\begin{theorem}[Poisson formula]\label{theorem:poisson}\lean{MetaplecticPoisson.Datum.poisson}
  There is a unique linear functional $B$ on $b(\mathcal{S}_{\mathbb{A}})$ such that for every $\Phi \in \mathcal{S}_{\mathbb{A}}$,
  \begin{equation*}
    \sum_{v \in F^2 - \{0\}} \Phi(v) + B(b \Phi)
    =
    \sum_{v \in F^2 - \{0\}} \mathcal{F} \Phi(v) + B(b \mathcal{F} \Phi),
  \end{equation*}
  with both sums converging absolutely.
\end{theorem}
For $n = 1$, we have $B(b \Phi) = \prod_{\mathfrak{p}} \Phi_{\mathfrak{p}}(0)$ for each pure tensor $\Phi$, by \eqref{eq:intro-poisson-A} and the uniqueness.

\begin{corollary}\label{corollary:poisson-vanishing}\lean{MetaplecticPoisson.Datum.poisson_of_leading_eq_zero}
  Let $\Phi \in \mathcal{S}_{\mathbb{A}}$ with $b \Phi = 0$ and $b \mathcal{F} \Phi = 0$.  Then
  \begin{equation*}
    \sum_{v \in F^2 - \{0\}} \Phi(v)
    =
    \sum_{v \in F^2 - \{0\}} \mathcal{F} \Phi(v).
  \end{equation*}
  This holds in particular if $\Phi$ is a pure tensor such that $b \Phi_{\mathfrak{p}} = 0$ for some place $\mathfrak{p}$ and $b \mathcal{F} \Phi_{\mathfrak{q}} = 0$ for some place $\mathfrak{q}$\lean{MetaplecticPoisson.Datum.poisson_pure_of_local_vanishing}; for instance, if some factor $\Phi_{\mathfrak{p}}$ lies in $\mathcal{D}$ and satisfies the vanishing condition.
\end{corollary}
\begin{proof}
  Apply Theorem \ref{theorem:poisson}, observing that $b \Phi_{\mathfrak{p}} = b \mathcal{F} \Phi_{\mathfrak{p}} = 0$ for a test function $\Phi_{\mathfrak{p}} \in \mathcal{D}$ satisfying the vanishing condition (\S\ref{sec:germs}).
\end{proof}

We prove Theorem \ref{theorem:poisson} in Appendix \ref{sec:proof-poisson}.  The proof uses the Poisson summation formulas on $\mathbb{A}$ and on $\mathbb{A}^\times / \mathbb{A}^{\times n}$, the local relations of Appendix \ref{sec:marginals}, and reduction theory.

When $\mathcal{O}_S$ is a principal ideal domain, the sums in Theorem \ref{theorem:poisson} take the following classical form, in terms of the power residue symbol.
\begin{definition}\label{definition:residue-symbol}\lean{NFoldCover.residue}
  Let $m \geq 1$ with $\mu_m \subset F$.  For $c \in \mathcal{O}_S - \{0\}$ and $d \in F$, we define the power residue symbol
  \begin{equation*}
    \qr{d}{c}_m :=
    \begin{cases}
      \prod_{\mathfrak{p} \mid c} \{c, d\}_{\mathfrak{p}, m}
      & \text{ if $d$ is a unit at each $\mathfrak{p}$ dividing $c$,} \\
      0 & \text{ otherwise,}
    \end{cases}
  \end{equation*}
  where $\mathfrak{p}$ runs over the primes not in $S$, and $\{\cdot, \cdot\}_{\mathfrak{p}, m}$ denotes the Hilbert symbol with values in $\mu_m$, regarded as complex-valued as in \S\ref{sec:hilbert-symbol}.
\end{definition}
\begin{example}\label{example:S-arithmetic}
  Suppose that $\mathcal{O}_S$ is a principal ideal domain.  Let $\Phi$ be a pure tensor with $\Phi_{\mathfrak{p}} = \Phi_{\mathfrak{p}}^0$ for each $\mathfrak{p} \notin S$.  We write $\lvert \tau \rvert_S := \prod_{\mathfrak{p} \in S} \lvert \tau \rvert_{\mathfrak{p}}$ for $\tau \in F^\times$.  Let $v \in F^2 - \{0\}$ with $\Phi(v) \neq 0$.  For each $\mathfrak{p} \notin S$, the vector $v$ is then $\varpi_{\mathfrak{p}}^{k_{\mathfrak{p}}}$ times a primitive vector, with $k_{\mathfrak{p}} \in n \mathbb{Z}_{\geq 0}$.  The ideal of $\mathcal{O}_S$ generated by the coordinates of $v$ is thus the $n$th power of a principal ideal $\tau \mathcal{O}_S$, and $v = \tau^n (c, d)$ with $c, d \in \mathcal{O}_S$ coprime.  We have $q_{\mathfrak{p}}^{(n-1) k_{\mathfrak{p}} / n} = \lvert \tau \rvert_{\mathfrak{p}}^{1-n}$ for $\mathfrak{p} \notin S$.  The definition of $\Phi^0$, together with the identities $\{\tau^n c, \tau^n d\}_{\mathfrak{p}, n} = \{c, d\}_{\mathfrak{p}, n}$ and $\prod_{\mathfrak{p} \notin S} \lvert \tau \rvert_{\mathfrak{p}}^{1-n} = \lvert \tau \rvert_S^{n-1}$, thus gives
  \begin{equation*}
    \Phi(v) =
    \begin{cases}
      \lvert \tau \rvert_S^{n-1} \qr{d}{c}_n \prod_{\mathfrak{p} \in S} \Phi_{\mathfrak{p}, X}(v) & \text{ if } c \neq 0, \\
      \lvert \tau \rvert_S^{n-1} \prod_{\mathfrak{p} \in S} \Phi_{\mathfrak{p}, Y}(v) & \text{ if } c = 0.
    \end{cases}
  \end{equation*}
  The decomposition $v = \tau^n (c, d)$ is unique once $\tau$ is chosen from a fixed set of representatives for the nonzero elements of $\mathcal{O}_S$ modulo units.  This gives \eqref{eq:intro-theta-S}.  If $\psi_{\mathfrak{p}}$ is unramified for each $\mathfrak{p} \notin S$, then $\mathcal{F} \Phi$ is of the same form, by Proposition \ref{proposition:basic-vector-fixed}.
\end{example}

\section{Fourier coefficients and sums of Gauss sums}\label{sec:proof-main}
We now deduce from Theorem \ref{theorem:poisson} a summation formula on $F^\times$, by taking a Fourier coefficient with respect to the unipotent translations.  Let $\Psi$ be a pure tensor.  For $u \in \mathbb{A}$, the tensor $n(u) \Psi := \otimes_{\mathfrak{p}} n(u_{\mathfrak{p}}) \Psi_{\mathfrak{p}}$ is again a pure tensor in $\mathcal{S}_{\mathbb{A}}$, by Corollary \ref{corollary:metaplectic} and Lemma \ref{lemma:basic-vector-invariant}.  By \eqref{eq:unipotent-charts} and \eqref{eq:unipotent-axis}, it has the same values as $\Psi$ at the points $(0, y)$, while its value at a point $(x, y)$ with $x \neq 0$ is $\Psi_X(x, y + u x)$.  We multiply the sum of its values over $F^2 - \{0\}$ by $\psi(-u)$, and integrate over $u \in \mathbb{A} / F$.  The terms with $x = 0$ do not depend upon $u$, so they drop out.  The terms with $x = c \neq 0$ unfold:
\begin{equation}\label{eq:fourier-coefficient}
  \begin{aligned}
    \int_{\mathbb{A} / F} \sum_{v \in F^2 - \{0\}} (n(u) \Psi)(v) \, \psi(-u) \, d u
    &= \sum_{c \in F^\times} \int_{\mathbb{A}} \Psi_X(c, u c) \psi(-u) \, d u \\
    &= \sum_{c \in F^\times} \prod_{\mathfrak{p}} \int_{F_{\mathfrak{p}}} \Psi_{\mathfrak{p}, X}(c, y) \psi_{\mathfrak{p}}(- y / c) \, d y.
  \end{aligned}
\end{equation}
In the last step, we substituted $y = u c$ and applied the product formula.  All sums and integrals converge absolutely.  Indeed, the same unfolding gives
\begin{equation*}
  \sum_{c \in F^\times} \int_{\mathbb{A}} \lvert \Psi_X(c, u c) \rvert \, d u
  = \int_{\mathbb{A} / F} \sum_{c \in F^\times} \sum_{y \in F} \lvert (n(u) \Psi)_X(c, y) \rvert \, d u.
\end{equation*}
For $u$ in a compact fundamental domain, the pure tensors $n(u) \Psi$ satisfy the bounds \eqref{eq:local-bounds} of Appendix \ref{sec:proof-poisson} with uniform constants, so the right-hand side is finite by Lemma \ref{lemma:theta-bounds}.

The Fourier coefficient is thus a sum over $c \in F^\times$ of products of local integrals.  We study those integrals in \S\S\ref{sec:jacquet}--\ref{sec:jacquet-gauss}, and apply \eqref{eq:fourier-coefficient} to the two sides of the Poisson formula in \S\ref{sec:line-formula}.  In \S\S\ref{sec:local-computations}--\ref{sec:quadratic-dual}, we make the resulting formula explicit as a summation formula for Gauss sums.

\subsection{Jacquet integrals}\label{sec:jacquet}
Let $F$ be a local field containing $\mu_n$, and let $\psi$ be a nontrivial unitary character of $F$.  The local integrals in \eqref{eq:fourier-coefficient} are the following.  On homogeneous functions, they are the functionals considered by Deligne \cite[\S\S6.4--6.5]{SB_1978-1979__21__244_0}.
\begin{definition}
  For $\Phi \in \mathcal{S}$ and $c \in F^\times$, we set
  \begin{equation*}
    J \Phi(c) := \int_F \Phi_X(c, y) \psi(- y / c) \, d y.
  \end{equation*}
  We write
  \begin{equation*}
    \mathcal{W} := J \mathcal{S}
  \end{equation*}
  for the space of the functions $J \Phi$ on $F^\times$.
\end{definition}
The integral converges absolutely, and $J \Phi$ is continuous: the function $\Phi_X$ is continuous on $U_X$ and, as noted in \eqref{eq:growth}, we have $\Phi_X(c, y) \ll (1 + \lvert y \rvert)^{-2}$ for $c$ in a compact subset of $F^\times$.  The substitution $y \mapsto y - u c$ gives
\begin{equation}\label{eq:jacquet-equivariance}
  J(n(u) \Phi) = \psi(u) J \Phi
  \quad \text{ for } u \in F.
\end{equation}

Every test function on $F^\times$ is a Jacquet integral.
\begin{lemma}\label{lemma:jacquet-surjective}
  Let $W \in C_c^\infty(F^\times)$.  Choose $\beta \in C_c^\infty(F)$ with $\int_F \beta(u) \psi(-u) \, d u = 1$.  The formulas
  \begin{equation*}
    \Phi_X(x, y) := \lvert x \rvert^{-1} W(x) \beta(y/x),
    \qquad
    \Phi_Y(x, y) := \{y, x\}_n \Phi_X(x, y)
  \end{equation*}
  (with $\Phi_Y(0, y) := 0$) define an element $\Phi \in \mathcal{D}$ with $J \Phi = W$, such that $\Phi_X$ vanishes outside a compact subset of $U_X$.  With $P$ as in \S\ref{sec:germs}, we have
  \begin{equation}\label{eq:lift-average}
    (P \Phi)_X(x, y) = \frac{\beta(y/x)}{\lvert x \rvert^{1 + 1/n}} \int_{x F^{\times n}} W(t) \lvert t \rvert^{1/n} \, d^\times t.
  \end{equation}
  In particular, if the integral of $W(x) \lvert x \rvert^{1/n}$ over each coset of $F^{\times n}$ in $F^\times$ vanishes, then $\Phi$ satisfies the vanishing condition.
\end{lemma}
\begin{proof}
  The first assertion is immediate.  For the second, let $a \in F^{\times n}$.  Then $\Phi_X(a x, a y) = \lvert a x \rvert^{-1} W(a x) \beta(y / x)$.  We multiply by $\lvert a \rvert^{1 + 1/n}$, integrate over $a$, and substitute $t = a x$.
\end{proof}

\subsection{The Fourier transform on Jacquet integrals}\label{sec:line-transform}
Let $\Phi \in \mathcal{S}$.  We express $J(\mathcal{F} \Phi)$ in terms of $W := J \Phi$.  Fix $c \in F^\times$ and $t \in F$.  Before integrating against $\psi(t)$, consider the Jacquet integral at $c$ of the line integrals with parameter $t$.  Formally, Definition \ref{definition:lines} gives
\begin{equation*}
  \int_F R_X \Phi((c, y'), t) \psi(-y'/c) \, d y'
  = \int_{F^2} \Phi_X \left(x, \frac{x y' - t^n}{c}\right) \{x, -c\}_n \psi(-y'/c) \, \frac{d x \, d y'}{\lvert c \rvert}.
\end{equation*}
Change variables from $y'$ to $y = (x y' - t^n)/c$.  Then $d x \, d y' / \lvert c \rvert = d y \, d^\times x$, and
\begin{equation*}
  y' = \frac{c y + t^n}{x},
  \qquad
  \psi(- y' / c) = \psi(- y / x) \, \psi(- t^n / c x).
\end{equation*}
Integrating $\Phi_X(x, y) \psi(-y/x)$ over $y$ gives $W(x)$.  The remaining integral is therefore
\begin{equation}\label{eq:jacquet-kernel}
  H_c(t) := \int_{F^\times} W(x) \{x, -c\}_n \psi(- t^n / c x) \, d^\times x.
\end{equation}
For homogeneous functions and $t = 1$, this is Deligne's calculation \cite[(6.5.4)]{SB_1978-1979__21__244_0} of the action of the standard intertwining operator on the functionals of \S\ref{sec:jacquet}; see also \cite{Szpruch2019, GaoShahidiSzpruch2023}.  The Fourier transform in $t$ is the remaining operation in our normalization, and we are led to expect that
\begin{equation}\label{eq:jacquet-formal}
  J(\mathcal{F} \Phi)(c) = \lvert n \rvert^{-1/2} \int_F \psi(t) H_c(t) \, d t.
\end{equation}

The absolute integral over $x$ and $y$ at fixed $t$ is
\begin{equation}\label{eq:x-integrable}
  \int_{F^\times} \int_F \lvert \Phi_X(x, y) \rvert \, d y \, d^\times x,
\end{equation}
which can diverge near the line $x = 0$.  To justify the calculation, we modify $\Phi$ near that line.  Choose $u \in F$ with $\psi(u) \neq 1$, and set
\begin{equation}\label{eq:shear-difference}
  \Delta := \frac{n(u) - 1}{\psi(u) - 1},
  \qquad
  \Xi := \Delta^3 \Phi.
\end{equation}
Then $\Xi \in \mathcal{S}$ (Corollary \ref{corollary:metaplectic}).  By \eqref{eq:jacquet-equivariance}, we have $J \Delta \Psi = J \Psi$ for $\Psi \in \mathcal{S}$, and $\Delta$ commutes with $\mathcal{F}$ on $\mathcal{S}$ (Theorem \ref{theorem:schwartz-space}).  Thus
\begin{equation}\label{eq:shear-jacquet}
  J \Xi = J \Phi,
  \qquad
  J(\mathcal{F} \Xi) = J(\mathcal{F} \Phi).
\end{equation}
Deligne \cite[(5.15.1)]{SB_1978-1979__21__244_0} uses the difference $\Delta$ to regularize these functionals on homogeneous functions.  Here the Jacquet integrals already converge.  The shear fixes the line $x = 0$ and preserves the $Y$-component there, by \eqref{eq:unipotent-axis}, so $\Xi_Y(0, y) = 0$.  On $x = c$ it translates $y$ by $u c$.  Taking three differences gives the estimate below over both $\mathbb{R}$ and $\mathbb{C}$.

\begin{lemma}\label{lemma:shear-integrable}
  Let $\Phi \in \mathcal{S}$, and let $\Xi$ be as in \eqref{eq:shear-difference}.  Then for each $A \geq 0$ and all $c \in F^\times$,
  \begin{equation}\label{eq:shear-bound}
    \int_F \lvert \Xi_X(c, y) \rvert \, d y
    \ll \min \left( \lvert c \rvert^{1/n}, \lvert c \rvert^{-A} \right),
  \end{equation}
  with the implied constant depending only upon $\Phi$, $u$ and $A$.
\end{lemma}
\begin{proof}
  By \eqref{eq:unipotent-charts}, we have
  \begin{equation}\label{eq:third-difference}
    \Xi_X(c, y) = (\psi(u) - 1)^{-3} \sum_{k = 0}^{3} (-1)^{3 - k} \binom{3}{k} \Phi_X(c, y + k u c).
  \end{equation}
  The integral in \eqref{eq:shear-bound} is thus at most a constant multiple of the integral of $\lvert \Phi_X(c, y) \rvert$ over $y$, which by \eqref{eq:growth} is $O(\lvert c \rvert^{1 - N})$ for $\lvert c \rvert \geq 1$ and each $N \geq 2$.  This gives \eqref{eq:shear-bound} for $\lvert c \rvert \geq 1$.  Also by \eqref{eq:growth}, each term in \eqref{eq:third-difference} is $O(\lvert c \rvert^{-1 + 1/n})$, since the points $(c, y + k u c)$ have norm at least $\lvert c \rvert$.  For each $C > 0$, the integral of $\lvert \Xi_X(c, y) \rvert$ over the $y$ with $\lvert y \rvert \leq C \lvert c \rvert$ is thus $O(\lvert c \rvert^{1/n})$.  It remains to bound the integral over the $y$ with $\lvert y \rvert > C \lvert c \rvert$, for some $C$ depending only upon $\Phi$ and $u$.

  Suppose first that $F$ is non-archimedean.  We show that the remaining integral vanishes for suitable $C$.  For $j \in \mathbb{Z}$, set $\Xi_j := \lvert \varpi \rvert^{j (n-1)} \Xi(\varpi^{n j} \, \cdot)$; its $Y$-component again vanishes on the line $x = 0$.  Let $A_0 := \{v : \lvert \varpi \rvert^n < \lVert v \rVert \leq 1\}$.  The function $\Xi$ vanishes outside a bounded set, and agrees near the origin with its leading term $b \Xi$, which satisfies $(b \Xi)(\mu^n \, \cdot) = \lvert \mu \rvert^{1-n} \, b \Xi$ (Proposition \ref{proposition:finite-germs}).  Thus $\Xi_j$ vanishes on $A_0$ when $j$ is sufficiently negative, and agrees with $b \Xi$ on $A_0$ when $j$ is sufficiently large.  The restrictions to $A_0$ of the functions $\Xi_j$ thus form a finite set.  The $Y$-component of each is locally constant, and vanishes on the compact set of points $(0, y)$ of $A_0$; it therefore vanishes at the points $(x, y)$ of $A_0$ with $\lvert x \rvert \leq \delta$, for some $\delta \in (0, \lvert \varpi \rvert^n]$.  Since each $v \in U$ is of the form $\varpi^{n j} v_0$ with $v_0 \in A_0$, the function $\Xi$ vanishes at each $v = (x, y)$ with $\lvert x \rvert \leq \delta \lVert v \rVert$.  In particular, $\Xi_X(c, y) = 0$ when $\lvert y \rvert > \lvert c \rvert / \delta$.

  Suppose now that $F$ is archimedean.  Set $d := [F : \mathbb{R}]$, and let $\ell$ denote the ordinary absolute value on $F$, so that $\lvert x \rvert = \ell(x)^d$.  For $\mu \in F^\times$, set $\Phi_\mu := \lvert \mu \rvert^{n-1} \Phi(\mu^n \, \cdot)$.  Let $K \subset U_Y$ be the compact set of points $(x, y)$ with $\ell(x) \leq 1$ and $1/2 \leq \ell(y) \leq 3/2$.  By Proposition \ref{proposition:schwartz-origin}, the partial derivatives of order at most three of the functions $(\Phi_\mu)_Y$, with respect to the real coordinates, are bounded on $K$, uniformly in $\mu \in F^\times$.

  Let $y \in F$ with $\ell(y) > C \ell(c)$, where $C := \max(1, 6 \ell(u))$.  For $s \in [0, 3]$, we have $\ell(s u c / y) < 1/2$, so $1 + s u c / y$ lies in $F^{\times n}$ (it is positive if $F = \mathbb{R}$), and thus $\{c, y + s u c\}_n = \{c, y\}_n$.  The terms in \eqref{eq:third-difference} are therefore, up to the common factor $\{c, y\}_n$, the values $\Phi_Y(c, y + k u c)$.  Set $r := \ell(y)$ and $\mu := r^{1/n}$, so that $\mu^n = r$ and $\lvert \mu \rvert^{n-1} = \lvert y \rvert^{1 - 1/n}$.  Then
  \begin{equation*}
    \Phi_Y(c, y + s u c) = \lvert y \rvert^{-1 + 1/n} g(s),
    \qquad
    g(s) := (\Phi_\mu)_Y \left( \frac{c}{r}, \frac{y + s u c}{r} \right),
  \end{equation*}
  and the points at which $(\Phi_\mu)_Y$ is evaluated lie in $K$.  The third difference $\sum_{k} (-1)^{3-k} \binom{3}{k} g(k)$ is bounded by the supremum of $\lvert g''' \rvert$ on $[0, 3]$, which is therefore $O(\ell(c / y)^3)$.  Hence
  \begin{equation*}
    \Xi_X(c, y) \ll \lvert y \rvert^{-1 + 1/n} \, \ell(c / y)^3
    = \ell(c)^3 \, \ell(y)^{-d + d/n - 3}.
  \end{equation*}
  Since $d / n < 3$, the integral of the right-hand side over the $y$ with $\ell(y) > C \ell(c)$ is $O(\ell(c)^{d/n}) = O(\lvert c \rvert^{1/n})$.
\end{proof}

Integrating this estimate against $d^\times c$ shows that the integral \eqref{eq:x-integrable} is finite for $\Xi$ in place of $\Phi$.

\begin{corollary}\label{corollary:jacquet-decay}\lean{MetaplecticPoisson.Datum.jacquet_decay_arch}
  Let $W \in \mathcal{W}$.  Then for each $A \geq 0$,
  \begin{equation*}
    W(c) \ll \min \left( \lvert c \rvert^{1/n}, \lvert c \rvert^{-A} \right)
    \quad \text{ for } c \in F^\times,
  \end{equation*}
  with the implied constant depending only upon $W$ and $A$.  In particular, $W$ is integrable with respect to $d^\times c$.
\end{corollary}
\begin{proof}
  Write $W = J \Phi$ with $\Phi \in \mathcal{S}$, and let $\Xi$ be as in \eqref{eq:shear-difference}.  By \eqref{eq:shear-jacquet}, we have $\lvert W(c) \rvert = \lvert J \Xi(c) \rvert \leq \int_F \lvert \Xi_X(c, y) \rvert \, d y$.  We apply Lemma \ref{lemma:shear-integrable}.
\end{proof}

To justify the Fourier transform in \eqref{eq:jacquet-formal} for general $W \in \mathcal{W}$, we first test against a smooth function of compact support in $F^\times$.

\begin{proposition}\label{proposition:jacquet-weak}
  Let $\Phi \in \mathcal{S}$ and $c \in F^\times$, and set $W := J \Phi$.  The integral \eqref{eq:jacquet-kernel} converges absolutely, and defines a bounded continuous function $H_c$ on $F$.  For each $\rho \in C_c^\infty(F^\times)$,
  \begin{equation}\label{eq:jacquet-weak}
    \int_{F} \lvert \mu \rvert^{-1} J(\mathcal{F} \Phi)(\mu^n c) \, \rho(\mu) \, d \mu
    =
    \lvert n \rvert^{-1/2} \int_F H_c(t) \hat{\rho}(t) \, d t.
  \end{equation}
\end{proposition}
\begin{proof}
  The assertions concerning $H_c$ follow from the integrability of $W$ (Corollary \ref{corollary:jacquet-decay}), the continuity by dominated convergence.

  Let $\Xi$ be as in \eqref{eq:shear-difference}.  By Lemma \ref{lemma:shear-integrable}, the integral \eqref{eq:x-integrable} is finite for $\Xi$ in place of $\Phi$; let $C$ denote its value.  The calculation at the beginning of this subsection, applied to $\Xi$, thus involves an absolutely convergent integral over $x$ and $y$.  Since $J \Xi = W$, it gives, for $t \in F$,
  \begin{equation}\label{eq:jacquet-lines}
    \int_F R_X \Xi((c, y'), t) \psi(- y' / c) \, d y' = H_c(t),
    \qquad
    \int_F \lvert R_X \Xi((c, y'), t) \rvert \, d y' \leq C.
  \end{equation}

  Let $\mu \in F^\times$.  By Lemma \ref{lemma:transform-fourier}, the substitution $y \mapsto \mu^n y'$ in the definition of $J(\mathcal{F} \Xi)(\mu^n c)$ gives
  \begin{equation*}
    \lvert \mu \rvert^{-1} J(\mathcal{F} \Xi)(\mu^n c)
    = \lvert n \rvert^{-1/2} \int_F \psi(- y' / c) \left( \int_F R_X \Xi((c, y'), t) \psi(\mu t) \, d t \right) \, d y'.
  \end{equation*}
  The inner integral is $\lvert n \rvert^{1/2} \lvert \mu \rvert^{n-1} (\mathcal{F} \Xi)_X(\mu^n c, \mu^n y')$.  For $\mu$ in the support of $\rho$, it is $O((1 + \lvert y' \rvert)^{-2})$, because $\mathcal{F} \Xi$ satisfies \eqref{eq:growth}.  This bound allows us to integrate over $\mu$ before $y'$.  For each $y'$, the function $R_X \Xi((c, y'), \cdot)$ is integrable by Lemma \ref{lemma:lines}\eqref{item:lines-convergence}, so we may then integrate over $\mu$ before $t$.  This gives
  \begin{equation*}
    \int_{F} \lvert \mu \rvert^{-1} J(\mathcal{F} \Xi)(\mu^n c) \, \rho(\mu) \, d \mu
    = \lvert n \rvert^{-1/2} \int_F \psi(- y' / c) \left( \int_F R_X \Xi((c, y'), t) \hat{\rho}(t) \, d t \right) \, d y'.
  \end{equation*}
  Since $\hat{\rho}$ is integrable, the bound in \eqref{eq:jacquet-lines} shows that this double integral converges absolutely.  Integrating over $y'$ first, we obtain \eqref{eq:jacquet-weak} with $\Xi$ in place of $\Phi$ on the left-hand side.  We conclude by \eqref{eq:shear-jacquet}.
\end{proof}

It follows that $\mathcal{F}$ induces an operator on $\mathcal{W}$.

\begin{corollary}\label{corollary:line-transform}
  Let $\Phi \in \mathcal{S}$ with $J \Phi = 0$.  Then $J(\mathcal{F} \Phi) = 0$.
\end{corollary}
\begin{proof}
  Let $c \in F^\times$.  The function $H_c$ vanishes, so by \eqref{eq:jacquet-weak}, the continuous function $\mu \mapsto \lvert \mu \rvert^{-1} J(\mathcal{F} \Phi)(\mu^n c)$ on $F^\times$ has integral zero against each $\rho \in C_c^\infty(F^\times)$.  It therefore vanishes.  We take $\mu = 1$.
\end{proof}

\begin{definition}\label{definition:transform}
  For $W \in \mathcal{W}$, we define $W^\vee \in \mathcal{W}$ by
  \begin{equation*}
    W^\vee := J(\mathcal{F} \Phi),
  \end{equation*}
  where $\Phi$ is any element of $\mathcal{S}$ with $J \Phi = W$.  We write $W^{\vee, \psi}$ for $W^\vee$ when we wish to indicate its dependence upon $\psi$.
\end{definition}
By Corollary \ref{corollary:line-transform}, the function $W^\vee$ does not depend upon the choice of $\Phi$.  It satisfies the bound of Corollary \ref{corollary:jacquet-decay}.  Set $\varepsilon := (-1)^{n+1}$.  We have $\mathcal{F}^2 \Phi = h(\varepsilon) \Phi$ (Theorem \ref{theorem:schwartz-space}), and Lemma \ref{lemma:torus} gives $J(h(\varepsilon) \Phi)(c) = \{c, \varepsilon\}_n \{-1, \varepsilon\}_n J \Phi(\varepsilon c)$.  Hence
\begin{equation}\label{eq:jacquet-inversion}
  W^{\vee \vee}(c) = \{-c, \varepsilon\}_n \, W(\varepsilon c).
\end{equation}
In particular, the transform $W \mapsto W^\vee$ maps $\mathcal{W}$ bijectively onto itself.

In this notation, \eqref{eq:jacquet-weak} says that, for each $W \in \mathcal{W}$ and $c \in F^\times$, the function $\lambda \mapsto \lvert \lambda \rvert^{-1} W^\vee(c \lambda^n)$ on $F^\times$ is $\lvert n \rvert^{-1/2}$ times the Fourier transform of $H_c$, in the sense of distributions.  For $W \in C_c^\infty(F^\times)$ the identity holds pointwise.

\begin{proposition}\label{proposition:transform-formula}\lean{NFoldCover.dualWeight}\lean{MetaplecticPoisson.Datum.root_fourier}
  Let $W \in C_c^\infty(F^\times)$ and $c \in F^\times$.  Then the function $H_c$ defined by \eqref{eq:jacquet-kernel} is a Schwartz--Bruhat function on $F$, and for $\lambda \in F^\times$,
  \begin{equation}\label{eq:root-fourier}
    W^\vee(c \lambda^n) = \lvert n \rvert^{-1/2} \lvert \lambda \rvert \, \hat{H}_c(\lambda).
  \end{equation}
  In particular,
  \begin{equation}\label{eq:transform-formula}
    W^\vee(c) = \lvert n \rvert^{-1/2} \int_F \psi(t) \left(
      \int_{F^\times} W(x)
      \{x, -c\}_n \psi(-t^n / c x)
      \, d^\times x
    \right) \, d t.
  \end{equation}
  For fixed $c$, as $W$ varies over a bounded subset of $C_c^\infty(F^\times)$, the functions $H_c$ form a bounded set of Schwartz--Bruhat functions.
\end{proposition}
\begin{proof}
  The function $z \mapsto \int_{F^\times} W(x) \{x, -c\}_n \psi(- z / c x) \, d^\times x$ is the Fourier transform of $x \mapsto \lvert x \rvert^{-1} W(1/x) \{x, -c\}_n^{-1}$, evaluated at $- z / c$.  The latter function lies in $C_c^\infty(F^\times)$, and remains in a bounded subset when $W$ does.  This gives the assertions concerning $H_c$.

  We have $W \in \mathcal{W}$, by Lemma \ref{lemma:jacquet-surjective}.  Since $H_c$ is integrable, the right-hand side of \eqref{eq:jacquet-weak} is the integral against $\rho(\mu)$ of the continuous function $\lvert n \rvert^{-1/2} \hat{H}_c$.  Since \eqref{eq:jacquet-weak} holds for every $\rho$, and $W^\vee$ is continuous, we obtain \eqref{eq:root-fourier}.  The case $\lambda = 1$ is \eqref{eq:transform-formula}.
\end{proof}
\begin{example}
  If $n = 1$, then $W^\vee(c) = W(1/c)$.
\end{example}

By \eqref{eq:root-fourier}, the transform $W^\vee$ is a Fourier transform in an $n$th root variable, as is $\mathcal{F}$.  For test functions $W$, this makes the bound of Corollary \ref{corollary:jacquet-decay} for $W^\vee$ uniform.

\begin{corollary}\label{corollary:decay}
  Let $\mathcal{B}$ be a bounded subset of $C_c^\infty(F^\times)$.  Then for each $A \geq 0$, we have
  \begin{equation*}
    W^\vee(c) \ll \min \left( \lvert c \rvert^{1/n}, \lvert c \rvert^{-A} \right)
  \end{equation*}
  for all $W \in \mathcal{B}$ and $c \in F^\times$, with the implied constant depending only upon $\mathcal{B}$ and $A$.  If $F$ is non-archimedean, then there is a compact subset of $F$, depending only upon $\mathcal{B}$, outside of which $W^\vee$ vanishes for every $W \in \mathcal{B}$.
\end{corollary}
\begin{proof}
  Each $c \in F^\times$ is of the form $d \lambda^n$, with $d$ in a fixed finite set of representatives for $F^\times / F^{\times n}$, and then $\lvert \lambda \rvert \asymp \lvert c \rvert^{1/n}$.  By Proposition \ref{proposition:transform-formula}, we have $W^\vee(c) = \lvert n \rvert^{-1/2} \lvert \lambda \rvert \hat{H}_d(\lambda)$, where the functions $H_d$ attached to the elements of $\mathcal{B}$ form a bounded set of Schwartz--Bruhat functions.  Their Fourier transforms are uniformly bounded and decay rapidly, uniformly; in the non-archimedean case, they vanish outside a fixed compact set.
\end{proof}

At a finite place, we may describe the elements of $\mathcal{W}$ near zero more precisely.

\begin{lemma}\label{lemma:jacquet-germ}
  Suppose that $F$ is non-archimedean.
  \begin{enumerate}[(i)]
  \item\label{item:jacquet-test} We have $J \mathcal{D} = C_c^\infty(F^\times)$.
  \item\label{item:jacquet-germ} Let $W \in \mathcal{W}$.  Then $W$ is locally constant, and vanishes at $c$ when $\lvert c \rvert$ is sufficiently large.  For each $d \in F^\times$, there is a constant $C_d$ such that $W(d \lambda^n) = C_d \lvert \lambda \rvert$ when $\lvert \lambda \rvert$ is sufficiently small.  In particular, $W(\varpi^n c) = \lvert \varpi \rvert \, W(c)$ when $\lvert c \rvert$ is sufficiently small.
  \end{enumerate}
\end{lemma}
\begin{proof}
  Let $\Psi \in \mathcal{S}$.  The function $\Psi_X$ is locally constant on $U_X$, and vanishes outside a bounded set (Proposition \ref{proposition:finite-germs}).  Hence $J \Psi$ is locally constant, and vanishes at $c$ when $\lvert c \rvert$ is sufficiently large.

  \eqref{item:jacquet-test} By Lemma \ref{lemma:jacquet-surjective}, it remains to show that for $\Phi \in \mathcal{D}$, we have $J \Phi(c) = 0$ when $\lvert c \rvert$ is sufficiently small.  Let $\Delta$ be as in \eqref{eq:shear-difference}.  Then $J \Delta \Phi = J \Phi$, and $\Delta \Phi \in \mathcal{D}$.  The $Y$-component of $\Delta \Phi$ is locally constant and vanishes on the line $x = 0$, by \eqref{eq:unipotent-axis}.  The support of $\Delta \Phi$ is thus a compact subset of $U_X$.  Its projection onto the $x$-axis is a compact subset of $F^\times$, outside of which $J \Delta \Phi$ vanishes.

  \eqref{item:jacquet-germ} By \eqref{item:jacquet-test} and the definition of the transform, we may write $W = W_1 + W_2^\vee$ with $W_1, W_2 \in C_c^\infty(F^\times)$.  By Proposition \ref{proposition:transform-formula}, we have $W_2^\vee(d \lambda^n) = \lvert n \rvert^{-1/2} \lvert \lambda \rvert \hat{H}_d(\lambda)$, where $H_d$ is the Schwartz--Bruhat function attached to $W_2$ by \eqref{eq:jacquet-kernel}.  Its Fourier transform is constant near zero.  This gives the assertion concerning $W(d \lambda^n)$, with $C_d = \lvert n \rvert^{-1/2} \hat{H}_d(0)$.  The last assertion follows upon taking $d$ in a finite set of representatives for $F^\times / F^{\times n}$.
\end{proof}

\subsection{The basic function}\label{sec:jacquet-gauss}
Suppose that $F$ is non-archimedean, of residue characteristic not dividing $n$.
\begin{definition}\label{definition:basic-function}
  The \emph{basic function} is the Jacquet integral of the basic vector:
  \begin{equation*}
    W_\psi^0 := J \Phi^0 \in \mathcal{W}.
  \end{equation*}
\end{definition}
Suppose now that $(F, \psi, n)$ is good (Definition~\ref{definition:good}).  As anticipated in \S\ref{sec:introduction}, the basic function is then given by Gauss sums.

\begin{definition}\label{definition:gauss-ramanujan}
  For $(F, \psi, n)$ good and $c \in F^\times$, we define the Gauss--Ramanujan sum
  \begin{equation*}
    g_{\psi}(c) :=
    \begin{cases}
      1
      & \text{ if } c \in \mathfrak{o}^\times, \\
      \sum_{d \in(\mathfrak{o} / c)^\times} \{c, d\}_n \psi(- d/c)
      & \text{ if } c \in \mathfrak{o} - \mathfrak{o}^\times - \{0\}, \\
      0    
      & \text{ if } c \notin \mathfrak{o}.
    \end{cases}
  \end{equation*}
\end{definition}
\begin{lemma}\label{lemma:jacquet-basic}
  Suppose $(F, \psi, n)$ is good.  Then, for $c \in F^\times$,
  \begin{equation*}
    W_\psi^0(c) = \vol(\mathfrak{o}) \lvert c \rvert \sum_{k \geq 0} q^{(n - 1) k} g_\psi(c / \varpi^{n k}).
  \end{equation*}
\end{lemma}
\begin{proof}
  We split the integral defining $J \Phi^0(c)$ according to the integer $k \geq 0$ for which $(c, y) \in \varpi^{n k} \cdot (\text{primitive})$.  Writing $c' := c / \varpi^{n k}$ and $y = \varpi^{n k} y'$, the contribution of $k$ is
  \begin{equation*}
    q^{(n-1) k} q^{- n k} \int \Phi^0_X(c', y') \psi(- y' / c') \, d y',
  \end{equation*}
  where $y'$ ranges over $\mathfrak{o}$ if $c' \in \mathfrak{o}^\times$, over $\mathfrak{o}^\times$ if $c' \in \mathfrak{o} - \mathfrak{o}^\times$, and over the empty set otherwise.  In the first case, the integrand is $1$.  In the second, it is $\{c', y'\}_n \psi(- y'/c')$, which depends only upon the class of $y'$ modulo $c'$.  In either case, the integral is $\vol(\mathfrak{o}) \lvert c' \rvert g_\psi(c')$.  Summing over $k$ gives the formula.
\end{proof}
\begin{proposition}\label{proposition:basic-function-dual}
  Suppose $(F, \psi, n)$ is good.  Choose $b \in \mathfrak{o} - \{0\}$ so that $b^{-1} \mathfrak{o}$ is the largest fractional ideal on which $\psi$ is trivial.  Then
  \begin{equation*}
    (W_\psi^0)^\vee(c) = W^0_\psi(b^n c).
  \end{equation*}
  In particular, $(W_\psi^0)^\vee = W_\psi^0$ if $\psi$ is unramified.
\end{proposition}
\begin{proof}
  By definition, $(W_\psi^0)^\vee = J(\mathcal{F} \Phi^0)$.  By Corollary \ref{corollary:basic-vector-general}, the substitution $y \mapsto y / b^n$ gives $J(\mathcal{F} \Phi^0)(c) = J \Phi^0(b^n c)$.
\end{proof}

\subsection{The main formula}\label{sec:line-formula}
We return to the global setting of \S\ref{sec:main-formula}.  Let $W : \mathbb{A}^\times \rightarrow \mathbb{C}$ be a function of the form
\begin{equation*}
  W(x) = \prod_{\mathfrak{p}} W_{\mathfrak{p}}(x_{\mathfrak{p}}),
\end{equation*}
with $W_{\mathfrak{p}} \in C_c^\infty(F_{\mathfrak{p}}^\times)$ for each $\mathfrak{p} \in S$ and $W_{\mathfrak{p}} = W_{\psi_\mathfrak{p}}^0$ for each $\mathfrak{p} \notin S$.  We define
\begin{equation*}
  W^\vee(x) := \prod_{\mathfrak{p}} W_{\mathfrak{p}}^\vee(x_{\mathfrak{p}}),
\end{equation*}
where each local transform is taken with respect to $\psi_{\mathfrak{p}}$.

\begin{theorem}\label{theorem:main}\lean{MetaplecticPoisson.Datum.lineSummation}
  Let $W$ be as above.  The sums of $W(c)$ and $W^\vee(c)$ over $c \in F^\times$ converge absolutely.  There are unique constants $\rho_S(\beta)$, independent of $W$ and indexed by the tuples $\beta = (\beta_{\mathfrak{p}})_{\mathfrak{p} \in S}$ of cosets $\beta_{\mathfrak{p}}$ of $F_{\mathfrak{p}}^{\times n}$ in $F_{\mathfrak{p}}^\times$, such that
  \begin{equation*}
    \sum_{c \in F^\times} W(c)
    =
    \sum_{c \in F^\times} W^\vee(c)
    +
    \sum_{\beta} \rho_S(\beta)
    \prod_{\mathfrak{p} \in S}
    \int_{\beta_{\mathfrak{p}}} W_{\mathfrak{p}}(x)
    \lvert x \rvert_{\mathfrak{p}}^{1/n} \, d^\times x.
  \end{equation*}
\end{theorem}
The constants $\rho_S(\beta)$ depend upon $\psi$ and $S$.  The sum over $\beta$, which we call the \emph{main term}, is the Fourier coefficient of the boundary term in Theorem \ref{theorem:poisson}; see \eqref{eq:main-term}.  It vanishes if $n = 1$.  It also vanishes if, for some $\mathfrak{p} \in S$, the integrals of $W_{\mathfrak{p}}(x) \lvert x \rvert_{\mathfrak{p}}^{1/n}$ over the cosets of $F_{\mathfrak{p}}^{\times n}$ all vanish.  For instance, if $\mathfrak{p}$ is complex, then $F_{\mathfrak{p}}^{\times n} = F_{\mathfrak{p}}^\times$, and this holds if $W_{\mathfrak{p}}(e^{i \theta} x) = e^{i \ell \theta} W_{\mathfrak{p}}(x)$ for some integer $\ell \neq 0$.
\begin{proof}
  For $\mathfrak{p} \in S$, let $\Phi_{\mathfrak{p}} \in \mathcal{D}$ be attached to $W_{\mathfrak{p}}$ by Lemma \ref{lemma:jacquet-surjective}; for $\mathfrak{p} \notin S$, let $\Phi_{\mathfrak{p}} := \Phi_{\mathfrak{p}}^0$.  Let $\Phi := \otimes_{\mathfrak{p}} \Phi_{\mathfrak{p}}$.  By Lemma \ref{lemma:jacquet-surjective} and the definitions of the basic function and of the transform, we have $J \Phi_{\mathfrak{p}} = W_{\mathfrak{p}}$ and $J(\mathcal{F} \Phi_{\mathfrak{p}}) = W_{\mathfrak{p}}^\vee$ for all $\mathfrak{p}$.  By \eqref{eq:fourier-coefficient}, applied to $\Phi$ and to $\mathcal{F} \Phi$, the sums of $W(c)$ and of $W^\vee(c)$ converge absolutely, and their difference is
  \begin{equation*}
    \int_{\mathbb{A} / F} \Bigl( \sum_{v \in F^2 - \{0\}} (n(u) \Phi)(v) - \sum_{v \in F^2 - \{0\}} (n(u) \mathcal{F} \Phi)(v) \Bigr) \psi(-u) \, d u.
  \end{equation*}

  Let $u \in \mathbb{A}$.  We apply Theorem \ref{theorem:poisson} to $n(u) \Phi$.  We have $\mathcal{F} n(u) \Phi = n(u) \mathcal{F} \Phi$ (Theorem \ref{theorem:schwartz-space}).  Since $S$ is nonempty and $n(u_{\mathfrak{p}}) \Phi_{\mathfrak{p}} \in \mathcal{D}$ for $\mathfrak{p} \in S$, we have $b (n(u) \Phi) = 0$.  The difference of the two sums in the integrand is thus $B(b \, n(u) \mathcal{F} \Phi)$, and so
  \begin{equation}\label{eq:main-term}
    \sum_{c \in F^\times} W(c) = \sum_{c \in F^\times} W^\vee(c) + M(W),
    \qquad
    M(W) := \int_{\mathbb{A} / F} B(b \, n(u) \mathcal{F} \Phi) \, \psi(-u) \, d u.
  \end{equation}
  For $n = 1$, the leading term $b \, n(u) \mathcal{F} \Phi$ does not depend upon $u$, since the local leading terms are then constant functions (\S\ref{sec:germs}); hence $M(W) = 0$.

  It remains to express $M(W)$ in terms of the integrals
  \begin{equation*}
    I_{\mathfrak{p}}(\beta_{\mathfrak{p}}) := \int_{\beta_{\mathfrak{p}}} W_{\mathfrak{p}}(x) \lvert x \rvert_{\mathfrak{p}}^{1/n} \, d^\times x.
  \end{equation*}
  At each $\mathfrak{p} \in S$, fix the auxiliary function in Lemma \ref{lemma:jacquet-surjective} independently of $W_{\mathfrak{p}}$.  By \eqref{eq:lift-average} and \eqref{eq:germ-via-average}, $b \mathcal{F} \Phi_{\mathfrak{p}}$ depends linearly on the vector of integrals $I_{\mathfrak{p}}(\beta_{\mathfrak{p}})$.  Since $b$ commutes locally with $n(u_{\mathfrak{p}})$, \eqref{eq:main-term} expresses $M(W)$ as a multilinear form in these vectors.  Each vector is arbitrary, since each coset supports a test function of weighted integral one.  The coefficients of the form in these coordinates are the required constants $\rho_S(\beta)$.
\end{proof}

\subsection{Local weights and their transforms}\label{sec:local-computations}
To make Theorem \ref{theorem:main} explicit, we must describe the factors $W_{\mathfrak{p}}^\vee$ and $W_{\psi_{\mathfrak{p}}}^0$.  We record how the transform $W \mapsto W^\vee$ behaves under dilation and change of character, describe it by Mellin transforms, evaluate the basic function, and compute the transforms of characters of the unit group, which will serve to twist the sums.  In this subsection, $F$ is a local field containing $\mu_n$, and $\psi$ is a nontrivial unitary character of $F$.

\subsubsection{Dilation and change of character}\label{sec:dilation-and-change-of-character}
In applications, the test functions are dilated, and the character $\psi$ is replaced by $\psi(b \, \cdot)$.  We record the effect on the transform.

\begin{lemma}\label{lemma:dilation}
  Let $W \in C_c^\infty(F^\times)$ and $X \in F^\times$.  Then
  \begin{equation*}
    \left( W(\cdot / X) \right)^\vee(c) = \{X, -c\}_n \left( W \{\cdot, X\}_n^{-1} \right)^\vee(X c).
  \end{equation*}
\end{lemma}
\begin{proof}
  We substitute $x \mapsto X x$ in the inner integral of \eqref{eq:transform-formula} and use that
  \begin{equation*}
    \{X x, -c\}_n = \{X, -c\}_n \{x, - X c\}_n \{x, X\}_n^{-1}.
    \qedhere
  \end{equation*}
\end{proof}

\begin{lemma}\label{lemma:change-of-character}
  Let $W \in C_c^\infty(F^\times)$ and $b \in F^\times$.  Set $\psi_b := \psi(b \, \cdot)$.  Then
  \begin{equation*}
    W^{\vee, \psi_b}(c) = \left( W \{\cdot, b\}_n \right)^{\vee, \psi}(b^{n-1} c).
  \end{equation*}
\end{lemma}
\begin{proof}
  Replacing $\psi$ by $\psi_b$ in \eqref{eq:transform-formula} multiplies each of the two measures by $\lvert b \rvert^{1/2}$.  The substitution $t \mapsto t/b$ cancels their product and changes the inner phase to $-t^n/(b^{n-1} c x)$.  We conclude using $\{x, -c\}_n = \{x, -b^{n-1} c\}_n \{x, b\}_n$.
\end{proof}

\begin{corollary}\label{corollary:uniform-dilates}
  Let $\mathcal{B}$ be a bounded subset of $C_c^\infty(F^\times)$.  Then for each $A \geq 0$, we have
  \begin{equation*}
    \left( W(\cdot / X) \right)^\vee(c) \ll \min \left( \lvert X c \rvert^{1/n}, \lvert X c \rvert^{-A} \right)
  \end{equation*}
  for all $W \in \mathcal{B}$, $X \in F^\times$ and $c \in F^\times$, with the implied constant depending only upon $\mathcal{B}$ and $A$.
\end{corollary}
\begin{proof}
  By Lemma \ref{lemma:dilation},
  \begin{equation*}
    \lvert (W(\cdot / X))^\vee(c) \rvert = \lvert (W \{\cdot, X\}_n^{-1})^\vee(X c) \rvert.
  \end{equation*}
  The character $\{\cdot, X\}_n$ depends only upon the class of $X$ in the finite group $F^\times / F^{\times n}$, and is locally constant.  The functions $W \{\cdot, X\}_n^{-1}$ with $W \in \mathcal{B}$ and $X \in F^\times$ thus form a bounded subset of $C_c^\infty(F^\times)$, and we conclude by Corollary \ref{corollary:decay}.
\end{proof}

\subsubsection{The Mellin transform}\label{sec:mellin}
For a continuous function $f$ on $F^\times$ and a unitary character $\omega$ of $F^\times$, we write
\begin{equation*}
  Z(f, \omega, z) := \int_{F^\times} f(\lambda) \omega(\lambda) \lvert \lambda \rvert^{z} \, d^\times \lambda,
\end{equation*}
when this converges absolutely.  We recall Tate's local functional equation \cite{MR0217026}: there is a meromorphic function $\gamma(z, \omega, \psi)$ such that for each Schwartz--Bruhat function $g$ on $F$ and $0 < \Re(z) < 1$,
\begin{equation}\label{eq:tate}
  Z(\hat{g}, \omega^{-1}, 1 - z) = \gamma(z, \omega, \psi) Z(g, \omega, z).
\end{equation}

Each $W \in C_c^\infty(F^\times)$ is the sum of its restrictions to the finitely many classes $\beta F^{\times n}$.  For $W$ supported on one class, the Hilbert symbol $\{x, -c\}_n$ in the formula \eqref{eq:transform-formula} for $W^\vee$ equals $\{\beta, -c\}_n$ on the support of $W$, so it may be taken outside the integral over $x$.  As a function of $c$, it changes the Mellin character from $\omega$ to $\eta := \omega \, \{\beta, \cdot\}_n$.  The Fourier transform in $1/x$ contributes the Tate factor for $\eta$, and the Fourier transform in the $n$th root variable $t$ contributes the one for $\omega^{-n}$.
\begin{proposition}\label{proposition:mellin}
  Let $\beta \in F^\times$, and let $W \in C_c^\infty(F^\times)$ be supported on the class $\beta F^{\times n}$.  Let $\omega$ be a unitary character of $F^\times$, and set $\eta := \omega \, \{\beta, \cdot\}_n$.  Then for $-1/n < \Re(s) < 0$,
  \begin{equation*}
    Z(W^\vee, \omega, s) = \lvert n \rvert^{-1/2} \, \eta(-1) \, \gamma(- n s, \omega^{-n}, \psi) \, \gamma(1 + s, \eta, \psi) \, Z(W, \omega^{-1}, -s).
  \end{equation*}
\end{proposition}
\begin{proof}
  The integral defining $Z(W^\vee, \omega, s)$ converges absolutely, by Corollary \ref{corollary:decay}.  Set
  \begin{equation*}
    V(z) := \int_{F^\times} W(x) \psi(- z / x) \, d^\times x.
  \end{equation*}
  We have $V(z) = \hat{g}(-z)$, where $g(x) := \lvert x \rvert^{-1} W(1/x)$, so $V$ is a Schwartz--Bruhat function on $F$.  We have $Z(g, \eta, 1 - z) = Z(W, \eta^{-1}, z)$.  By \eqref{eq:tate}, we thus have for $0 < \Re(z) < 1$
  \begin{equation}\label{eq:mellin-first}
    Z(V, \eta^{-1}, z) = \eta(-1) \, \gamma(1 - z, \eta, \psi) \, Z(W, \eta^{-1}, z).
  \end{equation}

  Let $d \in F^\times$, and set $V_d(\tau) := V(\tau^n / d)$.  For $x$ in the support of $W$, we have $\{x, -d\}_n = \{\beta, -d\}_n$, so the function $H_d$ of \eqref{eq:jacquet-kernel} is $\{\beta, -d\}_n V_d$.  By \eqref{eq:root-fourier}, we have for $\mu \in F^\times$
  \begin{equation*}
    W^\vee(d \mu^n) = \lvert n \rvert^{-1/2} \{\beta, -d\}_n \lvert \mu \rvert \, \hat{V}_d(\mu).
  \end{equation*}
  By \eqref{eq:haar-on-powers} and \eqref{eq:tate}, the integral over the class $d F^{\times n}$ in $\lvert n \rvert^{1/2} Z(W^\vee, \omega, s)$ is
  \begin{align*}
    \tfrac{\lvert n \rvert}{n} \{\beta, -d\}_n \, \omega(d) \lvert d \rvert^s \,
    &Z(\hat{V}_d, \omega^n, n s + 1) \\
    &= \tfrac{\lvert n \rvert}{n} \{\beta, -d\}_n \, \omega(d) \lvert d \rvert^s \, \gamma(- n s, \omega^{-n}, \psi) \, Z(V_d, \omega^{-n}, - n s) \\
    &= \{\beta, -d\}_n \, \gamma(- n s, \omega^{-n}, \psi) \int_{d^{-1} F^{\times n}} V(z) \, \omega^{-1}(z) \lvert z \rvert^{-s} \, d^\times z.
  \end{align*}
  For $z \in d^{-1} F^{\times n}$, we have $\{\beta, -d\}_n = \{\beta, -1\}_n \{\beta, z\}_n^{-1}$.  Summing over the classes, we obtain
  \begin{equation*}
    \lvert n \rvert^{1/2} Z(W^\vee, \omega, s) = \{\beta, -1\}_n \, \gamma(- n s, \omega^{-n}, \psi) \, Z(V, \eta^{-1}, - s).
  \end{equation*}
  We apply \eqref{eq:mellin-first} with $z = -s$.  For $x$ in the support of $W$, we have $\{\beta, x\}_n = \{\beta, \beta\}_n = \{\beta, -1\}_n$, so $Z(W, \eta^{-1}, -s) = \{\beta, -1\}_n Z(W, \omega^{-1}, -s)$.  Since $\{\beta, -1\}_n^2 = 1$, the proposition follows.
\end{proof}
For $n = 1$, the multiplier is $1$, in accordance with the formula $W^\vee(c) = W(1/c)$.  For $F = \mathbb{C}$, there is a single class, $\eta = \omega$, and $\lvert n \rvert^{-1/2} = 1/n$.  For $p$-adic $F$, compare with Szpruch \cite{Szpruch2019}, who expresses the entries of the local coefficients matrix of a genuine principal series through Tate $\gamma$-factors.

\begin{remark}\label{remark:kubota-transform}
  For $F = \mathbb{C}$ and $n = 3$, the transform $W \mapsto W^\vee$ agrees, after a change of variables, with the transform $f \mapsto f^*$ of Kubota \cite[(5)]{Kubota+1974+180+189}, which acts on functions on $\mathbb{C}$ invariant under $\mu_3$.  Take $\psi(z) := e^{2 \pi i \Re(z)}$, and write $\lvert z \rvert = z \bar{z}$ for the normalized absolute value.  For a $\mu_3$-invariant function $f \in C_c^\infty(\mathbb{C}^\times)$, define $I f \in C_c^\infty(\mathbb{C}^\times)$ by
  \begin{equation*}
    (I f)(x) := \lvert x \rvert^{1/3} f(x^{1/3} / \sqrt{3}).
  \end{equation*}
  Then $(I f)^\vee = I(f^*)$.  Indeed, let $\omega(z) := (z / \lvert z \rvert^{1/2})^m$ with $m \in \mathbb{Z}$, and set $u := \lvert m \rvert / 2$.  By the explicit formula for the local factors over $\mathbb{C}$ \cite{MR0217026} and the triplication formula for the gamma function, the factor relating $Z(W^\vee, \omega, s)$ to $Z(W, \omega^{-1}, -s)$ in Proposition \ref{proposition:mellin} is
  \begin{equation*}
    3^{6 s} (-1)^m \pi^{-4 s}
    \frac{\Gamma(\tfrac{1}{3} + s + u) \, \Gamma(\tfrac{2}{3} + s + u)}
    {\Gamma(\tfrac{1}{3} - s + u) \, \Gamma(\tfrac{2}{3} - s + u)}.
  \end{equation*}
  The factor following $3^{6 s}$ is the one relating the Mellin transforms of $f$ and $f^*$ in angular weight $3 m$, at the radial parameter $2 + 6 s$ \cite[p.~187]{Kubota+1974+180+189}.  The dilation by $1 / \sqrt{3}$ in the definition of $I$ accounts for $3^{6 s}$.
\end{remark}

\subsubsection{Gauss--Ramanujan sums}\label{sec:gauss-ramanujan-sums}
We evaluate the basic function.

\begin{lemma}\label{lemma:gauss-ramanujan}
  Suppose $(F, \psi, n)$ is good (Definition~\ref{definition:good}).  For $c \in F^\times$ and $u \in \mathfrak{o}^\times$, we have
  \begin{equation*}
    g_\psi(u c) = \{c, u\}_n \, g_\psi(c),
    \qquad
    W_\psi^0(u c) = \{c, u\}_n W_\psi^0(c).
  \end{equation*}
  In particular, both functions are invariant under $(\mathfrak{o}^\times)^n$.
\end{lemma}
\begin{proof}
  For the first identity, we substitute $d \mapsto u d$ and use that $\{u c, u d\}_n = \{c, u\}_n \{c, d\}_n$, by property \eqref{item:hilbert-units} of the Hilbert symbol.  The second follows from the first and Lemma \ref{lemma:jacquet-basic}, since $\{c / \varpi^{n k}, u\}_n = \{c, u\}_n$.
\end{proof}

\begin{example}\label{example:unramified-basic-function}
  Suppose $(F, \psi, n)$ is good and $\psi$ is unramified.  Then $g_\psi(c) = 1$ if $v(c) = 0$, and $g_\psi(c) = 0$ if $v(c) \geq 2$.  If $v(c) = 1$, then $g_\psi(c) = -1$ for $n = 1$, while $\lvert g_\psi(c) \rvert = q^{1/2}$ for $n \geq 2$.  Consequently, if $n \geq 2$, then writing $c = u \varpi^{n k + r}$ with $u \in \mathfrak{o}^\times$ and $0 \leq r < n$, we have
  \begin{equation*}
    W_\psi^0(c)
    = 1_{k \geq 0} \, q^{- k - r} \left( 1_{r = 0} + 1_{r = 1} g_\psi(u \varpi) \right),
  \end{equation*}
  while if $n = 1$, then $W_\psi^0 = 1_{\mathfrak{o}^\times}$.
\end{example}

\subsubsection{Characters of the unit group}\label{sec:twisted-vectors}
To twist the sums in Theorem \ref{theorem:main} by characters, we will take some components of $W$ to be characters of the unit group.  We compute their transforms.  We assume here that $(F, \psi, n)$ is good (Definition~\ref{definition:good}) and that $\psi$ is unramified.  For a character $\theta$ of $\mathfrak{o}^\times$ and $\alpha \in F^\times$, we set
\begin{equation*}
  G(\theta, \alpha) := \int_{\mathfrak{o}^\times} \theta(x) \psi(\alpha x) \, d x.
\end{equation*}
We have $G(\theta, u \alpha) = \overline{\theta(u)} G(\theta, \alpha)$ for $u \in \mathfrak{o}^\times$.  We recall that if $\theta$ has conductor exponent $e \geq 1$, then
\begin{equation}\label{eq:gauss-sum-facts}
  G(\theta, \alpha) = 0 \text{ unless } v(\alpha) = -e,
  \qquad 
  \lvert G(\theta, \alpha) \rvert = q^{-e/2} \text{ if } v(\alpha) = -e,
\end{equation}
while for the trivial character,
\begin{equation}\label{eq:ramanujan-facts}
  G(1, \alpha) =
  \begin{cases}
    1 - q^{-1} & \text{ if } v(\alpha) \geq 0, \\
    - q^{-1} & \text{ if } v(\alpha) = -1, \\
    0 & \text{ if } v(\alpha) \leq -2.
  \end{cases}
\end{equation}

\begin{lemma}\label{lemma:twisted-reduction}
  Let $\chi$ be a character of $\mathfrak{o}^\times$.  Set $W := 1_{\mathfrak{o}^\times} \chi$.  Then for $c \in F^\times$,
  \begin{equation*}
    W^\vee(c) = \int_F \psi(t) \, G( \overline{\chi} \{c, \cdot\}_n, - t^n / c) \, d t.
  \end{equation*}
  Moreover, $W^\vee(u c) = \overline{\chi(u)} \{c, u\}_n W^\vee(c)$ for $u \in \mathfrak{o}^\times$.
\end{lemma}
\begin{proof}
  For $x \in \mathfrak{o}^\times$, we have $\{1/x, -c\}_n = \{c, x\}_n$, by property \eqref{item:hilbert-units} of the Hilbert symbol.  We substitute $x \mapsto 1/x$ in the inner integral of \eqref{eq:transform-formula}.  The second assertion follows from the first, since $G(\theta, \alpha / u) = \theta(u) G(\theta, \alpha)$.
\end{proof}

For a character $\chi$ of $\mathfrak{o}^\times$, either $\chi^n$ is ramified or $\chi^n = 1$.  We treat the two cases in turn.  In the first, the transform of $1_{\mathfrak{o}^\times} \chi$ is supported on a single valuation:

\begin{proposition}\label{proposition:twisted-vector}
  Let $\chi$ be a character of $\mathfrak{o}^\times$ such that $\chi^n$ is ramified, with conductor exponent $e$.  Set $W := 1_{\mathfrak{o}^\times} \chi$.  Then $W^\vee$ is supported on $\{c : v(c) = - (n-1) e\}$, and for such $c$ and any $t \in F^\times$ with $v(t) = -e$, we have
  \begin{equation*}
    W^\vee(c) =
    q^{e} \,
    G(\chi^n, t) \,
    G( \overline{\chi} \{c, \cdot\}_n, - t^n / c).
  \end{equation*}
  In particular, $\lvert W^\vee(c) \rvert = 1$ for such $c$.
\end{proposition}
\begin{proof}
  Set $\rho_c := \overline{\chi} \{c, \cdot\}_n$.  On each group $1 + \varpi^j \mathfrak{o}$ with $j \geq 1$, the $n$th power map is an automorphism, since the residue characteristic does not divide $n$.  Thus $\chi$ and $\chi^n$ are trivial on the same such groups.  The symbol $\{c, \cdot\}_n$ is trivial on $1 + \varpi \mathfrak{o}$ by property \eqref{item:hilbert-tame} of the Hilbert symbol, so $\rho_c$ and $\chi^n$ are trivial on the same such groups.  Also, $\rho_c^n = \overline{\chi^n} \neq 1$.  Hence $\rho_c$ has conductor exponent $e$, including when $e = 1$.

  By Lemma \ref{lemma:twisted-reduction} and \eqref{eq:gauss-sum-facts}, the integrand vanishes unless $t$ satisfies $n v(t) = v(c) - e$.  Fix such a $t$.  For $s \in \mathfrak{o}^\times$, we have $G(\rho_c, - (s t)^n / c) = \chi^n(s) G(\rho_c, - t^n / c)$, so
  \begin{equation*}
    W^\vee(c) = \lvert t \rvert \, G(\chi^n, t) \, G(\rho_c, - t^n / c).
  \end{equation*}
  By \eqref{eq:gauss-sum-facts} again, this vanishes unless $v(t) = - e$, that is, unless $v(c) = -(n-1) e$.
\end{proof}

\begin{lemma}\label{lemma:twisted-vanishing}
  With notation as in Proposition \ref{proposition:twisted-vector}, the integral of $W(x) \lvert x \rvert^{1/n}$ over each coset of $F^{\times n}$ in $F^\times$ vanishes.
\end{lemma}
\begin{proof}
  The intersection of a coset of $F^{\times n}$ with $\mathfrak{o}^\times$ is either empty or of the form $c (\mathfrak{o}^\times)^n$ with $c \in \mathfrak{o}^\times$, on which $W(c s^n) = \chi(c) \chi^n(s)$.  Since $\chi^n$ is nontrivial on $\mathfrak{o}^\times$, the integral vanishes.
\end{proof}

If $\chi^n = 1$, then $\chi = \{a, \cdot\}_n$ for some $a \in F^\times$, by property \eqref{item:hilbert-tame} of the Hilbert symbol, and the transform of \(1_{\mathfrak o^\times}\chi\) is as follows:

\begin{proposition}\label{proposition:untwisted-vector}
  Let $a \in F^\times$, and set $W := 1_{\mathfrak{o}^\times} \{a, \cdot\}_n$.  Let $c \in F^\times$.
  \begin{enumerate}[(i)]
  \item\label{item:untwisted-same-class} If $v(c) \equiv v(a) \pmod{n}$, then
    \begin{equation*}
      W^\vee(c) = (1 - q^{-1}) \int_{n v(t) \geq v(c)} \psi(t) \, d t - q^{-1} \int_{n v(t) = v(c) - 1} \psi(t) \, d t.
    \end{equation*}
  \item\label{item:untwisted-other-class} If $v(c) \not\equiv v(a) \pmod{n}$, then $W^\vee(c) = 0$ unless $v(c) = n j + 1$ for some $j \in \mathbb{Z}$, in which case, for any $t \in F^\times$ with $v(t) = j$,
    \begin{equation*}
      W^\vee(c) = G(\{c/a, \cdot\}_n, - t^n / c) \int_{v(t') = j} \psi(t') \, d t'.
    \end{equation*}
  \end{enumerate}
\end{proposition}
\begin{proof}
  With $\chi := \{a, \cdot\}_n$, the character $\rho_c := \overline{\chi} \{c, \cdot\}_n = \{c / a, \cdot\}_n$ is trivial in case \eqref{item:untwisted-same-class}, and of conductor exponent one in case \eqref{item:untwisted-other-class}.  We apply Lemma \ref{lemma:twisted-reduction}.  In the first case, we conclude by \eqref{eq:ramanujan-facts}.  In the second, we use \eqref{eq:gauss-sum-facts}, noting that $G(\rho_c, - (s t)^n / c) = G(\rho_c, - t^n / c)$ for $s \in \mathfrak{o}^\times$.
\end{proof}
The integrals of $\psi$ appearing in Proposition \ref{proposition:untwisted-vector} are elementary: $\int_{v(t) \geq j} \psi(t) \, d t$ equals $q^{-j}$ if $j \geq 0$ and vanishes otherwise, while $\int_{v(t) = j} \psi(t) \, d t$ equals $q^{-j} (1 - q^{-1})$ if $j \geq 0$, equals $-1$ if $j = -1$, and vanishes otherwise.

\begin{example}\label{example:coprimality}
  Suppose $n \geq 2$.  Taking $a = 1$ and comparing with Example \ref{example:unramified-basic-function}, we obtain, for any $t \in F^\times$ with $v(t) = -1$,
  \begin{equation*}
    (1_{\mathfrak{o}^\times})^\vee(c) = (1 - q^{-1}) W_\psi^0(c) - 1_{v(c) = 1 - n} \, G(\{c, \cdot\}_n, - t^n / c).
  \end{equation*}
  Taking instead $v(a) \equiv 1 \pmod{n}$, we obtain $(1_{\mathfrak{o}^\times} \{a, \cdot\}_n)^\vee = 1_{v(c) = 1 - n}$.
\end{example}

\begin{corollary}\label{corollary:twisted-bounds}
  Let $\chi$ be a character of $\mathfrak{o}^\times$.  Let $e$ denote the conductor exponent of $\chi^n$ if that character is ramified, and set $e := 1$ otherwise.  Then $(1_{\mathfrak{o}^\times} \chi)^\vee$ is bounded in magnitude by one, and vanishes at $c$ unless $v(c) \geq -(n-1) e$.
\end{corollary}
\begin{proof}
  This follows from Propositions \ref{proposition:twisted-vector} and \ref{proposition:untwisted-vector}, together with \eqref{eq:gauss-sum-facts} and the evaluation of the integrals of $\psi$ recorded after the latter proposition.
\end{proof}

\subsection{Voronoi-type formulas for twisted Gauss sums}\label{sec:gauss-sums}
We now write Theorem \ref{theorem:main} in classical terms.  We retain the setting of \S\ref{sec:main-formula}, and assume that $S$ is large enough that
\begin{itemize}
\item $\psi_{\mathfrak{p}}$ is unramified for each $\mathfrak{p} \notin S$, and
\item the ring $\mathcal{O}_S \subset F$ of $S$-integers is a principal ideal domain.
\end{itemize}
The quotient
\begin{equation*}
  \mathbb{N}_S :=(\mathcal{O}_S - \{0\}) / \mathcal{O}_S^\times
\end{equation*}
is then in bijection with the set of nonzero ideals of $\mathcal{O}_S$.  We write
\begin{equation*}
  \lvert x \rvert_S := \prod_{\mathfrak{p} \in S} \lvert x \rvert_{\mathfrak{p}},
  \qquad
  \psi_S := \prod_{\mathfrak{p} \in S} \psi_{\mathfrak{p}},
  \qquad
  \{x, y\}_{S} := \prod_{\mathfrak{p} \in S} \{x_{\mathfrak{p}}, y_{\mathfrak{p}}\}_{\mathfrak{p}, n}.
\end{equation*}
The character $\psi_S$ of $F_S$ is trivial on $\mathcal{O}_S$.

In Theorem \ref{theorem:main}, the summand $W(c)$ is the product of the factors $W_{\mathfrak{p}}(c)$ with $\mathfrak{p} \in S$ and of the basic functions at the places outside $S$.  The basic functions are local generating functions for Gauss sums (Lemma \ref{lemma:jacquet-basic}).  We now express their product by global Gauss sums.  We use the power residue symbol of Definition~\ref{definition:residue-symbol}.

\begin{definition}\label{definition:global-gauss-sums}\lean{NFoldCover.normalizedGauss}
  For $b, c \in \mathcal{O}_S - \{0\}$, we define
  \begin{equation*}
    g_\psi^S(b, c) := \sum_{d \in(\mathcal{O}_S / c)^\times} \qr{d}{c}_n \psi_S(b d/c),
    \qquad
    \tilde{g}_\psi^S(b, c) :=
    \frac{1}{\lvert c \rvert_S}
    \sum_{
      \substack{
        \tau \in \mathbb{N}_S  \\
        \tau^n \mid c      
      }
    }
    \lvert \tau \rvert_S^{n - 1} g_\psi^S(b, c / \tau^n).
  \end{equation*}
\end{definition}
The second sum makes sense because $g_\psi^S(b, u^n c) = g_\psi^S(b, c)$ for $u \in \mathcal{O}_S^\times$.  If $c$ is not divisible by the $n$th power of any nonunit, then $\tilde{g}_\psi^S(b, c) = g_\psi^S(b, c) / \lvert c \rvert_S$.  In general, the sum over $\tau$ corresponds to the sums over $k$ in the basic functions, which come from the vectors $\tau^n (c, d)$ of Example \ref{example:S-arithmetic}.  The precise relation is as follows:

\begin{lemma}\label{lemma:product-of-basic-functions}
  Let $b \in \mathcal{O}_S - \{0\}$, and set $\psi_b := \psi(b \, \cdot)$.  For $c \in F^\times$, the product of the basic functions $W_{\psi_{b,\mathfrak{p}}}^0(c)$ over $\mathfrak{p} \notin S$ vanishes unless $c \in \mathcal{O}_S$, in which case
  \begin{equation*}
    \prod_{\mathfrak{p} \notin S} W_{\psi_{b,\mathfrak{p}}}^0(c) = \lvert b \rvert_S^{-1/2} \, \tilde{g}_\psi^S(b, c).
  \end{equation*}
\end{lemma}
\begin{proof}
  If $c$ is nonintegral at some $\mathfrak{p} \notin S$, the basic function there vanishes by Lemma \ref{lemma:jacquet-basic}.  Suppose now that $c \in \mathcal{O}_S - \{0\}$.  For $x \in F$, we have $\psi_S(x) = \prod_{\mathfrak{p} \notin S} \psi_{\mathfrak{p}}(-x)$.  By the Chinese remainder theorem, it follows that
  \begin{equation*}
    g_\psi^S(b, c) = \prod_{\mathfrak{p} \notin S} g_{\psi_{b,\mathfrak{p}}}(c).
  \end{equation*}
  We expand the product over $\mathfrak{p}$ of the sums over $k$ of Lemma \ref{lemma:jacquet-basic}.  The tuples $(k_{\mathfrak{p}})_{\mathfrak{p} \notin S}$ correspond to the classes $\tau \in \mathbb{N}_S$, via $v_{\mathfrak{p}}(\tau) = k_{\mathfrak{p}}$.  We have $\prod_{\mathfrak{p}} q_{\mathfrak{p}}^{(n-1) k_{\mathfrak{p}}} = \lvert \tau \rvert_S^{n-1}$ and, by Lemma \ref{lemma:gauss-ramanujan}, $g_{\psi_{b,\mathfrak{p}}}(c / \varpi_{\mathfrak{p}}^{n k_{\mathfrak{p}}}) = g_{\psi_{b,\mathfrak{p}}}(c / \tau^n)$.  By the displayed identity, the product over $\mathfrak{p}$ of the latter factors is $g_\psi^S(b, c / \tau^n)$ if $\tau^n$ divides $c$, and vanishes otherwise.  Finally, $\prod_{\mathfrak{p} \notin S} \lvert c \rvert_{\mathfrak{p}} = \lvert c \rvert_S^{-1}$ and $\prod_{\mathfrak{p} \notin S} \vol(\mathfrak{o}_{\mathfrak{p}}) = \prod_{\mathfrak{p} \notin S} \lvert b \rvert_{\mathfrak{p}}^{1/2} = \lvert b \rvert_S^{-1/2}$.
\end{proof}

For $W_S \in C_c^\infty(F_S^\times)$ of the form $W_S = \prod_{\mathfrak{p} \in S} W_{\mathfrak{p}}$, with each $W_{\mathfrak{p}} \in C_c^\infty(F_{\mathfrak{p}}^\times)$, we write $W_S^\vee := \prod_{\mathfrak{p} \in S} W_{\mathfrak{p}}^\vee$.

\begin{example}\label{example:simplest}
  Applying Theorem \ref{theorem:main} with the basic functions outside $S$, which are fixed by the transform, and using Lemma \ref{lemma:product-of-basic-functions} with $b = 1$, we obtain
  \begin{multline*}
    \sum_{c \in \mathcal{O}_S - \{0\}}
    W_S(c) \, \tilde{g}_\psi^S(1, c)
    \\
    =
    \sum_{c \in \mathcal{O}_S - \{0\}}
    W_S^\vee (c) \, \tilde{g}_\psi^S(1, c)
    +
    \sum_{\beta} \rho_S(\beta)
    \prod_{\mathfrak{p} \in S}
    \int_{\beta_{\mathfrak{p}}} W_{\mathfrak{p}}(x)
    \lvert x \rvert_{\mathfrak{p}}^{1/n} \, d^\times x.
  \end{multline*}
\end{example}
For $n = 3$ over $\mathbb{Q}(\mu_3)$, the main term for the sum over $c \equiv 1 \pmod{3}$ follows from Patterson's evaluation of the residues of the Dirichlet series of cubic Gauss sums \cite[Theorem~9.1]{Patterson1977CubicTheta}, \cite{Patterson1977CubicThetaII}.  Over number fields containing $\mu_3$, Kazhdan and Patterson \cite[\S II.3, pp.~134--135]{PMIHES_1984__59__35_0} give relations among these residues as the additive frequency varies.

We now twist the sums by characters, by taking the components of $W$ at finitely many further places $\mathfrak{p}$ to be of the form $1_{\mathfrak{o}_{\mathfrak{p}}^\times} \chi$, with $\chi$ a character of the unit group $\mathfrak{o}_{\mathfrak{p}}^\times$.  The characters for which $\chi^n$ is ramified are the simplest to treat: the transform is then a product of two Gauss sums, supported on a single valuation (Proposition \ref{proposition:twisted-vector}), and the main term vanishes (Lemma \ref{lemma:twisted-vanishing}).  We consider these first, and discuss the remaining characters in Remarks \ref{remark:further-conditions} and \ref{remark:frequency-twist}.

Let $T$ be a finite set of finite places, disjoint from $S$.  For each $\mathfrak{p} \in T$, let $\chi_{\mathfrak{p}}: \mathfrak{o}_{\mathfrak{p}}^\times \rightarrow \mathbb{C}^\times$ be a character such that $\chi_{\mathfrak{p}}^n$ is ramified.  Let $e_{\mathfrak{p}}$ denote the conductor exponent of $\chi_{\mathfrak{p}}^n$.  For $c \in \mathcal{O}_S - \{0\}$, we set
\begin{equation*}
  \chi(c) :=
  \begin{cases}
    \prod_{\mathfrak{p} \in T} \chi_{\mathfrak{p}}(c) & \text{ if } c \in \mathfrak{o}_{\mathfrak{p}}^\times \text{ for all } \mathfrak{p} \in T, \\
    0 & \text{ otherwise.}
  \end{cases}
\end{equation*}
We choose a generator $f \in \mathcal{O}_S$ of the ideal $\prod_{\mathfrak{p} \in T} \mathfrak{p}^{e_{\mathfrak{p}}}$; if $T$ is empty, we take $f = 1$.

Let $b \in \mathcal{O}_S - \{0\}$ be coprime to $f$.  The transforms at the places in $T$ will contribute the constant
\begin{equation}\label{eq:epsilon}
  \varepsilon := \prod_{\mathfrak{p} \in T}
  \bigl( 1_{\mathfrak{o}_{\mathfrak{p}}^\times} \chi_{\mathfrak{p}} \bigr)^\vee \bigl( (b f^{n-1})^{-1} \bigr),
\end{equation}
where the transforms are taken with respect to the characters $\psi_{\mathfrak{p}}$.  At each $\mathfrak{p} \in T$, the argument $(b f^{n-1})^{-1}$ has valuation $-(n-1)e_{\mathfrak{p}}$, the valuation on which the local transform is supported.  Explicitly, for $\mathfrak{p} \in T$, write $G_{\mathfrak{p}}$ for the Gauss sums of \S\ref{sec:twisted-vectors}, formed with respect to $\psi_{\mathfrak{p}}$, and set $\rho_{\mathfrak{p}} := \overline{\chi_{\mathfrak{p}}} \{f, \cdot\}_{\mathfrak{p}, n}$.  Since $b$ is a unit at $\mathfrak{p}$, Proposition \ref{proposition:twisted-vector} with $t = 1/f$ gives
\begin{equation}\label{eq:epsilon-gauss}
  \varepsilon = \prod_{\mathfrak{p} \in T}
  q_{\mathfrak{p}}^{e_{\mathfrak{p}}} \,
  G_{\mathfrak{p}}(\chi_{\mathfrak{p}}^n, 1/f) \,
  G_{\mathfrak{p}}(\rho_{\mathfrak{p}}, -b/f).
\end{equation}

\begin{theorem}\label{theorem:gauss-sums}\lean{NFoldCover.twistedGaussSummation}
  Let $T$, $\chi$, $f$, $b$ and $\varepsilon$ be as above.  Let $W_S = \prod_{\mathfrak{p} \in S} W_{\mathfrak{p}}$ with each $W_{\mathfrak{p}} \in C_c^\infty(F_{\mathfrak{p}}^\times)$.  Then $\lvert \varepsilon \rvert = 1$, and
  \begin{multline*}
    \sum_{c \in \mathcal{O}_S - \{0\}}
    W_S(c) \,
    \chi(c) \,
    \tilde{g}_\psi^S(b, c)
    \\
    =
    \varepsilon
    \sum_{c \in \mathcal{O}_S - \{0\}}
    \left( W_S \{\cdot, b\}_S \right)^\vee \left( \frac{c}{b f^{n-1}} \right)
    \overline{\chi(c)}
    \qr{c}{f}_n
    \qr{f}{c}_n
    \tilde{g}_\psi^S(b, c)
    + M_b(W_S),
  \end{multline*}
  where $M_b(W_S) := 0$ if $T$ is nonempty, while if $T$ is empty,
  \begin{equation*}
    M_b(W_S) :=
    \lvert b \rvert_S^{1/2}
    \sum_{\beta} \rho_S(\beta)
    \prod_{\mathfrak{p} \in S}
    \int_{\beta_{\mathfrak{p}}} W_{\mathfrak{p}}(x)
    \lvert x \rvert_{\mathfrak{p}}^{1/n} \, d^\times x,
  \end{equation*}
  where the $\rho_S(\beta)$ are the constants of Theorem \ref{theorem:main} for the character $\psi_b := \psi(b \, \cdot)$, and the measures $d^\times x$ in this display are likewise those attached to the characters $\psi_{b, \mathfrak{p}}$, namely $\lvert b \rvert_{\mathfrak{p}}^{1/2}$ times those attached to $\psi_{\mathfrak{p}}$.
\end{theorem}

\begin{proof}
  We apply Theorem \ref{theorem:main} to the character $\psi_b$, the set $S \cup T$, and the function $W$ with the given components at $\mathfrak{p} \in S$ and with $W_{\mathfrak{p}} := 1_{\mathfrak{o}_{\mathfrak{p}}^\times} \chi_{\mathfrak{p}}$ for $\mathfrak{p} \in T$.  If $T$ is nonempty, the main term vanishes, by Lemma \ref{lemma:twisted-vanishing}.  If $T$ is empty, it equals $\lvert b \rvert_S^{-1/2} M_b(W_S)$.

  The $S$-integers that are units at each $\mathfrak{p} \in T$ are the same as the $(S \cup T)$-integers with that property, and for such $c$, the quantity $\tilde{g}_\psi^S(b, c)$ is unchanged upon replacing $S$ by $S \cup T$.  Lemma \ref{lemma:product-of-basic-functions} thus gives
  \begin{equation*}
    \sum_{c \in F^\times} W(c) = \lvert b \rvert_S^{-1/2} \sum_{c \in \mathcal{O}_S - \{0\}} W_S(c) \chi(c) \tilde{g}_\psi^S(b, c).
  \end{equation*}
  On the dual side, the transforms are taken with respect to $\psi_b$.  Since $b$ is a unit at each $\mathfrak{p} \in T$, Lemma \ref{lemma:change-of-character} and Proposition \ref{proposition:basic-function-dual} give
  \begin{equation*}
    W_{\mathfrak{p}}^{\vee,\psi_b}(c) =
    \begin{cases}
      (W_{\mathfrak{p}} \{\cdot, b\}_n)^{\vee,\psi}(b^{n-1} c)
      & \text{ if } \mathfrak{p} \in S, \\
      W_{\mathfrak{p}}^{\vee,\psi}(b^{n-1} c)
      & \text{ if } \mathfrak{p} \in T, \\
      W_{\psi_{b,\mathfrak{p}}}^0(b^n c)
      & \text{ if } \mathfrak{p} \notin S \cup T.
    \end{cases}
  \end{equation*}
  After the substitution $c \mapsto c / b^n$, we obtain
  \begin{equation}\label{eq:dual-side}
    \sum_{c \in F^\times} W^\vee(c) = \sum_{c \in F^\times}
    \left( W_S \{\cdot, b\}_S \right)^\vee (c / b)
    \prod_{\mathfrak{p} \in T} W_{\mathfrak{p}}^{\vee} (c / b)
    \prod_{\mathfrak{p} \notin S \cup T} W_{\psi_{b,\mathfrak{p}}}^0(c).
  \end{equation}
  If $T$ is empty, we conclude by Lemma \ref{lemma:product-of-basic-functions}.  Suppose $T$ is nonempty.  By Proposition \ref{proposition:twisted-vector}, the summand in \eqref{eq:dual-side} vanishes unless $c = c_1 / f^{n-1}$ for some $c_1 \in \mathcal{O}_S - \{0\}$ that is a unit at each $\mathfrak{p} \in T$.  For such $c$, we evaluate the two products.
  \begin{itemize}
  \item Let $\mathfrak{p} \in T$.  Set $c_0 := (b f^{n-1})^{-1}$, so that $c / b = c_0 c_1$.  Since $b$ is a unit at $\mathfrak{p}$, we have $\{c_0, u\}_{\mathfrak{p}, n} = \{f, u\}_{\mathfrak{p}, n}$ for $u \in \mathfrak{o}_{\mathfrak{p}}^\times$.  Lemma \ref{lemma:twisted-reduction} thus gives $W_{\mathfrak{p}}^\vee(c_0 c_1) = \rho_{\mathfrak{p}}(c_1) W_{\mathfrak{p}}^\vee(c_0)$, and $W_{\mathfrak{p}}^\vee(c_0)$ is the $\mathfrak{p}$-factor of $\varepsilon$ in \eqref{eq:epsilon}.  Since $\rho_{\mathfrak{p}}(c_1) = \overline{\chi_{\mathfrak{p}}(c_1)} \{f, c_1\}_{\mathfrak{p},n}$, the product over $\mathfrak{p} \in T$ is $\varepsilon \, \overline{\chi(c_1)} \qr{c_1}{f}_n$.
  \item Let $\mathfrak{p} \notin S \cup T$.  By Lemma \ref{lemma:gauss-ramanujan}, we have $W_{\psi_{b,\mathfrak{p}}}^0(c) = \{c_1, f\}_{\mathfrak{p}, n} W_{\psi_{b,\mathfrak{p}}}^0(c_1)$.  The symbol is trivial unless $\mathfrak{p}$ divides $c_1$, so by Lemma \ref{lemma:product-of-basic-functions}, the product over such $\mathfrak{p}$ is $\lvert b \rvert_S^{-1/2} \qr{f}{c_1}_n \tilde{g}_\psi^S(b, c_1)$.
  \end{itemize}
  The theorem follows.  The modulus of $\varepsilon$ is one by \eqref{eq:epsilon-gauss} and \eqref{eq:gauss-sum-facts}.
\end{proof}

\begin{example}
  If $n = 1$, then $\tilde{g}_\psi^S(b, c) = 1_{c \mid b}$ and $\varepsilon = \chi(b)$, and the theorem reads
  \begin{equation*}
    \sum_{c \mid b} W_S(c) \chi(c) = \sum_{c \mid b} W_S(b/c) \chi(b/c).
  \end{equation*}
\end{example}

\begin{remark}\label{remark:further-conditions}
  Theorem \ref{theorem:gauss-sums} requires that $\chi_{\mathfrak{p}}^n$ be ramified for each $\mathfrak{p} \in T$.  The remaining characters (including the trivial character, which imposes the condition that $c$ be prime to $\mathfrak{p}$) are covered by the same theorem, because the set $S$ may be enlarged at will: for an arbitrary character $\chi_{\mathfrak{p}}$ of $\mathfrak{o}_{\mathfrak{p}}^\times$ at a place $\mathfrak{p} \notin S \cup T$, we adjoin $\mathfrak{p}$ to $S$ and take $W_{\mathfrak{p}} := 1_{\mathfrak{o}_{\mathfrak{p}}^\times} \chi_{\mathfrak{p}}$.  The transforms of such functions are computed in \S\ref{sec:twisted-vectors}, and bounded, uniformly in $\mathfrak{p}$ and $\chi_{\mathfrak{p}}$, in Corollary \ref{corollary:twisted-bounds}.
\end{remark}

\begin{remark}\label{remark:frequency-twist}
  Twists by characters $\chi_{\mathfrak{p}}$ with $\chi_{\mathfrak{p}}^n = 1$ may also be expressed through the frequency $b$.  Let $\mathfrak{p} \notin S$ be a prime with generator $\pi \in \mathcal{O}_S$, and set $\chi_\pi := \{\pi, \cdot\}_{\mathfrak{p}, n}$.  By property \eqref{item:hilbert-tame} of the Hilbert symbol, the characters $\chi$ of $\mathfrak{o}_{\mathfrak{p}}^\times$ with $\chi^n = 1$ are the powers of $\chi_\pi$.  Let $c \in \mathcal{O}_S - \{0\}$ be prime to $\pi$.  For any nonzero $t \in \mathcal{O}_S$ prime to $c$, the substitution $d \mapsto t^{-1} d$ in each Gauss sum of Definition \ref{definition:global-gauss-sums} gives
  \begin{equation}\label{eq:frequency-twist}
    \tilde{g}_\psi^S(t b, c) = \qr{t}{c}_n^{-1} \tilde{g}_\psi^S(b, c),
  \end{equation}
  since $\qr{t}{c / \tau^n}_n = \qr{t}{c}_n$ for each $\tau$ with $\tau^n \mid c$.  Taking $t = \pi^{n-a}$ and applying the reciprocity law in the form $\qr{\pi}{c}_n = \chi_\pi(c) \{\pi, c\}_S$, we obtain, for $0 \leq a \leq n$,
  \begin{equation*}
    \chi_\pi(c)^a \, \tilde{g}_\psi^S(b, c) = \{\pi, c\}_S^{-a} \, \tilde{g}_\psi^S(\pi^{n-a} b, c).
  \end{equation*}
  Thus such a twist amounts to a change of frequency, together with the condition that $c$ be prime to $\pi$ and a factor depending only upon the image of $c$ in $F_S^\times$.
\end{remark}

\begin{remark}\label{remark:sawin-whitehead}
  The coefficients considered by Sawin and Whitehead \cite[Proposition~2.6]{2025arXiv2507.08662} are, up to a factor depending only upon $\deg c$, of the form $\chi(c) \, \tilde{g}_\psi^S(b, c)$, with $F = \mathbb{F}_q(T)$ and $S = \{\infty\}$, where $\chi$ is a product of residue characters modulo primes $\pi$ at which $\chi_{\mathfrak{p}}^n$ is ramified.  Such a prime may also divide the frequency $b$.  Since the coefficient vanishes unless $c$ is prime to $\pi$, the identity \eqref{eq:frequency-twist} then allows us to replace $b$ by its prime-to-$\pi$ part, at the cost of multiplying $\chi_{\mathfrak{p}}$ by a character of order dividing $n$ and the weight at infinity by a Hilbert symbol.  Theorem \ref{theorem:gauss-sums} thus applies.
\end{remark}

\subsection{Twists with quadratic duals}\label{sec:quadratic-dual}
For cubic Gauss sums, a sextic twist can become quadratic in the dual sum.  Two Hilbert symbols account for the change.  We first explain this at one prime dividing the modulus, for Gauss sums of any order $n$.

We retain the setting of \S\ref{sec:gauss-sums}.  We say that $h \in \mathcal{O}_S - \{0\}$ is \emph{squarefree} if $v_{\mathfrak{p}}(h) \leq 1$ for all $\mathfrak{p} \notin S$.

\subsubsection{The dual twist}\label{sec:two-symbols}
Let $h \in \mathcal{O}_S - \{0\}$ be squarefree and not a unit, and let $\mathfrak{p} \notin S$ be a prime dividing $h$.  We work at $\mathfrak{p}$, and omit the subscript $\mathfrak{p}$ from the local notation.  Let $\chi$ be a character of $\mathfrak{o}^\times$ such that $\chi^n$ has conductor exponent one, and set $W := 1_{\mathfrak{o}^\times} \chi$.

The first symbol comes from the transform on the line.  By Lemma \ref{lemma:twisted-reduction}, for $u \in \mathfrak{o}^\times$ and $c \in F_{\mathfrak{p}}^\times$,
\begin{equation*}
  W^\vee(u c) = \chi(u)^{-1} \{c, u\}_n W^\vee(c).
\end{equation*}
Indeed, in the integral defining $W^\vee(u c)$, the substitution $x \mapsto x / u$ replaces $W(x)$ by $\chi(u)^{-1} W(x)$ and $\{x, - u c\}_n$ by $\{c, u\}_n \{x, -c\}_n$.  By Proposition \ref{proposition:twisted-vector}, the function $W^\vee$ is supported on the elements of valuation $1 - n$.  With
\begin{equation*}
  a := h^{1-n},
\end{equation*}
these are the products $a u$ with $u \in \mathfrak{o}^\times$, and
\begin{equation}\label{eq:first-symbol}
  W^\vee(a u) = W^\vee(a) \, \chi(u)^{-1} \{a, u\}_n.
\end{equation}

The second symbol comes from the basic functions at the remaining primes, when we index the dual sum by integers.  Let $T$ denote the set of primes dividing $h$.  Consider the dual sum $\sum_{c \in F^\times} \prod_{\mathfrak{q}} W_{\mathfrak{q}}^\vee(c)$ of Theorem \ref{theorem:main}, with a component as above at each prime in $T$, and the basic function at each remaining prime outside $S$.  Its summand vanishes unless $c = a c_1$, where $c_1 \in \mathcal{O}_S - \{0\}$ is a unit at each prime in $T$.  Indeed, at those primes, this is the support condition above; at the remaining primes outside $S$, the element $a$ is a unit, and the basic function vanishes outside the ring of integers.

We substitute $c = a c_1$.  At a prime $\mathfrak{q} \notin S \cup T$, where $a$ is a unit, Lemma \ref{lemma:gauss-ramanujan} gives
\begin{equation*}
  W_{\psi_{\mathfrak{q}}}^0(a c_1) = \{c_1, a\}_{\mathfrak{q}, n} \, W_{\psi_{\mathfrak{q}}}^0(c_1).
\end{equation*}
By the reciprocity law, the product of these symbols is
\begin{equation}\label{eq:second-symbol}
  \prod_{\mathfrak{q} \notin S \cup T} \{c_1, a\}_{\mathfrak{q}, n}
  = \{a, c_1\}_S \prod_{\mathfrak{q} \in T} \{a, c_1\}_{\mathfrak{q}, n}.
\end{equation}
The reindexed summand is thus the product of the basic functions $W_{\psi_{\mathfrak{q}}}^0(c_1)$ at the primes $\mathfrak{q} \notin S \cup T$, of the factor $\{a, c_1\}_S W_S^\vee(a c_1)$ from the places in $S$, and of a weight at each prime in $T$.  At $\mathfrak{p}$, where $c_1$ is a unit $u$, that weight is, by \eqref{eq:first-symbol},
\begin{equation}\label{eq:two-symbols}
  \{a, u\}_n W^\vee(a u)
  = W^\vee(a) \, \chi(u)^{-1} \{a, u\}_n^2.
\end{equation}
One factor $\{a, u\}_n$ comes from the Hilbert symbol in the transform at $\mathfrak{p}$.  The other is the factor at $\mathfrak{p}$ on the right-hand side of \eqref{eq:second-symbol}: it comes from the basic functions at the primes not dividing $h$, and is carried to $\mathfrak{p}$ by the reciprocity law.  We have $\{a, u\}_n = \{h, u\}_n$, since $a h^{-1}$ is an $n$th power.  The dual sum is thus twisted at $\mathfrak{p}$ by the character $\chi^{-1} \kappa^2$ of $\mathfrak{o}^\times$, where $\kappa := \{h, \cdot\}_n$.

Suppose now that $n$ is odd, and that the residue characteristic of each place outside $S$ is odd: if $F$ is a number field, this means that $S$ contains the places above $2$, and if $F$ is a function field, it means that $F$ has odd characteristic.  Let $\eta := \{h, \cdot\}_2$ denote the quadratic character of $\mathfrak{o}^\times$, and take
\begin{equation*}
  \chi := \eta \kappa^2.
\end{equation*}
Then $\chi^n = \eta$ has conductor exponent one, so the above applies, and $\chi^{-1} \kappa^2 = \eta$: the dual twist is quadratic.  For $n = 3$, the character $\chi$ is sextic, and the function $u \mapsto W^\vee(a u)$ itself still has the sextic character $\chi^{-1} \kappa = \chi$; it is the reindexed weight \eqref{eq:two-symbols} that has the quadratic character $\eta$.

Proposition \ref{proposition:twisted-vector} evaluates the constant $W^\vee(a)$.  The product of these constants over the primes dividing $h$ is the constant $\varepsilon$ of \eqref{eq:epsilon}, with $b = 1$ and $f = h$.  In Theorem \ref{theorem:gauss-sums}, the two symbols appear as the residue symbols $\qr{c}{f}_n$ and $\qr{f}{c}_n$, which agree, by the reciprocity law, up to a Hilbert symbol at the places in $S$.

\begin{remark}\label{remark:lifted-dilations}
  On the plane, the square in \eqref{eq:two-symbols} is a commutator.  For $a \in F_{\mathfrak{p}}^\times$, let $d_1(a)$ and $d_2(a)$ be the lifts of the diagonal matrices considered in the proof of Lemma \ref{lemma:torus}, and set
  \begin{equation*}
    D_a := \lvert a \rvert \, d_2(a) \, d_1(a).
  \end{equation*}
  By the transition relation \eqref{eq:transition-relation}, as in that proof, $d_2(a)$ has the multiplier $\{a, x\}_n$ in the $X$-chart, and $d_1(a)$ has the multiplier $\{a, y\}_n$ in the $Y$-chart.  Since $\{a, a y\}_n = \{a, -y\}_n$, we obtain, for each genuine function $\Phi$,
  \begin{align*}
    (D_a \Phi)_X(x, y) &= \lvert a \rvert \{a, x\}_n \Phi_X(a x, a y), \\
    (D_a \Phi)_Y(x, y) &= \lvert a \rvert \{a, -y\}_n \Phi_Y(a x, a y).
  \end{align*}
  The operator $D_a$ lifts the dilation by $a$.  It preserves $\mathcal{D}$, and the substitution $x \mapsto x / a$ in Definition \ref{definition:lines} gives $\mathcal{F} D_a = D_{a^{-1}} \mathcal{F}$ on $\mathcal{D}$; thus $D_a$ also preserves $\mathcal{S}$.  The factor $\lvert a \rvert$ cancels the change of measure in the Jacquet integral: for $\Phi \in \mathcal{S}$,
  \begin{equation*}
    J(D_a \Phi)(c) = \{a, c\}_n J \Phi(a c),
  \end{equation*}
  which is the local weight considered above.  For $u \in F_{\mathfrak{p}}^\times$, comparing $X$-components gives $D_u D_a = \{a, u\}_n D_{a u}$ and likewise $D_a D_u = \{u, a\}_n D_{a u}$, so
  \begin{equation*}
    D_u D_a = \{a, u\}_n^2 \, D_a D_u.
  \end{equation*}
  With $W$ and $a$ as above, the element $\Phi$ of $\mathcal{D}$ with $\Phi_X(x, y) := W(x) 1_{\mathfrak{o}}(y)$ satisfies $J \Phi = W$ and $D_u \Phi = \chi(u) \Phi$ for $u \in \mathfrak{o}^\times$.  Thus $D_u \mathcal{F} \Phi = \chi(u)^{-1} \mathcal{F} \Phi$, and
  \begin{equation*}
    D_u (D_a \mathcal{F} \Phi) = \chi(u)^{-1} \{a, u\}_n^2 \, D_a \mathcal{F} \Phi.
  \end{equation*}
  Since $J(\mathcal{F} \Phi) = W^\vee$ by definition, taking Jacquet integrals at $c = 1$ recovers \eqref{eq:two-symbols}.
\end{remark}

\subsubsection{The summation formula}\label{sec:quadratic-dual-summation-formula}
We assume, as at the end of \S\ref{sec:two-symbols}, that $n$ is odd and that the residue characteristic of each place outside $S$ is odd.

\begin{corollary}\label{corollary:quadratic-dual}\lean{NFoldCover.quadraticDual}
  Let $h \in \mathcal{O}_S - \{0\}$ be squarefree and not a unit.  There is a complex number $\varepsilon$ of modulus one such that, for all $W_S = \prod_{\mathfrak{p} \in S} W_{\mathfrak{p}}$ with each $W_{\mathfrak{p}} \in C_c^\infty(F_{\mathfrak{p}}^\times)$, we have
  \begin{multline*}
    \sum_{c \in \mathcal{O}_S - \{0\}}
    W_S(c)
    \qr{c}{h}_2
    \qr{c}{h}_n^2 \,
    \tilde{g}_\psi^S(1, c)
    \\
    =
    \varepsilon
    \sum_{c \in \mathcal{O}_S - \{0\}}
    W_S^\vee \left( \frac{c}{h^{n-1}} \right)
    \{h, c\}_{S}
    \qr{c}{h}_2
    \tilde{g}_\psi^S(1, c).
  \end{multline*}
\end{corollary}
\begin{example}
  If $n = 3$, then $\qr{c}{h}_2 \qr{c}{h}_3^2 = \qr{c}{h}_6$ is the sextic residue symbol, defined with respect to an embedding $\mu_6 \hookrightarrow \mathbb{C}^\times$ extending the given one on $\mu_3$.\lean{MetaplecticPoisson.Datum.sextic_to_quadratic}
\end{example}
\begin{proof}[Proof of Corollary \ref{corollary:quadratic-dual}]
  We apply Theorem \ref{theorem:gauss-sums} with $b = 1$, with $T$ the set of primes dividing $h$, and with the characters $\chi_{\mathfrak{p}} := \{h, \cdot\}_{\mathfrak{p}, 2} \{h, \cdot\}_{\mathfrak{p}, n}^2$ of \S\ref{sec:two-symbols}, for which $\chi_{\mathfrak{p}}^n$ is the quadratic character of $\mathfrak{o}_{\mathfrak{p}}^\times$, of conductor exponent one.  Thus $f = h$ and $\chi(c) = \qr{c}{h}_2 \qr{c}{h}_n^2$, and the constant $\varepsilon$ of that theorem does not depend on $W_S$.  The reciprocity law gives $\qr{h}{c}_n = \qr{c}{h}_n \{h, c\}_S$, so the dual sum involves
  \begin{equation*}
    \overline{\chi(c)} \qr{c}{h}_n \qr{h}{c}_n = \qr{c}{h}_2 \{h, c\}_S.
    \qedhere
  \end{equation*}
\end{proof}

\appendix
\section{Proof of Fourier inversion}\label{sec:marginals}
We prove the inversion formula on $\mathcal{D}$ (Proposition \ref{proposition:inversion}), from which we deduced Theorem \ref{theorem:schwartz-space}.  We work in the local setting of \S\S\ref{sec:hilbert-symbol}--\ref{sec:germs}.

Let $\Phi \in \mathcal{S}$.  For $x, y \in F^\times$, we set
\begin{equation*}
  m_\Phi(x) := \int_F \Phi_X(x, y) \, d y,
  \qquad
  r_\Phi(y) := \Phi_Y(0, y).
\end{equation*}
We call $m_\Phi$ the \emph{marginal} of $\Phi$, and $r_\Phi$ its \emph{axis values}.  For $n = 1$, the Fourier transform on the plane exchanges integration along the lines $x = \mathrm{const}$ with restriction to the line $x = 0$:
\begin{equation*}
  r_{\mathcal{F} \Phi}(y') = \int_F m_\Phi(x) \psi(x y') \, d x,
  \qquad
  m_{\mathcal{F} \Phi}(x') = \int_F r_\Phi(y) \psi(- y x') \, d y.
\end{equation*}
The first identity is immediate from the definition of the symplectic Fourier transform, and the second follows from the Fourier inversion formula on $F$.  We establish the corresponding relations for general $n$ (Lemmas \ref{lemma:marginal-axis} and \ref{lemma:axis-marginal}), in which the Fourier transforms are taken in an $n$th root variable.  The first is again immediate from the definition of $\mathcal{F}$, while the second requires an argument.  Together, they give the inversion formula, and they are the local input to the proof of the Poisson formula in Appendix \ref{sec:proof-poisson}.

\subsection{From marginals to axis values}\label{sec:from-marginals-to-axis-values}
Let $\Phi$ be a genuine function satisfying \eqref{eq:growth}, and let $y', \tau \in F^\times$.  For $v' = (0, y')$, the line in Definition \ref{definition:lines} is given by $x = \tau^n / y'$.  For such $x$, we have $\Phi_Y(x, y) = \{y, x\}_n \Phi_X(x, y)$ and $\{y, x\}_n = \{y, y'\}_n^{-1}$, so
\begin{equation}\label{eq:axis-line-integral}
  R_Y \Phi((0, y'), \tau) = \frac{\{-1, y'\}_n}{\lvert y' \rvert} \int_F \Phi_X(\tau^n / y', y) \, d y.
\end{equation}

\begin{lemma}\label{lemma:marginal-axis}
  Let $\Phi \in \mathcal{S}$ and $y' \in F^\times$.  For $\tau \in F$, set
  \begin{equation*}
    \nu(\tau) := R_Y \Phi((0, y'), \tau).
  \end{equation*}
  Then $\nu$ is continuous, with
  \begin{align*}
    \nu(\tau) &= \{-1, y'\}_n \lvert y' \rvert^{-1} m_\Phi(\tau^n / y')
    \qquad (\tau \neq 0), \\
    \nu(0) &= \{-1, y'\}_n \lvert y' \rvert^{-1} \int_F \{y, y'\}_n r_\Phi(y) \, d y,
  \end{align*}
  and for $\mu \in F^\times$,
  \begin{equation*}
    r_{\mathcal{F} \Phi}(\mu^n y') = \lvert n \rvert^{-1/2} \lvert \mu \rvert^{1 - n} \int_F \nu(\tau) \psi(\mu \tau) \, d \tau.
  \end{equation*}
  If $\Phi \in \mathcal{D}$, then $\nu$ is a Schwartz--Bruhat function on $F$.
\end{lemma}
\begin{proof}
  The continuity is part of Lemma \ref{lemma:lines}\eqref{item:lines-convergence}.  The formula for $\tau \neq 0$ is \eqref{eq:axis-line-integral}, and that for $\tau = 0$ is Definition \ref{definition:lines}, since $\Phi_Y(0, y) = r_\Phi(y)$.  The formula for $r_{\mathcal{F} \Phi}$ is Lemma \ref{lemma:transform-fourier}.  The last assertion is part of Lemma \ref{lemma:lines}\eqref{item:lines-smooth}.
\end{proof}

\subsection{From axis values to marginals}\label{sec:from-axis-values-to-marginals}
The converse relation is the main local computation.  It reduces, by the Steinberg relation, to the Fourier inversion formula on $F$.

\begin{lemma}\label{lemma:axis-marginal}
  Let $\Phi \in \mathcal{D}$ and $x' \in F^\times$.  Define a Schwartz--Bruhat function $\sigma$ on $F$ by
  \begin{equation*}
    \sigma(\tau) := \lvert n \rvert \lvert x' \rvert^{-1} \lvert \tau \rvert^{n-1} r_\Phi(- \tau^n / x'),
    \qquad \sigma(0) := 0.
  \end{equation*}
  Then for $\mu \in F^\times$,
  \begin{equation*}
    m_{\mathcal{F} \Phi}(\mu^n x') = \lvert n \rvert^{-1/2} \int_F \sigma(\tau) \psi(\mu \tau) \, d \tau.
  \end{equation*}
\end{lemma}
\begin{proof}
  Write $s := x'$.  The substitution $y = - a^n / s$, for which $d y = \lvert n \rvert \lvert s \rvert^{-1} \lvert a \rvert^{n-1} \, d a$, shows that the assertion is equivalent to the identity
  \begin{equation}\label{eq:axis-marginal-roots}
    \lvert n \rvert^{1/2} \, m_{\mathcal{F} \Phi}(\mu^n s) = \int_F r_\Phi(y) \sum_{a^n = - s y} \psi(\mu a) \, d y,
  \end{equation}
  where $a$ runs over the $n$th roots of $- s y$ in $F$; there are $n$ of them if $- s y \in F^{\times n}$, and none otherwise.  Both sides of \eqref{eq:axis-marginal-roots} are continuous functions of $\mu \in F^\times$.  It thus suffices to show that they have the same integral against $\rho(\mu) \, d^\times \mu$ for each $\rho \in C_c^\infty(F^\times)$.  Fix such a $\rho$, and set $h := \hat{\rho}$.  Then $h$ is a Schwartz--Bruhat function on $F$, and $\int_F h = \rho(0) = 0$.

  \emph{The $Y$-chart.}  Let $\mu \in F^\times$.  By Lemma \ref{lemma:transform-fourier}, the substitution $y \mapsto \mu^n y'$ in the definition of the marginal gives
  \begin{equation*}
    \lvert n \rvert^{1/2} \, m_{\mathcal{F} \Phi}(\mu^n s) = \lvert \mu \rvert \int_F \left( \int_F R_X \Phi((s, y'), \tau) \psi(\mu \tau) \, d \tau \right) \, d y'.
  \end{equation*}
  The inner integral is $\lvert n \rvert^{1/2} \lvert \mu \rvert^{n-1} (\mathcal{F} \Phi)_X(\mu^n s, \mu^n y')$.  For $\mu$ in the support of $\rho$, it is $O((1 + \lvert y' \rvert)^{-2})$, by Corollary \ref{corollary:transform-decay}.  This bound allows us to integrate over $\mu$ before $y'$.  For each $y'$, the function $R_X \Phi((s, y'), \cdot)$ is integrable by Lemma \ref{lemma:lines}\eqref{item:lines-convergence}, so we may then integrate over $\mu$ before $\tau$.  This gives
  \begin{equation}\label{eq:axis-marginal-tested}
    \lvert n \rvert^{1/2} \int_{F^\times} m_{\mathcal{F} \Phi}(\mu^n s) \rho(\mu) \, d^\times \mu
    = \int_F \left( \int_F R_X \Phi((s, y'), \tau) h(\tau) \, d \tau \right) \, d y',
  \end{equation}
  where the integral over $y'$ converges absolutely.  Let $y', \tau \in F^\times$.  By Lemma \ref{lemma:lines}\eqref{item:lines-compatibility} and Definition \ref{definition:lines}, since $\{s, y'\}_n \{y, y'\}_n \{-1, y'\}_n = \{- s y, y'\}_n$, we have
  \begin{equation*}
    R_X \Phi((s, y'), \tau) = \int_F \Phi_Y \left( \frac{\tau^n + s y}{y'}, y \right) \{- s y, y'\}_n \, \frac{d y}{\lvert y' \rvert}.
  \end{equation*}
  For $y \in F^\times$, we write
  \begin{equation*}
    f_y(x) := \Phi_Y(x, y),
    \qquad
    \eta_y(v) := \{v, - s y\}_n.
  \end{equation*}
  Thus $f_y$ is a Schwartz--Bruhat function on $F$, with $f_y(0) = r_\Phi(y)$, and $\eta_y$ is a character of $F^\times$.  With $v := 1 / y'$, we have $\{- s y, y'\}_n = \eta_y(v)$ and $d y' / \lvert y' \rvert = d^\times v$.  Let $V$ be a compact subset of $F^\times$.  Since $\int_F h = 0$, the contribution to the right-hand side of \eqref{eq:axis-marginal-tested} of the $y'$ with $1 / y' \in V$ is
  \begin{equation}\label{eq:axis-marginal-truncated}
    \int_{F^\times} \int_F h(\tau) \int_V \eta_y(v) \left( f_y((\tau^n + s y) v) - f_y(v) \right) \, d^\times v \, d \tau \, d y.
  \end{equation}
  These integrals converge absolutely, because $V$ is compact.

  \emph{Convergence.}  We claim that there is a bounded function $C$ on $F^\times$, decaying rapidly at infinity, such that for $y, b \in F^\times$,
  \begin{equation}\label{eq:log-bound}
    \int_{F^\times} \lvert f_y(b v) - f_y(v) \rvert \, d^\times v \leq C(y) (1 + \lvert \log \lvert b \rvert \rvert).
  \end{equation}
  To see this, let $\chi$ denote the characteristic function of $\{v : \lvert v \rvert \leq 1\}$, and write $f_y = r_\Phi(y) \chi + g_y$.  The integral of $\lvert \chi(b v) - \chi(v) \rvert$ is $O(\lvert \log \lvert b \rvert \rvert)$, and that of $\lvert g_y(b v) - g_y(v) \rvert$ is at most twice the integral of $\lvert g_y \rvert$.  By the definition of $\mathcal{D}$, the integral of $\lvert g_y(v) \rvert$ with respect to $d^\times v$ is bounded, and decays rapidly as $\lvert y \rvert \rightarrow \infty$: for $v$ small relative to $y$, we use the local constancy of $\Phi_Y$ or the mean value theorem, according as $F$ is non-archimedean or archimedean, and for the remaining $v$, we use the decay of $\Phi$ at the origin and at infinity.

  The function $(\tau, y) \mapsto h(\tau) C(y) (1 + \lvert \log \lvert \tau^n + s y \rvert \rvert)$ is integrable on $F \times F^\times$, because the logarithmic singularity is integrable in $y$.  By \eqref{eq:log-bound}, we may thus let $V$ increase to $F^\times$ in \eqref{eq:axis-marginal-truncated}.  We obtain
  \begin{equation}\label{eq:axis-marginal-L}
    \lvert n \rvert^{1/2} \int_{F^\times} m_{\mathcal{F} \Phi}(\mu^n s) \rho(\mu) \, d^\times \mu
    = \int_{F^\times} \left( \int_F h(\tau) L_y(\tau^n + s y) \, d \tau \right) \, d y,
  \end{equation}
  where, for $b \in F^\times$,
  \begin{equation*}
    L_y(b) := \int_{F^\times} \eta_y(v) \left( f_y(b v) - f_y(v) \right) \, d^\times v.
  \end{equation*}

  \emph{The symbol.}  Fix $y \in F^\times$, and write $L := L_y$, $\eta := \eta_y$ and $f := f_y$.  For $b, c \in F^\times$, the substitution $v \mapsto v / c$ gives
  \begin{equation}\label{eq:log-cocycle}
    L(b c) = \eta(c)^{-1} L(b) + L(c).
  \end{equation}
  Exchanging $b$ and $c$, we see that $(\eta(c)^{-1} - 1) L(b) = (\eta(b)^{-1} - 1) L(c)$.  Suppose that $- s y \notin F^{\times n}$.  Then $\eta$ is nontrivial, by property \eqref{item:hilbert-powers} of the Hilbert symbol, and so $L(b) = 0$ whenever $\eta(b) = 1$.  By property \eqref{item:hilbert-steinberg} of the Hilbert symbol, applied after dividing both entries by $\tau^n$ when $\tau \neq 0$, we have $\eta(\tau^n + s y) = \{\tau^n + s y, - s y\}_n = 1$ for all $\tau \in F$.  The inner integral in \eqref{eq:axis-marginal-L} thus vanishes.

  \emph{The roots.}  Suppose now that $- s y \in F^{\times n}$.  Then $\eta$ is trivial, and \eqref{eq:log-cocycle} reads $L(b c) = L(b) + L(c)$.  Since $\tau^n + s y$ is the product of the factors $\tau - a$ over the roots $a$ of $- s y$, we have $L(\tau^n + s y) = \sum_{a^n = - s y} L(\tau - a)$.  For each such $a$, we claim that
  \begin{equation}\label{eq:linear-factor}
    \int_F h(\tau) L(\tau - a) \, d \tau = f(0) \int_{F^\times} \rho(\mu) \psi(\mu a) \, d^\times \mu.
  \end{equation}
  This is the case $n = 1$ of the lemma.  By \eqref{eq:log-bound}, the double integral on the left-hand side converges absolutely, so we may integrate over $\tau$ first.  Since $\int_F h = 0$, the term $f(v)$ then drops out.  By the Parseval formula, the left-hand side is thus
  \begin{equation*}
    \int_{F^\times} \left( \int_F h(\tau) f((\tau - a) v) \, d \tau \right) \, d^\times v
    =
    \int_{F^\times} \left( \int_F \rho(\mu) \psi(\mu a) \lvert v \rvert^{-1} \hat{f}(\mu / v) \, d \mu \right) \, d^\times v.
  \end{equation*}
  The double integral on the right converges absolutely, and $\int_{F^\times} \lvert v \rvert^{-1} \hat{f}(\mu / v) \, d^\times v = \lvert \mu \rvert^{-1} \int_F \hat{f} = \lvert \mu \rvert^{-1} f(0)$.  This gives \eqref{eq:linear-factor}.

  \emph{Conclusion.}  Since $f_y(0) = r_\Phi(y)$, we have shown that the right-hand side of \eqref{eq:axis-marginal-L} is
  \begin{equation*}
    \int_{F^\times} r_\Phi(y) \sum_{a^n = - s y} \left( \int_{F^\times} \rho(\mu) \psi(\mu a) \, d^\times \mu \right) \, d y,
  \end{equation*}
  which is the integral of the right-hand side of \eqref{eq:axis-marginal-roots} against $\rho(\mu) \, d^\times \mu$.  By \eqref{eq:axis-marginal-L}, the same holds for the left-hand side.
\end{proof}

\subsection{The inversion formula}\label{sec:inversion-formula}
\begin{proposition}\label{proposition:inversion}
  Set $\varepsilon := (-1)^{n+1}$.  For $\Phi \in \mathcal{D}$, we have $\mathcal{F} \mathcal{F} \Phi = h(\varepsilon) \Phi$.
\end{proposition}
The idea is as follows.  By Lemmas \ref{lemma:marginal-axis} and \ref{lemma:axis-marginal}, the transform exchanges marginals and axis values, each time by a Fourier transform in an $n$th root variable.  Thus $\mathcal{F}^2$ acts on axis values by Fourier inversion; the two factors $\lvert n \rvert^{-1/2}$ cancel the Jacobian $\lvert n \rvert$ of the $n$th power map, which enters through $\sigma$.  This gives the formula on the axis $x = 0$, hence everywhere, since the translations $n^-(u)$ commute with $\mathcal{F}$ and move each point of $U_Y$ onto that axis.

\begin{proof}
  Let $\Phi \in \mathcal{D}$, and set $\Psi := \mathcal{F} \Phi \in \mathcal{S}$.

  We first compute the axis values of $\mathcal{F} \Psi$.  Let $y' \in F^\times$, and set $f(t) := \lvert t \rvert^{n-1} r_\Phi(- t^n y')$, so that $f$ is a Schwartz--Bruhat function on $F$.  By Lemma \ref{lemma:axis-marginal} with $x' := 1 / y'$, we have $m_\Psi(\tau^n / y') = \lvert n \rvert^{1/2} \lvert y' \rvert \hat{f}(\tau)$ for $\tau \neq 0$.  By Lemma \ref{lemma:marginal-axis}, applied to $\Psi$ with $\mu = 1$, and by Fourier inversion,
  \begin{equation*}
    r_{\mathcal{F} \Psi}(y')
    = \{-1, y'\}_n \int_F \hat{f}(\tau) \psi(\tau) \, d \tau
    = \{-1, y'\}_n f(-1)
    = \{-1, y'\}_n r_\Phi(\varepsilon y').
  \end{equation*}
  For odd $n$, we have $\{-1, y'\}_n = 1$, since $-1 = (-1)^n$.  In either case, the right-hand side is thus $r_{h(\varepsilon) \Phi}(y')$.

  Let $u \in F$.  By Corollary \ref{corollary:metaplectic}, we have $n^-(u) \Phi \in \mathcal{D}$ and $\mathcal{F} n^-(u) \Phi = n^-(u) \Psi$.  By Lemma \ref{lemma:unipotent-properties} for $\Psi$, and the above applied to $n^-(u) \Phi$, we obtain
  \begin{equation*}
    (\mathcal{F} \Psi)_Y(u y', y')
    = r_{\mathcal{F} n^-(u) \Psi}(y')
    = r_{h(\varepsilon) n^-(u) \Phi}(y')
    = (h(\varepsilon) \Phi)_Y(u y', y').
  \end{equation*}
  In the last step, we used that $h(\varepsilon)$ commutes with $n^-(u)$, by Lemma \ref{lemma:torus}, as $\varepsilon^2 = 1$.  Thus $\mathcal{F} \Psi$ and $h(\varepsilon) \Phi$ have the same $Y$-component.  Their $X$-components are continuous on $U_X$, and agree on the dense subset $U_X \cap U_Y$, where they are determined by the $Y$-components (Proposition \ref{proposition:fourier-transform-genuine}).  They therefore agree.
\end{proof}

\subsection{Marginals and axis values of Schwartz functions}\label{sec:marginals-and-axis-values-of-schwartz-functions}
We extend the two relations from $\mathcal{D}$ to $\mathcal{S}$.  In Appendix \ref{sec:proof-poisson}, we apply the Poisson summation formula on $\mathbb{A}$ to products of the functions $\nu$ and $\sigma$ considered below.  For that purpose, we need them to be Schwartz--Bruhat functions on $F$.

\begin{lemma}\label{lemma:schwartz-marginals}
  Let $\Phi \in \mathcal{S}$ and $x', y' \in F^\times$.
  \begin{enumerate}[(i)]
  \item\label{item:schwartz-nu} The function $\nu$ of Lemma \ref{lemma:marginal-axis} is a Schwartz--Bruhat function on $F$.
  \item\label{item:schwartz-sigma} The function $\sigma(\tau) := \lvert n \rvert \lvert x' \rvert^{-1} \lvert \tau \rvert^{n-1} r_\Phi(- \tau^n / x')$ on $F^\times$ extends to a Schwartz--Bruhat function on $F$, and $m_{\mathcal{F} \Phi}(\mu^n x') = \lvert n \rvert^{-1/2} \hat{\sigma}(\mu)$ for $\mu \in F^\times$.
  \end{enumerate}
\end{lemma}
\begin{proof}
  For $\Phi \in \mathcal{D}$, these are Lemmas \ref{lemma:marginal-axis} and \ref{lemma:axis-marginal}.  Both assertions are linear in $\Phi$, so it remains to treat $\Phi = \mathcal{F} \Phi_0$ with $\Phi_0 \in \mathcal{D}$.

  \eqref{item:schwartz-nu} Let $\sigma_0$ denote the function attached to $\Phi_0$ and $1 / y'$ by Lemma \ref{lemma:axis-marginal}, so that $m_\Phi(\tau^n / y') = \lvert n \rvert^{-1/2} \hat{\sigma}_0(\tau)$.  Then $\nu = \lvert n \rvert^{-1/2} \{-1, y'\}_n \lvert y' \rvert^{-1} \hat{\sigma}_0$ on $F^\times$, and the right-hand side is a Schwartz--Bruhat function on $F$.  Since $\nu$ is continuous, the two functions agree on $F$.

  \eqref{item:schwartz-sigma} Let $\nu_0$ denote the function attached to $\Phi_0$ and $-1/x'$ by Lemma \ref{lemma:marginal-axis}, so that $\nu_0(t) = \{-1, -1/x'\}_n \lvert x' \rvert \, m_{\Phi_0}(- t^n x')$ and $r_\Phi(- \tau^n / x') = \lvert n \rvert^{-1/2} \lvert \tau \rvert^{1-n} \hat{\nu}_0(\tau)$.  Then $\sigma = \lvert n \rvert^{1/2} \lvert x' \rvert^{-1} \hat{\nu}_0$ on $F^\times$, and the right-hand side is a Schwartz--Bruhat function on $F$.  By Fourier inversion,
  \begin{equation*}
    \lvert n \rvert^{-1/2} \hat{\sigma}(\mu) = \lvert x' \rvert^{-1} \nu_0(- \mu)
    = \{-1, -1/x'\}_n \, m_{\Phi_0}(\varepsilon \mu^n x').
  \end{equation*}
  On the other hand, Proposition \ref{proposition:inversion} and Lemma \ref{lemma:torus}\eqref{item:torus-formula} give
  \begin{equation*}
    m_{\mathcal{F} \Phi}(\mu^n x') = m_{h(\varepsilon) \Phi_0}(\mu^n x')
    = \{x', \varepsilon\}_n \{-1, \varepsilon\}_n \, m_{\Phi_0}(\varepsilon \mu^n x').
  \end{equation*}
  The two symbol factors agree: for odd $n$, both are trivial, since $-1$ is then an $n$th power, while for even $n$, we have $\varepsilon = -1$ and $\{-1, -1/x'\}_n = \{-1, -1\}_n \{x', -1\}_n$.
\end{proof}

\subsection{The transform of the basic vector}\label{sec:basic-vector-transform}
We prove Proposition \ref{proposition:basic-vector-fixed}\eqref{item:basic-fixed}.  Suppose, as in \S\ref{sec:basic-vector}, that $F$ is non-archimedean, of residue characteristic not dividing $n$.  By \eqref{eq:axis-line-integral}, the values of $\mathcal{F} \Phi^0$ on the line $x = 0$ are determined by the following integrals:

\begin{lemma}\label{lemma:basic-vector-marginal}
  Suppose that $\psi$ is unramified.  For $x \in F^\times$, we have
  \begin{equation*}
    \int_F \Phi^0_X(x, y) \, d y =
    \begin{cases}
      1 & \text{ if } v(x) \in n \mathbb{Z}_{\geq 0}, \\
      0 & \text{ otherwise.}
    \end{cases}
  \end{equation*}
\end{lemma}
\begin{proof}
  Set $m := v(x)$.  The range $v(y) \geq m$ contributes $q^{(n-1) m / n} q^{-m} = q^{- m / n}$ if $m \in n \mathbb{Z}_{\geq 0}$, and zero otherwise.  In the range $v(y) < m$, the integrand vanishes unless $v(y) = n j$ with $j \geq 0$.  For such $y$, write $y = \varpi^{n j} u$ with $u \in \mathfrak{o}^\times$.  The integrand is then $q^{(n-1) j} \{x, y\}_n$, and $\{x, y\}_n = \{x, u\}_n = \{\varpi, u\}_n^{m}$.  By property \eqref{item:hilbert-tame} of the Hilbert symbol, the contribution of such $y$ is thus $q^{- j} (1 - q^{-1})$ if $n$ divides $m$, and zero otherwise.  If $m = n k$ with $k \geq 0$, the total is
  \begin{equation*}
    q^{-k} + \sum_{0 \leq j < k} q^{-j} (1 - q^{-1}) = 1.
    \qedhere
  \end{equation*}
\end{proof}

\begin{proof}[Proof of Proposition \ref{proposition:basic-vector-fixed}\eqref{item:basic-fixed}]
  Suppose that $\psi$ is unramified, and let $y' \in F^\times$.  By \eqref{eq:axis-line-integral} and Lemma \ref{lemma:basic-vector-marginal}, the function $\tau \mapsto R_Y \Phi^0((0, y'), \tau)$ on $F^\times$ vanishes unless $v(y') = n k$ for some $k \in \mathbb{Z}$, in which case it equals $q^{n k} 1_{v(\tau) \geq k}$; here $\{-1, y'\}_n = 1$ by properties \eqref{item:hilbert-powers} and \eqref{item:hilbert-units} of the Hilbert symbol.  In that case, its integral against $\psi(\tau)$ is $q^{(n-1) k} 1_{k \geq 0}$.  Thus
  \begin{equation*}
    (\mathcal{F} \Phi^0)_Y(0, y') = \Phi^0_Y(0, y').
  \end{equation*}
  It remains to pass from the line $x = 0$ to the rest of $U$.  By Lemmas \ref{lemma:unipotent-properties} and \ref{lemma:basic-vector-invariant}, the difference $\Delta := \mathcal{F} \Phi^0 - \Phi^0$ satisfies $\Delta_X(x, y + u x) = \Delta_X(x, y)$ and $\Delta_Y(x + u y, y) = \Delta_Y(x, y)$ for $u \in \mathfrak{o}$.  The function $\lvert \Delta \rvert$ on $U$ is thus invariant under the substitutions $(x, y) \mapsto (x, y + u x)$ and $(x, y) \mapsto (x + u y, y)$ with $u \in \mathfrak{o}$.  These carry each $(x, y) \in U$ to the line $x = 0$: if $\lvert x \rvert \leq \lvert y \rvert$, we apply the second with $u = - x / y$; otherwise, we first apply the first with $u = 1$, which gives $(x, x + y)$ with $\lvert x + y \rvert = \lvert x \rvert$.  Since $\lvert \Delta \rvert$ vanishes on that line, it vanishes on $U$.
\end{proof}

\section{Proof of the Poisson formula}\label{sec:proof-poisson}
We prove Theorem \ref{theorem:poisson}.  We retain the setting and notation of \S\ref{sec:poisson}.

\subsection{Strategy}\label{sec:strategy}
The usual proof of the Poisson summation formula applies the Fourier inversion formula on a compact quotient group.  Genuine functions are not functions on a group, so we argue differently.  We first treat the pure tensors $\Phi$ having a factor with leading term zero and a factor whose transform has leading term zero (Proposition \ref{proposition:poisson-vanishing}).  The local metaplectic groups act on the factors of $\mathcal{S}_{\mathbb{A}}$, compatibly with $\mathcal{F}$.  Each $\gamma \in \mathrm{SL}_2(F)$ lifts to an operator $\tilde{\gamma}$ on $\mathcal{S}_{\mathbb{A}}$ satisfying $(\tilde{\gamma} \Psi)(v) = \Psi(v \gamma)$ for $v \in F^2 - \{0\}$, which therefore preserves the sum of the values at those points (Lemma \ref{lemma:rational-invariance}).  The difference of the two sides, evaluated on the translates of $\Phi$, thus defines, after taking absolute values, a function $\delta$ on $\mathrm{SL}_2(F) \backslash \mathrm{SL}_2(\mathbb{A})$.  We show that $\delta = 0$ in three steps.
\begin{enumerate}[(i)]
\item The constant term of the difference, defined by integrating over the translates by $n(u)$ with $u \in \mathbb{A} / F$, vanishes (\S\ref{sec:constant-terms}).  This is where the hypotheses on the leading terms are used.
\item Consequently, $\delta$ decays rapidly at the cusp (Lemma \ref{lemma:decay}).
\item Each side is a sum over $F^2 - \{0\}$, which is a single orbit of $\mathrm{SL}_2(F)$ with unipotent stabilizers.  Unfolding the integral of $\delta^2$ and applying the first step again, we obtain $\int \delta^2 = 0$ (\S\ref{sec:proof-poisson-conclusion}).
\end{enumerate}
The general case follows by applying this one to suitable combinations of a pure tensor and its dilate at one finite place, for which the relevant leading terms vanish (\S\ref{sec:boundary}).

\subsection{Convergence}\label{sec:theta-sums}
For a place $\mathfrak{p}$ and $v \in F_{\mathfrak{p}}^2$, we write $\lVert v \rVert_{\mathfrak{p}}$ for the quantity $\lVert v \rVert$ defined using the normalized absolute value on $F_{\mathfrak{p}}$.  We write $\lvert a \rvert$ for the norm of an idele $a$, and set $\kappa := 1 - 1/n$.

Let $\Psi$ be a pure tensor.  By \eqref{eq:growth}, Corollary \ref{corollary:transform-decay} and the definition of $\Phi^0$, there are constants $C_{\mathfrak{p}} \geq 1$, with $C_{\mathfrak{p}} = 1$ for almost all $\mathfrak{p}$, and $C_{\mathfrak{p}, N} \geq 0$ for $N \geq 1$, such that for all $v \in F_{\mathfrak{p}}^2 - \{0\}$,
\begin{equation}\label{eq:local-bounds}
  \lvert \Psi_{\mathfrak{p}}(v) \rvert \leq
  \beta_{\mathfrak{p}, N}(v) :=
  \begin{cases}
    C_{\mathfrak{p}} \lVert v \rVert_{\mathfrak{p}}^{-\kappa} 1_{\lVert v \rVert_{\mathfrak{p}} \leq C_{\mathfrak{p}}}
    & \text{ if $\mathfrak{p}$ is finite,} \\
    C_{\mathfrak{p}, N} \min ( \lVert v \rVert_{\mathfrak{p}}^{-\kappa}, \lVert v \rVert_{\mathfrak{p}}^{-N} )
    & \text{ if $\mathfrak{p}$ is archimedean and $N \geq 1$.}
  \end{cases}
\end{equation}
For a family $\xi = (\xi_{\mathfrak{p}})_{\mathfrak{p}}$ with $\xi_{\mathfrak{p}} \in F_{\mathfrak{p}}^2 - \{0\}$, we write $\beta_N(\xi) := \prod_{\mathfrak{p}} \beta_{\mathfrak{p}, N}(\xi_{\mathfrak{p}})$ and $\lvert \Psi \rvert(\xi) := \prod_{\mathfrak{p}} \lvert \Psi_{\mathfrak{p}} \rvert(\xi_{\mathfrak{p}})$, whenever almost all factors are at most one.  Thus $\lvert \Psi \rvert \leq \beta_N$.  The functions $\beta_{\mathfrak{p}, N}$ depend only upon the constants in \eqref{eq:local-bounds}, and are nonincreasing functions of $\lVert v \rVert_{\mathfrak{p}}$.

\begin{lemma}\label{lemma:theta-bounds}
  For each idele $a$, we have
  \begin{equation*}
    \sum_{v \in F^2 - \{0\}} \beta_4(a v) \ll \lvert a \rvert^{-2 - \kappa},
  \end{equation*}
  where the implied constant depends only upon the constants in \eqref{eq:local-bounds}.
\end{lemma}
\begin{proof}
  Let $t$ be an idele, with $t_{\mathfrak{p}} = 1$ if $\mathfrak{p}$ is archimedean or $C_{\mathfrak{p}} = 1$, and with $(q_{\mathfrak{p}} C_{\mathfrak{p}})^{-1} < \lvert t_{\mathfrak{p}} \rvert_{\mathfrak{p}} \leq C_{\mathfrak{p}}^{-1}$ otherwise.  For finite $\mathfrak{p}$, we then have $\beta_{\mathfrak{p}, 4}(v) \leq C_{\mathfrak{p}} \lVert t_{\mathfrak{p}} v \rVert_{\mathfrak{p}}^{-\kappa} 1_{\lVert t_{\mathfrak{p}} v \rVert_{\mathfrak{p}} \leq 1}$.  Replacing $a$ by $a t$, we may thus assume that $C_{\mathfrak{p}} = 1$ for all finite $\mathfrak{p}$, up to a constant factor in the final estimate.  Let $v \in F^2 - \{0\}$, and set $s_{\mathfrak{p}} := \lvert a_{\mathfrak{p}} \rvert_{\mathfrak{p}} \lVert v \rVert_{\mathfrak{p}}$.  By the product formula, $\prod_{\mathfrak{p}} s_{\mathfrak{p}} \geq \lvert a \rvert$.  Thus
  \begin{equation}\label{eq:theta-term-bound}
    \beta_4(a v) \ll \lvert a \rvert^{- \kappa}
    \prod_{\mathfrak{p} \mid \infty} \min(1, s_{\mathfrak{p}}^{-3})
    \prod_{\mathfrak{p} \nmid \infty} 1_{s_{\mathfrak{p}} \leq 1}.
  \end{equation}
  Let $\rho_{\mathfrak{p}} > 0$, with $\rho_{\mathfrak{p}} = 1$ for almost all $\mathfrak{p}$, and set $R := \prod_{\mathfrak{p}} \rho_{\mathfrak{p}}$.  The number of nonzero $x \in F$ with $\lvert x \rvert_{\mathfrak{p}} \leq \rho_{\mathfrak{p}}$ at every place is $O_F(R)$.  For $R < 1$, no such $x$ exists, by the product formula.  For $R \geq 1$, subdivide the ball at one fixed finite place of residue cardinality $q$ into $q^k$ equal balls.  For two points in the same resulting box, the product of the local bounds on their difference is $\ll_F R q^{-k}$, by the triangle inequalities at infinity and the ultrametric inequalities elsewhere.  We may thus take $q^k \ll_F R$ so that each box contains at most one rational point.

  Let $j = (j_{\mathfrak{p}})_{\mathfrak{p} \mid \infty}$ be a tuple of integers, with sum $\lvert j \rvert$; if $F$ is a function field, there is only the empty tuple.  The number of $v \in F^2 - \{0\}$ with $s_{\mathfrak{p}} \leq 2^{j_{\mathfrak{p}}}$ for $\mathfrak{p} \mid \infty$ and $s_{\mathfrak{p}} \leq 1$ for $\mathfrak{p} \nmid \infty$ is therefore $O((2^{\lvert j \rvert} / \lvert a \rvert)^2)$.  Decomposing the sum over $v$ according to the least such $j$ and applying \eqref{eq:theta-term-bound}, we obtain
  \begin{equation*}
    \sum_{v \in F^2 - \{0\}} \beta_4(a v) \ll \lvert a \rvert^{- 2 - \kappa}
    \sum_{j} \,
    \prod_{\mathfrak{p} \mid \infty} 2^{2 j_{\mathfrak{p}}} \min(1, 2^{3 (1 - j_{\mathfrak{p}})}).
  \end{equation*}
  The sum over $j$ converges.
\end{proof}

By Lemma \ref{lemma:theta-bounds} with $a = 1$, the sum of $\Psi(v)$ over $v \in F^2 - \{0\}$ converges absolutely for each pure tensor $\Psi$, hence for each $\Psi \in \mathcal{S}_{\mathbb{A}}$.  We denote it by $\Theta(\Psi)$.  This defines a linear functional $\Theta$ on $\mathcal{S}_{\mathbb{A}}$, and gives the convergence asserted in Theorem \ref{theorem:poisson}.

We will gain decay from the following observation.  Let $M \geq 0$.  For $v \in F_{\mathfrak{p}}^2 - \{0\}$ and $N \geq 1$, we have
\begin{equation}\label{eq:weighted-bounds}
  \max(1, \lVert v \rVert_{\mathfrak{p}})^M \, \beta_{\mathfrak{p}, N + M}(v) \leq \beta'_{\mathfrak{p}, N}(v),
\end{equation}
where $\beta'_{\mathfrak{p}, N}$ is defined like $\beta_{\mathfrak{p}, N}$, with the constants $C_{\mathfrak{p}}^{M+1}$ and $C_{\mathfrak{p}, N + M}$ in place of $C_{\mathfrak{p}}$ and $C_{\mathfrak{p}, N}$.  We write $\beta'_N$ for the product of the functions $\beta'_{\mathfrak{p}, N}$; Lemma \ref{lemma:theta-bounds} applies to it.

\subsection{Unfolding the constant term}\label{sec:unipotent-average}
Let $\Psi$ be a pure tensor.  For $x, y \in F^\times$, we set
\begin{gather*}
  m_\Psi(x) := \prod_{\mathfrak{p}} m_{\Psi_{\mathfrak{p}}}(x),
  \qquad
  r_\Psi(y) := \prod_{\mathfrak{p}} r_{\Psi_{\mathfrak{p}}}(y),
  \\
  C(\Psi) := \sum_{y \in F^\times} r_\Psi(y) + \sum_{x \in F^\times} m_\Psi(x).
\end{gather*}
For $u \in \mathbb{A}$, we write $n(u) \Psi := \otimes_{\mathfrak{p}} n(u_{\mathfrak{p}}) \Psi_{\mathfrak{p}}$.  It is again a pure tensor in $\mathcal{S}_{\mathbb{A}}$, by Corollary \ref{corollary:metaplectic} and Lemma \ref{lemma:basic-vector-invariant}.

\begin{lemma}\label{lemma:constant-term}
  The sums defining $C(\Psi)$ converge absolutely, and
  \begin{equation*}
    \int_{\mathbb{A} / F} \Theta(n(u) \Psi) \, d u = C(\Psi).
  \end{equation*}
\end{lemma}
\begin{proof}
  By \eqref{eq:unipotent-axis}, the terms with $x = 0$ in the sum defining $\Theta(n(u) \Psi)$ do not depend upon $u$; they contribute $\sum_{y} r_\Psi(y)$.  The terms with $x \neq 0$ contribute
  \begin{equation*}
    \sum_{x \in F^\times} \int_{\mathbb{A} / F} \sum_{y \in F} \Psi_X(x, y + u x) \, d u
    = \sum_{x \in F^\times} \int_{\mathbb{A}} \Psi_X(x, u x) \, d u
    = \sum_{x \in F^\times} m_\Psi(x).
  \end{equation*}
  The convergence follows from Lemma \ref{lemma:theta-bounds}, applied to the pure tensors $n(u) \Psi$ with $u$ in a compact fundamental domain for $\mathbb{A} / F$; these satisfy \eqref{eq:local-bounds} with uniform constants, since $\lvert n(u_{\mathfrak{p}}) \Psi_{\mathfrak{p}} \rvert(v) = \lvert \Psi_{\mathfrak{p}} \rvert(v \, \pi(n(u_{\mathfrak{p}})))$ (Corollary \ref{corollary:metaplectic}).
\end{proof}

\subsection{Applying Poisson summation to the constant term}\label{sec:constant-terms}
We carry out the first step of the strategy.  By Lemma \ref{lemma:constant-term}, we must compare sums of marginals and axis values.  We do so by applying the Poisson summation formula on $\mathbb{A}$ to each of the two relations of Lemma \ref{lemma:schwartz-marginals}.  Each application produces a term at the origin on either side.  Of these four terms, two are given by the leading terms of $\Phi$ and $\mathcal{F} \Phi$; the remaining two cancel, by a Poisson summation formula on $\mathbb{A}^\times / \mathbb{A}^{\times n}$ (Lemma \ref{lemma:isotropy}).

We write $\Delta := F^\times / F^{\times n}$.  The following is a form of the Poisson summation formula for the image of $F^\times$ in $\mathbb{A}^\times / \mathbb{A}^{\times n}$, which is its own orthogonal complement with respect to the Hilbert symbol \cite[XIII.5, Proposition 8]{MR0427267}.

\begin{lemma}\label{lemma:isotropy}
  For each place $\mathfrak{p}$, let $\beta_{\mathfrak{p}}$ be a function on $F_{\mathfrak{p}}^\times / F_{\mathfrak{p}}^{\times n}$.  Assume that for almost all $\mathfrak{p}$, the function $\beta_{\mathfrak{p}}$ is the characteristic function of the image of $\mathfrak{o}_{\mathfrak{p}}^\times$.  Set
  \begin{equation*}
    \tilde{\beta}_{\mathfrak{p}}(y) := \frac{1}{n} \sum_{c \in F_{\mathfrak{p}}^\times / F_{\mathfrak{p}}^{\times n}} \{c, y\}_{\mathfrak{p}, n} \, \beta_{\mathfrak{p}}(c).
  \end{equation*}
  Then
  \begin{equation*}
    \sum_{x \in \Delta} \prod_{\mathfrak{p}} \beta_{\mathfrak{p}}(x)
    =
    \sum_{y \in \Delta} \prod_{\mathfrak{p}} \tilde{\beta}_{\mathfrak{p}}(y).
  \end{equation*}
\end{lemma}
\begin{proof}
  The group $F_{\mathfrak{p}}^\times / F_{\mathfrak{p}}^{\times n}$ has order $n^2 / \lvert n \rvert_{\mathfrak{p}}$, and the Hilbert symbol identifies it with its dual.  The transform $\beta_{\mathfrak{p}} \mapsto \lvert n \rvert_{\mathfrak{p}}^{1/2} \tilde{\beta}_{\mathfrak{p}}$ is thus the Fourier transform for the self-dual measure.  By the product formula, the product of these local transforms is the Fourier transform on $\mathbb{A}^\times / \mathbb{A}^{\times n}$ for the self-dual measure.  The assertion is the Poisson summation formula for the discrete cocompact subgroup $\Delta$, which equals its own orthogonal complement.
\end{proof}

Let $s = \otimes_{\mathfrak{p}} s_{\mathfrak{p}}$ be a pure tensor in $b(\mathcal{S}_{\mathbb{A}})$; thus $s_{\mathfrak{p}} \in b(\mathcal{S}_{\mathfrak{p}})$ for every $\mathfrak{p}$, and $s_{\mathfrak{p}} = b \Phi_{\mathfrak{p}}^0$ for almost every $\mathfrak{p}$.  For $y \in F^\times$, we set $s_Y(0, y) := \prod_{\mathfrak{p}} (s_{\mathfrak{p}})_Y(0, y)$.  Almost all factors are equal to one, and the product depends only upon the class of $y$ in $\Delta$, by the homogeneity of the leading terms and the product formula.  It vanishes unless $n$ divides $v_{\mathfrak{p}}(y)$ for all $\mathfrak{p}$ outside a finite set depending only upon $s$ (\S\ref{sec:germs}).  The classes of such $y$ form a finite subgroup of $\Delta$, by the finiteness of the class group and the finite generation of the unit groups.  We may thus set
\begin{equation*}
  T(s) := \frac{1}{n} \sum_{y \in \Delta} s_Y(0, y).
\end{equation*}

\begin{proposition}\label{proposition:constant-terms}
  Let $\Phi$ be a pure tensor.  Then $C(\Phi) + T(b \Phi) = C(\mathcal{F} \Phi) + T(b \mathcal{F} \Phi)$.
\end{proposition}
\begin{proof}
  Let $y' \in F^\times$.  For each $\mathfrak{p}$, let $\nu_{\mathfrak{p}}$ be the Schwartz--Bruhat function attached to $\Phi_{\mathfrak{p}}$ and $y'$ by Lemmas \ref{lemma:marginal-axis} and \ref{lemma:schwartz-marginals}.  For almost all $\mathfrak{p}$, it is the characteristic function of $\mathfrak{o}_{\mathfrak{p}}$ (Lemma \ref{lemma:basic-vector-marginal}).  Let $\nu := \otimes \nu_{\mathfrak{p}}$.  By the reciprocity law and the product formula, which applies also to the factors $\lvert n \rvert_{\mathfrak{p}}^{-1/2}$, we have for $\tau, \mu \in F^\times$
  \begin{equation*}
    \nu(\tau) = m_\Phi(\tau^n / y'),
    \qquad
    \hat{\nu}(\mu) = r_{\mathcal{F} \Phi}(\mu^n y').
  \end{equation*}
  By continuity and the definition of the leading term, $\hat{\nu}_{\mathfrak{p}}(0) = \lvert n \rvert_{\mathfrak{p}}^{1/2} (b \mathcal{F} \Phi_{\mathfrak{p}})_Y(0, y')$ for each $\mathfrak{p}$, so $\hat{\nu}(0) = (b \mathcal{F} \Phi)_Y(0, y')$.  The Poisson summation formula on $\mathbb{A}$ thus gives
  \begin{equation*}
    \sum_{\tau \in F^\times} m_\Phi(\tau^n / y') + \nu(0) = \sum_{\mu \in F^\times} r_{\mathcal{F} \Phi}(\mu^n y') + (b \mathcal{F} \Phi)_Y(0, y').
  \end{equation*}
  We sum over $y' \in \Delta$, using the formula for $\nu_{\mathfrak{p}}(0)$ in Lemma \ref{lemma:marginal-axis}.  Since the map $\tau \mapsto \tau^n$ on $F^\times$ is $n$-to-one, we obtain
  \begin{equation}\label{eq:constant-terms-1}
    n \sum_{x \in F^\times} m_\Phi(x)
    + \sum_{y' \in \Delta} \prod_{\mathfrak{p}} \int_{F_{\mathfrak{p}}} \{y, y'\}_{\mathfrak{p}, n} r_{\Phi_{\mathfrak{p}}}(y) \, d y
    = n \sum_{y \in F^\times} r_{\mathcal{F} \Phi}(y) + n \, T(b \mathcal{F} \Phi).
  \end{equation}

  Let $x' \in F^\times$.  For each $\mathfrak{p}$, let $\sigma_{\mathfrak{p}}$ be the Schwartz--Bruhat function attached to $\Phi_{\mathfrak{p}}$ and $x'$ by Lemma \ref{lemma:schwartz-marginals}.  For almost all $\mathfrak{p}$, it is the characteristic function of $\mathfrak{o}_{\mathfrak{p}}$, by the definition of $\Phi^0$.  Let $\sigma := \otimes \sigma_{\mathfrak{p}}$.  By continuity and the definition of the leading term, $\sigma_{\mathfrak{p}}(0) = \lvert n \rvert_{\mathfrak{p}} \lvert x' \rvert_{\mathfrak{p}}^{-1} (b \Phi_{\mathfrak{p}})_Y(0, -1/x')$ for each $\mathfrak{p}$, so $\sigma(0) = (b \Phi)_Y(0, -1/x')$.  As above, we have for $\tau, \mu \in F^\times$
  \begin{equation*}
    \sigma(\tau) = r_\Phi(- \tau^n / x'),
    \qquad
    \hat{\sigma}(\mu) = m_{\mathcal{F} \Phi}(\mu^n x').
  \end{equation*}
  By \eqref{eq:haar-on-powers}, we have $\int_{F_{\mathfrak{p}}} \sigma_{\mathfrak{p}} = n \int_{- F_{\mathfrak{p}}^{\times n} / x'} r_{\Phi_{\mathfrak{p}}}(y) \, d y$.  Summing the Poisson summation formula for $\sigma$ over $x' \in \Delta$, we obtain
  \begin{equation}\label{eq:constant-terms-2}
    n \sum_{y \in F^\times} r_\Phi(y) + n \, T(b \Phi)
    =
    n \sum_{x \in F^\times} m_{\mathcal{F} \Phi}(x)
    + \sum_{x' \in \Delta} \prod_{\mathfrak{p}} n \int_{- F_{\mathfrak{p}}^{\times n} / x'} r_{\Phi_{\mathfrak{p}}}(y) \, d y.
  \end{equation}

  It remains to compare the sums over $\Delta$ in \eqref{eq:constant-terms-1} and \eqref{eq:constant-terms-2}.  Set $\beta_{\mathfrak{p}}(c) := n \int_{c F_{\mathfrak{p}}^{\times n}} r_{\Phi_{\mathfrak{p}}}(y) \, d y$.  This is the integral of the function $\sigma_{\mathfrak{p}}$ attached to $x' = -1/c$.  At a place where $\Phi_{\mathfrak{p}} = \Phi_{\mathfrak{p}}^0$ and $\psi_{\mathfrak{p}}$ is unramified, the definition of $\Phi^0$ shows that $\beta_{\mathfrak{p}}$ is the characteristic function of the image of $\mathfrak{o}_{\mathfrak{p}}^\times$.  With $\tilde{\beta}_{\mathfrak{p}}$ as in Lemma \ref{lemma:isotropy}, we have $\tilde{\beta}_{\mathfrak{p}}(y') = \int_{F_{\mathfrak{p}}} \{y, y'\}_{\mathfrak{p}, n} r_{\Phi_{\mathfrak{p}}}(y) \, d y$.  By that lemma and the reindexing $c = -1/x'$,
  \begin{align*}
    \sum_{y' \in \Delta} \prod_{\mathfrak{p}} \int_{F_{\mathfrak{p}}} \{y, y'\}_{\mathfrak{p}, n} r_{\Phi_{\mathfrak{p}}}(y) \, d y
    &= \sum_{y' \in \Delta} \prod_{\mathfrak{p}} \tilde{\beta}_{\mathfrak{p}}(y')
      = \sum_{c \in \Delta} \prod_{\mathfrak{p}} \beta_{\mathfrak{p}}(c) \\
    &= \sum_{x' \in \Delta} \prod_{\mathfrak{p}} n \int_{- F_{\mathfrak{p}}^{\times n} / x'} r_{\Phi_{\mathfrak{p}}}(y) \, d y.
  \end{align*}
  Thus the sums over $\Delta$ in \eqref{eq:constant-terms-1} and \eqref{eq:constant-terms-2} are equal.  Adding the two identities, we obtain $n (C(\Phi) + T(b \Phi)) = n (C(\mathcal{F} \Phi) + T(b \mathcal{F} \Phi))$.
\end{proof}

\subsection{Lifts}\label{sec:rational-lifts}
We make precise the action of $\mathrm{SL}_2(\mathbb{A})$, up to scalars, on $\mathcal{S}_{\mathbb{A}}$.  For each place $\mathfrak{p}$, let $\pi_{\mathfrak{p}} : \widetilde{\mathrm{SL}}_2(F_{\mathfrak{p}}) \rightarrow \mathrm{SL}_2(F_{\mathfrak{p}})$ be as in Proposition \ref{proposition:kubota}.  For $\mathfrak{p}$ finite, of residue characteristic not dividing $n$, the elements of $\widetilde{\mathrm{SL}}_2(F_{\mathfrak{p}})$ that are products of operators $n(u)$ and $n^-(u)$ with $u \in \mathfrak{o}_{\mathfrak{p}}$ fix $\Phi^0$ (Lemma \ref{lemma:basic-vector-invariant}); we call them \emph{integral}.  Their images under $\pi_{\mathfrak{p}}$ exhaust $\mathrm{SL}_2(\mathfrak{o}_{\mathfrak{p}})$.

A \emph{lift} of $m \in \mathrm{SL}_2(\mathbb{A})$ is a family $g = (g_{\mathfrak{p}})_{\mathfrak{p}}$ with $g_{\mathfrak{p}} \in \widetilde{\mathrm{SL}}_2(F_{\mathfrak{p}})$ and $\pi_{\mathfrak{p}}(g_{\mathfrak{p}}) = m_{\mathfrak{p}}$ for all $\mathfrak{p}$, and with $g_{\mathfrak{p}}$ integral for almost all $\mathfrak{p}$.  Lifts exist.  For a pure tensor $\Psi$, we write $g \Psi := \otimes_{\mathfrak{p}} g_{\mathfrak{p}} \Psi_{\mathfrak{p}}$.  It is again a pure tensor in $\mathcal{S}_{\mathbb{A}}$.  This defines a linear map $g$ on $\mathcal{S}_{\mathbb{A}}$, and $\mathcal{F} g = g \mathcal{F}$, by Theorem \ref{theorem:schwartz-space}.  If $g$ and $g'$ are lifts of the same $m$, then $g'_{\mathfrak{p}} = \zeta_{\mathfrak{p}} g_{\mathfrak{p}}$ with $\zeta_{\mathfrak{p}} \in \mu_n$ for each $\mathfrak{p}$, and $\zeta_{\mathfrak{p}} = 1$ whenever both components are integral, since integral elements fix $\Phi_{\mathfrak{p}}^0$.  The maps $g$ and $g'$ on $\mathcal{S}_{\mathbb{A}}$ thus differ by the scalar $\prod_{\mathfrak{p}} \zeta_{\mathfrak{p}} \in \mu_n$.

For $u \in \mathbb{A}$ and $a \in \mathbb{A}^\times$, the families $n(u) := (n(u_{\mathfrak{p}}))_{\mathfrak{p}}$ and $h(a) := (h(a_{\mathfrak{p}}))_{\mathfrak{p}}$ are lifts of the corresponding matrices, by \S\ref{sec:torus} and Lemma \ref{lemma:basic-vector-invariant}.

\begin{lemma}\label{lemma:rational-invariance}
  Let $\gamma \in \mathrm{SL}_2(F)$.  There is a lift $\tilde{\gamma}$ of $\gamma$ such that, for each $\Psi \in \mathcal{S}_{\mathbb{A}}$ and each $v \in F^2 - \{0\}$, we have $(\tilde{\gamma} \Psi)(v) = \Psi(v \gamma)$.  In particular, $\Theta(\tilde{\gamma} \Psi) = \Theta(\Psi)$.
\end{lemma}
\begin{proof}
  The group $\mathrm{SL}_2(F)$ is generated by the matrices $\pi(n(u))$ and $\pi(n^-(u))$ with $u \in F$.  It thus suffices to treat $\tilde{\gamma} = n(u)$ and $\tilde{\gamma} = n^-(u)$ with $u \in F$, in which case the assertion follows, for pure tensors, from \eqref{eq:unipotent-charts} and \eqref{eq:unipotent-axis}.
\end{proof}

\subsection{The case of vanishing leading terms}\label{sec:proof-poisson-conclusion}
For $\Psi \in \mathcal{S}_{\mathbb{A}}$, we set
\begin{equation*}
  D(\Psi) := \Theta(\Psi) - \Theta(\mathcal{F} \Psi).
\end{equation*}
This defines a linear functional $D$ on $\mathcal{S}_{\mathbb{A}}$.

\begin{proposition}\label{proposition:poisson-vanishing}\lean{MetaplecticPoisson.Datum.thetaDifference_vanishing}
  Let $\Phi$ be a pure tensor.  Suppose that $b \Phi_{\mathfrak{p}} = 0$ for some place $\mathfrak{p}$, and that $b \mathcal{F} \Phi_{\mathfrak{q}} = 0$ for some place $\mathfrak{q}$.  Then $D(\Phi) = 0$.
\end{proposition}
The proof occupies the remainder of this subsection.  Let $\Phi$ be as in the proposition.  For $m \in \mathrm{SL}_2(\mathbb{A})$, with lift $g$, we set
\begin{gather*}
  \delta(m) := \lvert D(g \Phi) \rvert,
  \\
  \theta(m) := \sum_{v \in F^2 - \{0\}} \left( \lvert g \Phi \rvert(v) + \lvert g \mathcal{F} \Phi \rvert(v) \right).
\end{gather*}
The functions $\delta$ and $\theta$ do not depend upon the choice of lift.  By Lemma \ref{lemma:rational-invariance}, they are left-invariant under $\mathrm{SL}_2(F)$.  We must show that $\delta(1) = 0$.

\begin{lemma}\label{lemma:translated-constant-terms}
  For every lift $g$, we have
  \begin{equation*}
    \int_{\mathbb{A} / F} D(n(u) g \Phi) \, d u = 0.
  \end{equation*}
\end{lemma}
\begin{proof}
  By Lemma \ref{lemma:rational-invariance}, the integrand is a function on $\mathbb{A} / F$.  Since $\mathcal{F} g \Phi = g \mathcal{F} \Phi$, the integral is $C(g \Phi) - C(\mathcal{F} g \Phi)$, by Lemma \ref{lemma:constant-term}.  By Proposition \ref{proposition:constant-terms}, this equals $T(b \mathcal{F} g \Phi) - T(b g \Phi)$.  The leading term commutes with the local operators (\S\ref{sec:germs}).  The pure tensor $b g \Phi$ thus has $\mathfrak{p}$-factor $g_{\mathfrak{p}} \, b \Phi_{\mathfrak{p}} = 0$, and the pure tensor $b \mathcal{F} g \Phi$ has $\mathfrak{q}$-factor $g_{\mathfrak{q}} \, b \mathcal{F} \Phi_{\mathfrak{q}} = 0$.  Both values of $T$ therefore vanish.
\end{proof}

To estimate $\delta$ and $\theta$, we recall some reduction theory \cite[I.2.1]{MR1361168}.  Fix a finite place $w$, and let $z$ denote the idele with component $\varpi_w^{-1}$ at $w$ and $1$ elsewhere.  Thus $\lvert z \rvert = q_w > 1$, and each idele is the product of an element of $F^\times$, a power of $z$, and an element of a fixed compact set.  There is a compact subset $\Omega$ of $\mathrm{SL}_2(\mathbb{A})$ such that every element of $\mathrm{SL}_2(\mathbb{A})$ is of the form $\gamma \, \pi(n(u) h(z^k)) \, \omega$ with $\gamma \in \mathrm{SL}_2(F)$, $u \in \mathbb{A}$, $k \geq 0$ and $\omega \in \Omega$.  For a nonnegative measurable function $f$ on $\mathrm{SL}_2(F) \backslash \mathrm{SL}_2(\mathbb{A})$, we have
\begin{equation}\label{eq:siegel}
  \int_{\mathrm{SL}_2(F) \backslash \mathrm{SL}_2(\mathbb{A})} f(m) \, d m
  \ll
  \sum_{k \geq 0} \lvert z \rvert^{-2 k}
  \int_{\mathbb{A} / F} \int_{\Omega} f(\pi(n(u) h(z^k)) \omega) \, d \omega \, d u,
\end{equation}
for a suitable finite measure $d \omega$ on $\Omega$.  The factor $\lvert z \rvert^{-2 k}$ arises because conjugation by $h(z^k)$ multiplies the Haar measure on the group of matrices $\pi(n(u))$ by $\lvert z \rvert^{2 k}$.

Let $\tilde{\Omega}$ denote the set of lifts of the elements of $\Omega$.  For each lift $g$ of $m \in \mathrm{SL}_2(\mathbb{A})$, each pure tensor $\Psi$ and each place $\mathfrak{p}$, we have $\lvert g_{\mathfrak{p}} \Psi_{\mathfrak{p}} \rvert(v) = \lvert \Psi_{\mathfrak{p}} \rvert(v \, m_{\mathfrak{p}})$ (Corollary \ref{corollary:metaplectic}).  Since $\Omega$ is compact, there is a finite set of places outside of which $m_{\mathfrak{p}} \in \mathrm{SL}_2(\mathfrak{o}_{\mathfrak{p}})$ for every $m \in \Omega$, so that $\lVert v \, m_{\mathfrak{p}} \rVert_{\mathfrak{p}} = \lVert v \rVert_{\mathfrak{p}}$; at each remaining place, $\lVert v \, m_{\mathfrak{p}} \rVert_{\mathfrak{p}} \asymp \lVert v \rVert_{\mathfrak{p}}$, uniformly for $m \in \Omega$.  The pure tensors $\omega \Phi$ and $\omega \mathcal{F} \Phi$ with $\omega \in \tilde{\Omega}$ thus satisfy \eqref{eq:local-bounds} with uniform constants.  The same applies to any compact subset of $\mathrm{SL}_2(\mathbb{A})$ in place of $\Omega$.

\begin{lemma}\label{lemma:growth}
  For $u \in \mathbb{A}$, $k \geq 0$ and $\omega \in \tilde{\Omega}$, we have $\theta(\pi(n(u) h(z^k) \omega)) \ll \lvert z \rvert^{3 k}$.
\end{lemma}
\begin{proof}
  By Lemma \ref{lemma:rational-invariance} and the compactness of $\mathbb{A} / F$, we may assume that $u$ lies in a fixed compact set.  By Lemma \ref{lemma:torus}, we have $n(u) h(z^k) = h(z^k) n(u')$, where $u' := u / z^{2 k}$ lies in a compact set independent of $k$.  Let $\Psi$ be of the form $n(u') \omega \Phi$ or $n(u') \omega \mathcal{F} \Phi$ with $\omega \in \tilde{\Omega}$; such pure tensors satisfy \eqref{eq:local-bounds} with uniform constants.  For $(x, y) \in F^2 - \{0\}$, we have
  \begin{equation}\label{eq:cusp-term}
    \lvert h(z^k) \Psi \rvert(x, y) = \lvert \Psi \rvert(z^k x, z^{-k} y) \leq \beta_4(z^k x, z^{-k} y) \leq \beta_4(z^{-k} (x, y)).
  \end{equation}
  Here $(z^k x, z^{-k} y)$ denotes the family with component $(\varpi_w^{-k} x, \varpi_w^k y)$ at $w$ and $(x, y)$ elsewhere.  The last inequality holds because $\beta_{w, 4}$ is a nonincreasing function of $\lVert \cdot \rVert_w$, and $\lVert (\varpi_w^{-k} x, \varpi_w^{k} y) \rVert_w \geq \lVert \varpi_w^{k} (x, y) \rVert_w$.  Lemma \ref{lemma:theta-bounds}, applied to the idele $z^{-k}$, thus gives $\theta(\pi(n(u) h(z^k) \omega)) \ll \lvert z \rvert^{(2 + \kappa) k}$.
\end{proof}

\begin{lemma}\label{lemma:decay}
  For each $M \geq 0$, and for $u \in \mathbb{A}$, $k \geq 0$ and $\omega \in \tilde{\Omega}$, we have $\delta(\pi(n(u) h(z^k) \omega)) \ll \lvert z \rvert^{-M k}$.
\end{lemma}
\begin{proof}
  Fix $k$ and $\omega \in \tilde{\Omega}$, and set $\phi(u) := D(n(u) h(z^k) \omega \Phi)$.  By Lemmas \ref{lemma:rational-invariance} and \ref{lemma:translated-constant-terms}, $\phi$ is a function on $\mathbb{A} / F$ of integral zero.  For a pure tensor $\Psi$, the terms with $x = 0$ in the sum defining $\Theta(n(u) \Psi)$ do not depend upon $u$, by \eqref{eq:unipotent-axis}.  Let $\phi_1(u)$ denote the contribution to $\phi(u)$ of the terms with $x \neq 0$ in the sums defining $\Theta(n(u) h(z^k) \omega \Phi)$ and $\Theta(n(u) h(z^k) \omega \mathcal{F} \Phi)$.  Then $\phi - \phi_1$ is constant.  Since $\phi$ has integral zero, that constant is minus the integral of $\phi_1$ over $\mathbb{A} / F$, and so
  \begin{equation*}
    \sup \lvert \phi \rvert \leq 2 \sup \lvert \phi_1 \rvert.
  \end{equation*}
  It thus suffices to show that $\phi_1(u) \ll \lvert z \rvert^{-M k}$ for $u$ in a fixed compact set.

  Let $\Psi$ be as in the proof of Lemma \ref{lemma:growth}.  By \eqref{eq:cusp-term}, we must bound the sum of $\lvert \Psi \rvert(z^k x, z^{-k} y)$ over $x \in F^\times$ and $y \in F$.  For such $x$, the product formula gives
  \begin{equation*}
    \prod_{\mathfrak{p}} \max(1, \lvert (z^k)_{\mathfrak{p}} \, x \rvert_{\mathfrak{p}}) \geq \prod_{\mathfrak{p}} \lvert (z^k)_{\mathfrak{p}} \, x \rvert_{\mathfrak{p}} = \lvert z \rvert^{k}.
  \end{equation*}
  Let $M' \geq 0$.  By \eqref{eq:weighted-bounds}, with $M'$ in place of $M$, we thus have
  \begin{equation*}
    \lvert \Psi \rvert(z^k x, z^{-k} y)
    \leq \beta_{4 + M'}(z^k x, z^{-k} y)
    \leq \lvert z \rvert^{- M' k} \, \beta'_4(z^k x, z^{-k} y).
  \end{equation*}
  As in the proof of Lemma \ref{lemma:growth}, the sum of $\beta'_4(z^k x, z^{-k} y)$ over $(x, y) \in F^2 - \{0\}$ is $O(\lvert z \rvert^{3 k})$.  We take $M' = M + 3$.
\end{proof}

We now carry out the last step of the strategy of \S\ref{sec:strategy}.  By \eqref{eq:siegel} and Lemmas \ref{lemma:growth} and \ref{lemma:decay}, the functions $\delta^2$ and $\theta \delta$ are integrable on $\mathrm{SL}_2(F) \backslash \mathrm{SL}_2(\mathbb{A})$.  Let $\Psi \in \{\Phi, \mathcal{F} \Phi\}$.  By Lemma \ref{lemma:rational-invariance}, we have for each lift $g$ the identities
\begin{equation*}
  \Theta(g \Psi) = \sum_{\gamma \in N(F) \backslash \mathrm{SL}_2(F)} (\tilde{\gamma} g \Psi)(0, 1),
  \qquad
  D(\tilde{\gamma} g \Phi) = D(g \Phi),
\end{equation*}
where $N(F)$ denotes the stabilizer of $(0, 1)$, which is the group of matrices $\pi(n(u))$ with $u \in F$.  Unfolding the sum over $\gamma$ against the integral over $\mathrm{SL}_2(F) \backslash \mathrm{SL}_2(\mathbb{A})$, we obtain
\begin{multline*}
  \int_{\mathrm{SL}_2(F) \backslash \mathrm{SL}_2(\mathbb{A})} \Theta(g \Psi) \overline{D(g \Phi)} \, d m
  \\
  =
  \int_{N(\mathbb{A}) \backslash \mathrm{SL}_2(\mathbb{A})} (g \Psi)(0, 1) \left( \int_{\mathbb{A} / F} \overline{D(n(u) g \Phi)} \, d u \right) \, d m.
\end{multline*}
Here $g$ denotes a lift of $m$.  The integrands do not depend upon its choice.  The manipulation is justified by the integrability of $\theta \delta$.  The inner integral vanishes by Lemma \ref{lemma:translated-constant-terms}.  Taking the difference over $\Psi = \Phi$ and $\Psi = \mathcal{F} \Phi$, we obtain
\begin{equation*}
  \int_{\mathrm{SL}_2(F) \backslash \mathrm{SL}_2(\mathbb{A})} \delta(m)^2 \, d m = 0.
\end{equation*}
Thus $\delta$ vanishes almost everywhere.

Let $K$ denote the group of lifts $k$, with trivial components at the archimedean places (if any), such that $k_{\mathfrak{p}} \Phi_{\mathfrak{p}} = \Phi_{\mathfrak{p}}$ for each finite $\mathfrak{p}$.  Such lifts also fix $\mathcal{F} \Phi$, so $\delta$ is right-invariant under the image of $K$.  For each finite $\mathfrak{p}$, the element $\Phi_{\mathfrak{p}}$ of $\mathcal{S}$ is fixed by $n(u)$ and $n^-(u)$ for $u$ sufficiently small: this is clear for elements of $\mathcal{D}$, and follows for their transforms because $\mathcal{F}$ commutes with these operators.  When $\Phi_{\mathfrak{p}} = \Phi_{\mathfrak{p}}^0$, it is fixed by the integral elements.  The image of $K$ thus contains an open subgroup of the group $\mathrm{SL}_2(\mathbb{A}_f)$ of elements with trivial archimedean components.  The function $\delta$ is also continuous in the archimedean components.  Indeed, let $\mathfrak{p}$ be an archimedean place.  Near a given element of $\mathrm{SL}_2(F_{\mathfrak{p}})$, we may choose lifts of the form $g_0 g_1$, where $g_0$ is fixed and $g_1$ is a product of operators $n(u)$ and $n^-(u)$ whose parameters depend continuously upon the element.  For $\Psi \in \mathcal{S}_{\mathfrak{p}}$, the functions $(u, v) \mapsto (n(u) \Psi)_X(v)$ and $(u, v) \mapsto (n(u) \Psi)_Y(v)$ are continuous on $F_{\mathfrak{p}} \times U_X$ and $F_{\mathfrak{p}} \times U_Y$, by \eqref{eq:unipotent-charts} and \eqref{eq:unipotent-other-chart}, and likewise for $n^-(u)$.  The terms of the sums defining $\Theta(g \Phi)$ and $\Theta(g \mathcal{F} \Phi)$ thus depend continuously upon the archimedean components, and these sums converge locally uniformly in those components, by Lemma \ref{lemma:theta-bounds}.  The function $\delta$ is thus continuous on $\mathrm{SL}_2(\mathbb{A})$.
Hence $\delta$ vanishes identically, and in particular $\delta(1) = 0$.  This proves Proposition \ref{proposition:poisson-vanishing}.

\subsection{The boundary term}\label{sec:boundary}
We now prove Theorem \ref{theorem:poisson}.  By Proposition \ref{proposition:poisson-vanishing}, the linear functional $D$ vanishes on each pure tensor having a factor with leading term zero and a factor whose transform has leading term zero.  To define $B$, we modify a given pure tensor at one place, without changing its leading term, so that the transform of its factor at that place has leading term zero.  This is possible because dilations scale the leading term of a function and that of its transform differently.

Fix a finite place $w$, and write $q$ and $\varpi$ for the residue field cardinality and a uniformizer at $w$.  Let $\Phi_w \in \mathcal{S}_w$.  By \S\ref{sec:germs} and Lemma \ref{lemma:transform-dilation}, we have $b (\Phi_w(\varpi^n \, \cdot)) = q^{n-1} \, b \Phi_w$ and $b \mathcal{F} (\Phi_w(\varpi^n \, \cdot)) = q^{n+1} \, b \mathcal{F} \Phi_w$.  The function
\begin{equation*}
  Q_w \Phi_w := \frac{q^{n+1} \Phi_w - \Phi_w(\varpi^n \, \cdot)}{q^{n+1} - q^{n-1}}
\end{equation*}
thus satisfies
\begin{equation}\label{eq:projector}
  b \, Q_w \Phi_w = b \Phi_w,
  \qquad
  b \mathcal{F} Q_w \Phi_w = 0.
\end{equation}
We write $Q$ for the linear map on $\mathcal{S}_{\mathbb{A}}$ given by $Q_w$ acting in the $w$-factor: for a pure tensor $\Phi$, the pure tensor $Q \Phi$ is obtained from $\Phi$ by replacing $\Phi_w$ with $Q_w \Phi_w$.

\begin{lemma}\label{lemma:boundary-well-defined}
  Let $\Phi \in \mathcal{S}_{\mathbb{A}}$ with $b \Phi = 0$.  Then $D(Q \Phi) = 0$.
\end{lemma}
\begin{proof}
  Choose a finite set $T$ of places large enough to represent $\Phi$: thus $\Phi$ is the tensor product of an element $\Phi_T$ of $\bigotimes_{\mathfrak{p} \in T} \mathcal{S}_{\mathfrak{p}}$ with the basic vectors at the places outside $T$.  Since the vectors $b \Phi_{\mathfrak{p}}^0$ are nonzero, the image of $\Phi_T$ in $\bigotimes_{\mathfrak{p} \in T} b(\mathcal{S}_{\mathfrak{p}})$ vanishes.  The kernel of a tensor product of finitely many linear maps is spanned by the pure tensors having a factor in the kernel of the corresponding map.  Thus $\Phi$ is a finite sum of pure tensors, each having a factor with leading term zero.  By linearity, we may suppose that $\Phi$ is such a pure tensor, with $b \Phi_{\mathfrak{p}} = 0$.  By \eqref{eq:projector}, the $\mathfrak{p}$-factor of $Q \Phi$ has leading term zero, whether or not $\mathfrak{p} = w$, and the transform of its $w$-factor has leading term zero.  Thus $D(Q \Phi) = 0$, by Proposition \ref{proposition:poisson-vanishing}.
\end{proof}

By Lemma \ref{lemma:boundary-well-defined}, the linear functional $D \circ Q$ vanishes on the kernel of the surjection $b$.  There is thus a unique linear functional $B$ on $b(\mathcal{S}_{\mathbb{A}})$ such that, for all $\Phi \in \mathcal{S}_{\mathbb{A}}$,
\begin{equation*}
  B(b \Phi) = - D(Q \Phi).
\end{equation*}

\emph{The formula.}  We use two observations.  First, $D(\mathcal{F} \Psi) = - D(\Psi)$ for each $\Psi \in \mathcal{S}_{\mathbb{A}}$.  Indeed, we may suppose that $\Psi$ is a pure tensor.  By Theorem \ref{theorem:schwartz-space}, we have $\mathcal{F} \mathcal{F} \Psi = h(\varepsilon) \Psi := \otimes_{\mathfrak{p}} h(\varepsilon) \Psi_{\mathfrak{p}}$.  Since $h(\varepsilon)$ is a product of operators $n(u)$ and $n^-(u)$ with $u \in \{\pm 1\}$, the proof of Lemma \ref{lemma:rational-invariance} shows that $(h(\varepsilon) \Psi)(v) = \Psi(\varepsilon v)$.  Hence $\Theta(\mathcal{F} \mathcal{F} \Psi) = \Theta(\Psi)$.

Second, let $\Psi$ be a pure tensor with $b \mathcal{F} \Psi_w = 0$.  By \eqref{eq:projector}, the $w$-factor $\Psi_w - Q_w \Psi_w$ of the pure tensor $\Psi - Q \Psi$ has leading term zero, and so does its transform.  Thus $D(\Psi - Q \Psi) = 0$, by Proposition \ref{proposition:poisson-vanishing}, and so
\begin{equation}\label{eq:boundary-observation}
  D(\Psi) = D(Q \Psi) = - B(b \Psi).
\end{equation}

Both sides of the formula of Theorem \ref{theorem:poisson} are linear in $\Phi$, so we may suppose that $\Phi$ is a pure tensor.  We apply \eqref{eq:boundary-observation} to $\Psi := \mathcal{F}(\Phi - Q \Phi)$.  Its hypothesis holds: by Theorem \ref{theorem:schwartz-space} and \eqref{eq:projector},
\begin{equation*}
  b \mathcal{F} \Psi_w = b \, h(\varepsilon) (\Phi_w - Q_w \Phi_w) = h(\varepsilon) \, b (\Phi_w - Q_w \Phi_w) = 0.
\end{equation*}
Moreover, $b \Psi = b \mathcal{F} \Phi$, since $b \mathcal{F} Q_w \Phi_w = 0$.  The two observations thus give
\begin{equation*}
  D(\Phi - Q \Phi) = - D(\Psi) = B(b \Psi) = B(b \mathcal{F} \Phi).
\end{equation*}
Adding the identity $D(Q \Phi) = - B(b \Phi)$, we obtain $D(\Phi) = B(b \mathcal{F} \Phi) - B(b \Phi)$, which is the formula.

\emph{Uniqueness.}  Let $B'$ be a linear functional on $b(\mathcal{S}_{\mathbb{A}})$ for which the formula of Theorem \ref{theorem:poisson} holds.  Let $\Phi$ be a pure tensor.  We apply that formula to $Q \Phi$.  By \eqref{eq:projector}, we have $b \, Q \Phi = b \Phi$, while $b \mathcal{F} Q \Phi = 0$, since its $w$-factor is zero.  Hence $B'(b \Phi) = - D(Q \Phi) = B(b \Phi)$.  Since the leading terms of the pure tensors span $b(\mathcal{S}_{\mathbb{A}})$, we conclude that $B' = B$.


\section{Statements checked in Lean}\label{sec:lean}
\sloppy
For number fields $F$, the results listed below are formalized in Lean~4 with Mathlib, in the directory \texttt{lean} of \url{\leanrepo}, using only the standard axioms of Lean.  Theorem~\ref{theorem:gauss-sums} and Corollary~\ref{corollary:quadratic-dual} are stated in \texttt{lean/Challenge.lean}, which imports only Mathlib, and proved in \texttt{lean/Solution.lean} from the formal development.  There the self-dual measures and the Hilbert symbols are arbitrary objects with their defining properties, which determine them, and $W^\vee$ is given by \eqref{eq:transform-formula}.  The other results are proved in the development itself.  Not formalized are the case of function fields, Remark~\ref{remark:kubota-transform}, and, at archimedean places, the smoothness in $\mu$ across $\mu = 0$ in Proposition~\ref{proposition:schwartz-origin}.  The second column links to the source of each declaration.

\medskip
\begin{center}
\small
\begin{tabular}{p{0.2\textwidth}p{0.4\textwidth}p{0.32\textwidth}}
Paper & Lean & Hypotheses and scope \\
\hline
Theorem~\ref{theorem:gauss-sums} & \leanentry{NFoldCover.twistedGaussSummation}{Challenge.lean}{344} & The setting of \S\ref{sec:gauss-sums}.  The constants $\rho(\beta)$ in $M_b$ are characterized as the unique family, independent of $T$, $\chi$, $f$ and $W_S$, for which the identity holds; they are not identified with $\lvert b \rvert_S$ times those of Theorem~\ref{theorem:main} for $\psi(b \, \cdot)$.  The absolute convergence of both sums is included. \\
Corollary~\ref{corollary:quadratic-dual} & \leanentry{NFoldCover.quadraticDual}{Challenge.lean}{377} & The setting of \S\ref{sec:quadratic-dual-summation-formula}; the absolute convergence of both sums is included. \\
Lemma~\ref{lemma:lines}\eqref{item:lines-smooth} & \leanentry{MetaplecticPoisson.Datum.radon_smooth_real}{MetaplecticSummation/ArchimedeanSmoothness.lean}{450}, \leanentry{MetaplecticPoisson.Datum.radon_smooth_complex}{MetaplecticSummation/ArchimedeanSmoothness.lean}{469} & Real and complex places: joint smoothness, and rapid decay in $\tau$ of every derivative, uniformly for $v'$ in compact sets. \\
Proposition~\ref{proposition:fourier-transform-genuine} & \leanentry{MetaplecticPoisson.Datum.arch_transform_smooth}{MetaplecticSummation/ArchimedeanSmoothness.lean}{338} & Real and complex places: the smoothness of $\mathcal{F} \Phi$ for $\Phi \in \mathcal{D}$. \\
Definition~\ref{definition:schwartz-space} & \leanentry{MetaplecticPoisson.Datum.schwartz_smooth_real}{MetaplecticSummation/ArchimedeanSmoothness.lean}{423}, \leanentry{MetaplecticPoisson.Datum.schwartz_smooth_complex}{MetaplecticSummation/ArchimedeanSmoothness.lean}{429} & Real and complex places: the elements of $\mathcal{S}$ are smooth. \\
Proposition~\ref{proposition:basic-vector-fixed} & \leanentry{MetaplecticPoisson.Datum.basicCharts_fourier_fixed}{MetaplecticSummation/BasicFixed.lean}{103}, \leanentry{MetaplecticPoisson.Datum.basic_mem}{MetaplecticSummation/BasicMembership.lean}{184} & Finite places prime to $n$. \\
Theorem~\ref{theorem:schwartz-space} & \leanentry{MetaplecticPoisson.Datum.schwartz_inversion}{MetaplecticSummation/Inversion.lean}{500} & Every place; $\widetilde{\mathrm{SL}}_2(F_{\mathfrak{p}})$ is the group generated by $\mu_n$, $n(u)$ and $n^-(u)$. \\
Proposition~\ref{proposition:schwartz-origin} & \leanentry{MetaplecticPoisson.Datum.schwartz_rayBound_real}{MetaplecticSummation/ArchimedeanRays.lean}{683}, \leanentry{MetaplecticPoisson.Datum.schwartz_rayBound_complex}{MetaplecticSummation/ArchimedeanRays.lean}{690}, \leanentry{MetaplecticPoisson.Datum.leading_smooth_real}{MetaplecticSummation/ArchimedeanSmoothness.lean}{436}, \leanentry{MetaplecticPoisson.Datum.leading_smooth_complex}{MetaplecticSummation/ArchimedeanSmoothness.lean}{442} & Real and complex places: the functions \eqref{eq:normalized-ray} are smooth in $v$, and each derivative in $v$ decays rapidly in $\mu$, uniformly for $v$ in compact sets; the leading term $b \Psi$ is smooth.  The smoothness of the extensions in $\mu$ across $\mu = 0$ is not formalized. \\
Theorem~\ref{theorem:poisson} & \leanentry{MetaplecticPoisson.Datum.poisson}{MetaplecticSummation/Poisson.lean}{85} & The reference vector at the places dividing $n$ is $0$, which does not change $\mathcal{S}_{\mathbb{A}}$. \\
Corollary~\ref{corollary:poisson-vanishing} & \leanentry{MetaplecticPoisson.Datum.poisson_of_leading_eq_zero}{MetaplecticSummation/Poisson.lean}{96}, \leanentry{MetaplecticPoisson.Datum.poisson_pure_of_local_vanishing}{MetaplecticSummation/Poisson.lean}{104} & As for Theorem~\ref{theorem:poisson}. \\
Corollary~\ref{corollary:jacquet-decay} & \leanentry{MetaplecticPoisson.Datum.jacquet_decay_arch}{MetaplecticSummation/Jacquet/ArchDecay.lean}{37} & Real and complex places. \\
Theorem~\ref{theorem:main} & \leanentry{MetaplecticPoisson.Datum.lineSummation}{MetaplecticSummation/Jacquet.lean}{133} & The setting of \S\ref{sec:main-formula}. \\
\end{tabular}
\end{center}

\medskip
\noindent Further declarations cited in the text: \leanentry{MetaplecticPoisson.Glues}{MetaplecticSummation/Core.lean}{19}, \leanentry{MetaplecticPoisson.Datum.testSet}{MetaplecticSummation/Schwartz/Charts.lean}{47}, \leanentry{MetaplecticPoisson.radon}{MetaplecticSummation/Core.lean}{26}, \leanentry{MetaplecticPoisson.transform}{MetaplecticSummation/Core.lean}{38}, \leanentry{MetaplecticPoisson.Datum.schwartz}{MetaplecticSummation/Schwartz/Charts.lean}{62}, \leanentry{MetaplecticPoisson.Datum.kubota_projection}{MetaplecticSummation/Cover/Projection.lean}{84}, \leanentry{MetaplecticPoisson.Datum.exists_leading}{MetaplecticSummation/Schwartz.lean}{9}, \leanentry{NFoldCover.residue}{Challenge.lean}{279}, \leanentry{NFoldCover.dualWeight}{Challenge.lean}{251}, \leanentry{MetaplecticPoisson.Datum.root_fourier}{MetaplecticSummation/Jacquet/LineWeights.lean}{24}, \leanentry{NFoldCover.normalizedGauss}{Challenge.lean}{299}, \leanentry{MetaplecticPoisson.Datum.sextic_to_quadratic}{MetaplecticSummation/Gauss.lean}{330}, \leanentry{MetaplecticPoisson.Datum.thetaDifference_vanishing}{MetaplecticSummation/Poisson.lean}{68}.

\bibliography{refs}{} \bibliographystyle{plain}
\end{document}
